% La pollution et le réchauffement climatique — Allemand Tle A — Cours signé OnBuch+
% Programme officiel MINESEC. Compiler avec : tectonic AL13-la-pollution-et-le-rechauffement-climatique.tex
\newcommand{\DOCMATIERE}{Allemand}
\newcommand{\DOCNIVEAU}{Tle A}
\newcommand{\DOCMODULE}{Chapitre 3 : Environnement, santé et bien-être}
\newcommand{\DOCLECON}{AL13}
\newcommand{\DOCTITRE}{La pollution et le réchauffement climatique}
% OnBuch+ — préambule « cours signés » ENRICHI (compilé avec tectonic / XeLaTeX)
% Système visuel « Pop » d'OnBuch+ : fond crème (jamais blanc), bordures encre
% épaisses, ombres « dures » décalées, titres Archivo Black en capitales.
% Version MATHÉMATIQUES : graphes (pgfplots/tikz), boîtes pédagogiques
% (définition, propriété, méthode, attention, le savais-tu, exercice résolu,
% exercice, TP, bilan), page de garde et signature OnBuch+ ancrée en bas.
%
% Macros attendues dans main.tex AVANT \input{preamble} :
%   \DOCMATIERE  \DOCNIVEAU  \DOCMODULE  \DOCLECON (numéro)  \DOCTITRE
\documentclass[11pt]{article}

\usepackage{fontspec}
\usepackage[ngerman,french]{babel} % français = langue principale ; l'allemand via \all{..} / textebox
\usepackage[a4paper,margin=2cm,top=3.4cm,bottom=2.8cm,headheight=1.4cm,footskip=1.3cm]{geometry}
\usepackage{amsmath,amssymb,amsfonts}
\usepackage{mathtools} % \xmapsto, \coloneqq... souvent utilisés par les modèles
\usepackage{cancel} % simplifications barrées (bilans de forces)
% \abs{z} (module, valeur absolue) et \norm{u} (norme), utilisés par les modèles
\providecommand{\abs}{}\let\abs\relax\DeclarePairedDelimiter{\abs}{\lvert}{\rvert}
\providecommand{\norm}{}\let\norm\relax\DeclarePairedDelimiter{\norm}{\lVert}{\rVert}
\usepackage[table]{xcolor} % [table] : fournit aussi \rowcolors (alternance de lignes)
\usepackage{pagecolor}
\usepackage{tikz}
\usetikzlibrary{arrows.meta,positioning,calc,shapes.geometric,shapes.misc,shapes.arrows,decorations.pathmorphing,decorations.pathreplacing,decorations.markings,patterns,shadows,mindmap,trees,fit,backgrounds,matrix,intersections,angles,quotes}
\usepackage{pgfplots}
\usepackage[version=4]{mhchem} % équations nucléaires \ce{^{235}_{92}U}
\usepackage[european,siunitx]{circuitikz} % schémas électriques
\pgfplotsset{compat=1.18}
\usepgfplotslibrary{fillbetween,groupplots} % aires entre courbes (calcul intégral)
\usepackage{enumitem}
\usepackage{fancyhdr}
% [most] charge toutes les bibliothèques tcolorbox (listings, minted, raster,
% documentation...) et gonfle inutilement la mémoire TeX sur un document long
% (~300 boîtes) : on ne charge que ce qui est réellement utilisé.
\usepackage[skins,breakable]{tcolorbox}
\usepackage{graphicx}
\usepackage{colortbl}
\usepackage{booktabs}
\usepackage{tabularx}
\usepackage{diagbox} % case d'en-tête diagonale des tableaux à double entrée
% Colonnes extensibles centrée / gauche / droite (C, L, R), souvent utilisées par les modèles
\newcolumntype{C}{>{\centering\arraybackslash}X}
\newcolumntype{L}{>{\raggedright\arraybackslash}X}
\newcolumntype{R}{>{\raggedleft\arraybackslash}X}
\usepackage{array}
\usepackage{multicol}
\usepackage{multirow}
\usepackage{siunitx}
% parse-numbers=false : accepte aussi \SI{2,5 \times 10^{-3}}{...}
\sisetup{locale=FR,output-decimal-marker={,},group-minimum-digits=4,parse-numbers=false}
\DeclareSIUnit\ohm{\ensuremath{\symup{\Omega}}} % U+2126 absent de la police
\DeclareSIUnit\division{div} % réglages d'oscilloscope (ms/div, V/div)
\DeclareSIUnit\tr{tr} % tours (vitesse de rotation, ex. \SI{1500}{\tr\per\minute})
\DeclareSIUnit\molar{M} % concentration molaire (ex. \SI{3}{\molar}, \SI{1}{\micro\molar})
\DeclareSIUnit\an{an} % année (le modèle confond parfois avec \year, un compteur TeX) — ex. \SI{5}{\kilo\meter\per\an}
\DeclareSIUnit\FCFA{FCFA} % franc CFA (ex. \SI{500}{\FCFA\per\kilo\gram})
\DeclareSIUnit\degreeBrix{\textdegree Brix} % teneur en sucre d'un sirop/jus (réfractométrie)
\DeclareSIUnit\month{mois} % le modèle utilise parfois \month, un compteur TeX, comme unité de durée
\usepackage{qrcode}
% \degree : commande fréquemment utilisée par les modèles pour un angle ou une
% température en dehors de \SI{}{} (non fournie par les paquets chargés).
\providecommand{\degree}{^{\circ}}
% \qed (fin de démonstration), utilisé par les modèles sans amsthm
\providecommand{\qed}{\hfill$\square$}
% \barre : double barre des valeurs interdites dans un tableau de variations (tkz-tab)
\providecommand{\barre}{\ensuremath{\|}}
% environnement proof (amsthm non chargé) utilisé par les modèles
\ifdefined\proof\else\newenvironment{proof}[1][Démonstration]{\par\noindent\textit{#1.}\ }{\qed\par}\fi
\usepackage{lastpage}
\usepackage{hyperref}
\hypersetup{hidelinks}

% ============================================================
% Palette « Pop » OnBuch+ — mêmes hex que la classe P (plus_ui.dart)
% ============================================================
\definecolor{poporange}{HTML}{F59321}
\definecolor{popdark}{HTML}{E07A0C}
\definecolor{popcream}{HTML}{FFF6E8}
\definecolor{popink}{HTML}{1C1714}
\definecolor{poppurple}{HTML}{7A5AE0}
\definecolor{popgreen}{HTML}{2E9E6A}
\definecolor{poppink}{HTML}{E0457B}
\definecolor{popblue}{HTML}{2F7DE1}
\definecolor{popgold}{HTML}{C98A12}
\definecolor{popmuted}{HTML}{8A7F76}
\definecolor{popline}{HTML}{EADFD2}
\definecolor{popgray}{HTML}{8A7F76}
\colorlet{popgrey}{popgray}
% Alias tolérants (le modèle nomme parfois les couleurs différemment)
\colorlet{popwhite}{white}
\colorlet{popblack}{popink}
\colorlet{popred}{poppink}
\colorlet{popyellow}{popgold}
% Teintes claires (fonds de tableaux/encadrés) — le modèle nomme parfois
% les couleurs avec un suffixe L (ex. popblueL) sans les avoir définies.
\colorlet{poporangeL}{poporange!20}
\colorlet{popdarkL}{popdark!20}
\colorlet{poppurpleL}{poppurple!20}
\colorlet{popgreenL}{popgreen!20}
\colorlet{poppinkL}{poppink!20}
\colorlet{popblueL}{popblue!20}
\colorlet{popgoldL}{popgold!20}
\colorlet{popredL}{popred!20}
\colorlet{popgrayL}{popgray!20}
% Teintes claires pour fonds de tableaux / zones
\colorlet{popblueL}{popblue!10!white}
\colorlet{poporangeL}{poporange!14!white}
\colorlet{popgreenL}{popgreen!12!white}
\colorlet{poppurpleL}{poppurple!12!white}
\colorlet{poppinkL}{poppink!12!white}
\colorlet{popgoldL}{popgold!14!white}

\pagecolor{popcream}

% ============================================================
% Typographie
% ============================================================
\defaultfontfeatures{Ligatures=TeX}
\setmainfont{PlusJakartaSans-Medium.ttf}[Path=../../../fonts/,BoldFont=PlusJakartaSans-Bold.ttf]
% Symboles, API et emojis absents de la police principale : repli sur DejaVu Sans (caractères actifs)
\newfontface\symfont{DejaVuSans.ttf}[Path=../../../fonts/]
\makeatletter
\newcommand\symactive[1]{\catcode#1=13 \begingroup\lccode`\~=#1 \lowercase{\endgroup\def~}{{\symfont\char#1}}}
\newcommand\symblank[1]{\catcode#1=13 \begingroup\lccode`\~=#1 \lowercase{\endgroup\def~}{}}
\newcommand\symhyph[1]{\catcode#1=13 \begingroup\lccode`\~=#1 \lowercase{\endgroup\def~}{\mbox{-}}}
\newcount\symc
\newcommand\symrange[2]{\symc=#1 \@whilenum\symc<#2 \do{\expandafter\symactive\expandafter{\the\symc}\advance\symc by 1 }}
\newcommand\symblankrange[2]{\symc=#1 \@whilenum\symc<#2 \do{\expandafter\symblank\expandafter{\the\symc}\advance\symc by 1 }}
\makeatother
\symrange{"0250}{"02FF}\symactive{"2713}\symactive{"2714}\symactive{"2717}\symactive{"2718}\symactive{"2610}\symactive{"2611}\symactive{"2612}
\symhyph{"2011}\symhyph{"2E1A}\symblankrange{"1F300}{"1FAFF}\symblank{"FE0F}
% Maths en sans-serif (Fira Math) avec chiffres et lettres droites de Plus
% Jakarta, pour que les formules se fondent dans le texte.
\usepackage[math-style=ISO,bold-style=ISO]{unicode-math}
\setmathfont{FiraMath-Regular.otf}[Path=../../../fonts/]
\setmathfont{PlusJakartaSans-Medium.ttf}[Path=../../../fonts/,range={up/num,up/{Latin,latin},\mathup}]
\usepackage{icomma} % virgule décimale sans espace en mode math
% Caractères mathématiques Unicode absents des polices de texte (Plus Jakarta,
% Archivo Black) : rendus par leur équivalent mathématique, en texte comme en
% formule — sinon ils s'affichent en carré vide (ex. « Arithmétique dans ℤ »).
\usepackage{newunicodechar}
\newunicodechar{ℝ}{\ensuremath{\mathbb{R}}}
\newunicodechar{ℤ}{\ensuremath{\mathbb{Z}}}
\newunicodechar{ℂ}{\ensuremath{\mathbb{C}}}
\newunicodechar{ℕ}{\ensuremath{\mathbb{N}}}
\newunicodechar{ℚ}{\ensuremath{\mathbb{Q}}}
\newunicodechar{𝔻}{\ensuremath{\mathbb{D}}}
\newunicodechar{α}{\ensuremath{\alpha}}
\newunicodechar{β}{\ensuremath{\beta}}
\newunicodechar{γ}{\ensuremath{\gamma}}
\newunicodechar{δ}{\ensuremath{\delta}}
\newunicodechar{ε}{\ensuremath{\varepsilon}}
\newunicodechar{θ}{\ensuremath{\theta}}
\newunicodechar{λ}{\ensuremath{\lambda}}
\newunicodechar{μ}{\ensuremath{\mu}}
\newunicodechar{σ}{\ensuremath{\sigma}}
\newunicodechar{τ}{\ensuremath{\tau}}
\newunicodechar{φ}{\ensuremath{\varphi}}
\newunicodechar{ψ}{\ensuremath{\psi}}
\newunicodechar{ω}{\ensuremath{\omega}}
\newunicodechar{Σ}{\ensuremath{\Sigma}}
% majuscules produites par \MakeUppercase dans les titres (α-aminés -> Α)
\newunicodechar{Α}{\ensuremath{\alpha}}
\newunicodechar{Β}{\ensuremath{\beta}}
\newunicodechar{Γ}{\ensuremath{\Gamma}}
\newunicodechar{Δ}{\ensuremath{\Delta}}
\newunicodechar{Ε}{\ensuremath{\varepsilon}}
\newunicodechar{Θ}{\ensuremath{\Theta}}
\newunicodechar{Λ}{\ensuremath{\Lambda}}
\newunicodechar{Μ}{\ensuremath{\mu}}
\newunicodechar{Τ}{\ensuremath{\tau}}
\newunicodechar{Φ}{\ensuremath{\Phi}}
\newunicodechar{Ψ}{\ensuremath{\Psi}}
\newunicodechar{Ω}{\ensuremath{\Omega}}
\newunicodechar{↦}{\ensuremath{\mapsto}}
\newunicodechar{∈}{\ensuremath{\in}}
\newunicodechar{∉}{\ensuremath{\notin}}
\newunicodechar{∧}{\ensuremath{\wedge}}
\newunicodechar{⇔}{\ensuremath{\Leftrightarrow}}
\newunicodechar{⟺}{\ensuremath{\Longleftrightarrow}}
\newunicodechar{⇒}{\ensuremath{\Rightarrow}}
\newunicodechar{⟹}{\ensuremath{\Longrightarrow}}
\newunicodechar{∘}{\ensuremath{\circ}}
\newunicodechar{∪}{\ensuremath{\cup}}
\newunicodechar{∩}{\ensuremath{\cap}}
\newunicodechar{≡}{\ensuremath{\equiv}}
\newunicodechar{⊂}{\ensuremath{\subset}}
\newunicodechar{⊥}{\ensuremath{\perp}}
\newunicodechar{∃}{\ensuremath{\exists}}
\newunicodechar{∀}{\ensuremath{\forall}}
\newunicodechar{∛}{\ensuremath{\sqrt[3]{\,}}}
\newunicodechar{∜}{\ensuremath{\sqrt[4]{\,}}}
\newunicodechar{⃗}{\ensuremath{{}^{\to}}}
\newunicodechar{ⁿ}{\ifmmode^{n}\else\textsuperscript{\ensuremath{n}}\fi}
\newunicodechar{ˣ}{\ifmmode^{x}\else\textsuperscript{\ensuremath{x}}\fi}
\newunicodechar{ᵇ}{\ifmmode^{b}\else\textsuperscript{\ensuremath{b}}\fi}
\newunicodechar{ʳ}{\ifmmode^{r}\else\textsuperscript{\ensuremath{r}}\fi}
\newunicodechar{ᵘ}{\ifmmode^{u}\else\textsuperscript{\ensuremath{u}}\fi}
\newunicodechar{ʸ}{\ifmmode^{y}\else\textsuperscript{\ensuremath{y}}\fi}
\newunicodechar{ᵃ}{\ifmmode^{a}\else\textsuperscript{\ensuremath{a}}\fi}
\newunicodechar{ᵉ}{\ifmmode^{e}\else\textsuperscript{\ensuremath{e}}\fi}
\newunicodechar{ᵏ}{\ifmmode^{k}\else\textsuperscript{\ensuremath{k}}\fi}
\newunicodechar{ᵖ}{\ifmmode^{p}\else\textsuperscript{\ensuremath{p}}\fi}
\newunicodechar{ᴺ}{\ifmmode^{N}\else\textsuperscript{\ensuremath{N}}\fi}
\newunicodechar{ᶜ}{\ifmmode^{c}\else\textsuperscript{\ensuremath{c}}\fi}
\newunicodechar{ʰ}{\ifmmode^{h}\else\textsuperscript{\ensuremath{h}}\fi}
\newunicodechar{ᵠ}{\ifmmode^{q}\else\textsuperscript{\ensuremath{q}}\fi}
\newunicodechar{ₙ}{\ifmmode_{n}\else\textsubscript{\ensuremath{n}}\fi}
\newunicodechar{ₐ}{\ifmmode_{a}\else\textsubscript{\ensuremath{a}}\fi}
\newunicodechar{ᵢ}{\ifmmode_{i}\else\textsubscript{\ensuremath{i}}\fi}
\newunicodechar{ᵣ}{\ifmmode_{r}\else\textsubscript{\ensuremath{r}}\fi}
\newunicodechar{ₖ}{\ifmmode_{k}\else\textsubscript{\ensuremath{k}}\fi}
\newunicodechar{ᵦ}{\ifmmode_{\beta}\else\textsubscript{\ensuremath{\beta}}\fi}
\newunicodechar{ₓ}{\ifmmode_{x}\else\textsubscript{\ensuremath{x}}\fi}
\newfontfamily\popdisplay{ArchivoBlack-Regular.ttf}[Path=../../../fonts/]
\newfontfamily\poplogo{SpaceGrotesk-Bold.ttf}[Path=../../../fonts/]
\newfontfamily\popsemi{PlusJakartaSans-SemiBold.ttf}[Path=../../../fonts/]
\newfontfamily\popxbold{PlusJakartaSans-ExtraBold.ttf}[Path=../../../fonts/]
\newfontfamily\popxboldls{PlusJakartaSans-ExtraBold.ttf}[Path=../../../fonts/,LetterSpace=3.0]

\setlist[enumerate]{leftmargin=1.7em,itemsep=5pt,topsep=4pt}
\setlist[itemize]{leftmargin=1.7em,itemsep=4pt,topsep=4pt,label={\color{poporange}\textbullet}}
\setlength{\parindent}{0pt}
\setlength{\parskip}{5pt}
\renewcommand{\arraystretch}{1.25}

% pgfplots : style Pop par défaut
\pgfplotsset{
  every axis/.append style={axis line style={popink,line width=1pt},tick style={popink},
    grid style={popline,line width=0.5pt},label style={font=\small},tick label style={font=\footnotesize}},
  popcurve/.style={poporange,line width=1.8pt,no markers,smooth},
}
% Même style disponible dans un \draw TikZ simple (hors pgfplots)
\tikzset{popcurve/.style={poporange,line width=1.8pt}}
\tikzset{popgrid/.style={popline,very thin}}
\colorlet{popcurve}{poporange} % le nom du style est parfois utilisé comme couleur

% ============================================================
% Pill / squiggle
% ============================================================
\newcommand{\poppill}[3]{%
  \tikz[baseline=-0.55ex]\node[draw=popink,line width=1.2pt,rounded corners=2.6pt,%
    fill=#2,inner xsep=6.5pt,inner ysep=3pt]{%
    \popxboldls\fontsize{7}{7}\selectfont\color{#3}\MakeUppercase{#1}};%
}
\newcommand{\popsquiggle}[1]{%
  \tikz[baseline=-0.4ex]{
    \draw[#1,line width=2.6pt,line cap=round]
      plot [smooth] coordinates {(0,0.09) (0.34,0.01) (0.68,0.21) (1.02,0.01) (1.36,0.10)};
  }%
}
% Mot-clé surligné (marqueur orange sous le texte)
\newcommand{\cle}[1]{{\popxbold\color{popdark}#1}}

% ============================================================
% En-tête / pied
% ============================================================
\pagestyle{fancy}
\fancyhf{}
\renewcommand{\headrulewidth}{0pt}
\renewcommand{\footrulewidth}{0pt}
\lhead{\raisebox{-0.15ex}{\poplogo\fontsize{13}{13}\selectfont\color{popink}OnBuch\color{poporange}+}}
\rhead{\poppill{\DOCMATIERE\ \textbullet\ \DOCNIVEAU\ \textbullet\ Leçon \DOCLECON}{white}{popink}}
\lfoot{\popxbold\fontsize{7.6}{7.6}\selectfont\color{popmuted}\MakeUppercase{Cours signé OnBuch+}\ \textbullet\ {\popsemi\DOCTITRE}}
\rfoot{\tikz[baseline=-0.6ex]\node[rounded corners=8.5pt,fill=popink,minimum height=17pt,minimum width=17pt,inner xsep=5pt,inner sep=0]{\popxbold\fontsize{8}{8}\selectfont\color{popcream}\thepage\,/\,\pageref*{LastPage}};}

% ============================================================
% Titres
% ============================================================
\newcommand{\chaptitre}[2]{%
  \begin{tcolorbox}[enhanced,colback=poporange,colframe=popink,boxrule=2.6pt,arc=11pt,
      drop shadow=popink,left=15pt,right=15pt,top=12pt,bottom=11pt,before skip=3pt,after skip=18pt]
    \poppill{#2}{popink}{popcream}\par\vspace{9pt}
    {\raggedright\hyphenpenalty=10000\exhyphenpenalty=10000\popdisplay\fontsize{22}{24}\selectfont\color{popcream}\MakeUppercase{#1}\par}
  \end{tcolorbox}
}
% Section numérotée : pastille orange avec le numéro + titre Archivo
\newcounter{popsec}
\newcommand{\coursec}[1]{%
  \stepcounter{popsec}\setcounter{popsub}{0}%
  \par\vspace{16pt}\noindent
  \begin{minipage}{\linewidth}
  \tikz[baseline=-0.7ex]\node[draw=popink,line width=1.6pt,fill=poporange,rounded corners=4pt,minimum size=22pt,inner sep=0]{\popdisplay\fontsize{12}{12}\selectfont\color{popcream}\thepopsec};%
  \hspace{8pt}{\popdisplay\fontsize{15}{17}\selectfont\color{popink}\MakeUppercase{#1}}\par
  \vspace{4pt}\noindent{\color{poporange}\rule{36pt}{3.6pt}}
  \end{minipage}\par\nopagebreak\vspace{8pt}
}
\newcounter{popsub}
\newcommand{\courssub}[1]{%
  \stepcounter{popsub}%
  \par\vspace{9pt}\noindent{\popxbold\fontsize{11.5}{13}\selectfont\color{popink}{\color{poporange}\thepopsec.\thepopsub}\hspace{6pt}#1}\par\nopagebreak\vspace{4pt}
}

% ============================================================
% Boîtes pédagogiques — même recette : fond blanc, bordure épaisse colorée,
% ombre dure encre, badge plein. Titre optionnel [#1].
% ============================================================
\tcbset{popbase/.style={breakable,enhanced,colback=white,boxrule=2.3pt,arc=9pt,
  shadow={3pt}{-3pt}{0pt}{popink},left=14pt,right=14pt,top=11pt,bottom=11pt,before skip=11pt,after skip=15pt,
  fontupper=\popsemi\fontsize{10.6}{14.5}\selectfont\color{popink},
  fontlower=\popsemi\fontsize{10.6}{14.5}\selectfont\color{popink}}}
% Coche dessinée (le glyphe ✓ n'existe pas dans les polices du document)
\newcommand{\popcheck}[1]{\tikz[baseline=-0.5ex]\draw[#1,line width=1.6pt,line cap=round,line join=round](-0.13,0.0)--(-0.04,-0.09)--(0.13,0.1);}
\providecommand{\coche}{\popcheck{popgreen}}
\providecommand{\resultat}[1]{\par\smallskip\noindent\fcolorbox{popgreen}{popgreenL}{\parbox{\dimexpr\linewidth-2\fboxsep-2\fboxrule\relax}{\textbf{Résultat.} #1}}\par\smallskip}
\newcommand{\popboxhead}[3]{\poppill{#1}{#2}{white}\ifx\relax#3\relax\else\hspace{6pt}{\popxbold\fontsize{10.6}{12}\selectfont\color{popink}#3}\fi\par\vspace{7pt}}

\newtcolorbox{aretenir}[1][]{popbase,colframe=popgreen,before upper={\popboxhead{À retenir}{popgreen}{#1}}}
% Texte en allemand (document de lecture, dialogue, extrait) : cadre
% d'étude. Un titre optionnel (source, auteur) s'affiche dans l'en-tête.
\newtcolorbox{textebox}[1][]{popbase,colframe=popblue,before upper={\popboxhead{Text}{popblue}{#1}\begin{otherlanguage}{ngerman}},after upper={\end{otherlanguage}}}
% Marques ✓ / ✗ (forme correcte / fautive) souvent utilisées par les modèles
\providecommand{\cross}{\textcolor{poppink}{\ensuremath{\times}}}
\providecommand{\xmark}{\cross}\providecommand{\crossmark}{\cross}\providecommand{\cmark}{\ensuremath{\checkmark}}
% Ligne de réponse / justification à compléter (inventée par les modèles)
\providecommand{\justification}[1]{\par\noindent\textit{Justification~:} \dotfill\par}
% \textipa (package tipa non chargé) : la police ne couvre pas l'API -> simple passage
\providecommand{\textipa}[1]{#1}
% Soulignement (ulem non chargé) : \uline, \ul
\providecommand{\uline}[1]{\underline{#1}}\providecommand{\ul}[1]{\underline{#1}}
% \glossaire{\item ...} inventée par les modèles : liste de glossaire
\providecommand{\glossaire}[1]{\begin{itemize}#1\end{itemize}}
% \alert (beamer) inventée par les modèles : texte mis en évidence
\providecommand{\alert}[1]{\textbf{\textcolor{poppink}{#1}}}
% variantes ASCII d'ordinaux écrites en macro
\providecommand{\iere}{\textsuperscript{re}}\providecommand{\ieme}{\textsuperscript{e}}\providecommand{\ere}{\textsuperscript{re}}\providecommand{\eme}{\textsuperscript{e}}\providecommand{\er}{\textsuperscript{er}}\providecommand{\ie}{\textsuperscript{e}}
% Ordinaux français écrits en macro par les modèles : 1\ière, 3\ième
\newcommand{\ière}{\textsuperscript{re}}\newcommand{\ième}{\textsuperscript{e}}
% \exemple (inventée par les modèles) : étiquette « Exemple : »
\providecommand{\exemple}{\textbf{Exemple~:}\ }
% \square utilisable aussi hors mode math (cases à cocher)
\AtBeginDocument{\let\squareMath\square\renewcommand{\square}{\ensuremath{\squareMath}}}
% Mot ou phrase en allemand à l'intérieur d'un paragraphe français
\newcommand{\all}[1]{\ifmmode\text{\textit{#1}}\else\begingroup\selectlanguage{ngerman}\itshape #1\endgroup\fi}
\let\esp\all % alias toléré
\newtcolorbox{exemplebox}[1][]{popbase,colframe=poporange,before upper={\popboxhead{Exemple}{poporange}{#1}}}
\newtcolorbox{definition}[1][]{popbase,colframe=popblue,before upper={\popboxhead{Définition}{popblue}{#1}}}
\newtcolorbox{propriete}[1][]{popbase,colframe=poppurple,before upper={\popboxhead{Propriété}{poppurple}{#1}}}
\newtcolorbox{methode}[1][]{popbase,colframe=popgold,before upper={\popboxhead{Méthode}{popgold}{#1}}}
\newtcolorbox{attention}[1][]{popbase,colframe=poppink,before upper={\popboxhead{Attention}{poppink}{#1}}}
\newtcolorbox{experience}[1][]{popbase,colframe=popblue,colback=popblueL,before upper={\popboxhead{Expérience}{popblue}{#1}}}
% « Le savais-tu ? » : carte violette pleine (contexte, culture, Cameroun)
\newtcolorbox{savaistu}[1][]{popbase,colback=poppurple,colframe=popink,
  fontupper=\popsemi\fontsize{10.4}{14.2}\selectfont\color{white},
  before upper={\poppill{Le savais-tu\,?}{white}{poppurple}\ifx\relax#1\relax\else\hspace{6pt}{\popxbold\color{white}#1}\fi\par\vspace{7pt}}}
% Exercice résolu : énoncé en haut, \tcblower puis la solution sur fond crème
\newcounter{popexo}\newcounter{popexr}
\newtcolorbox{exoresolu}[1][]{popbase,colframe=poporange,colbacklower=popcream,
  segmentation style={popink,line width=1.2pt,dashed},
  before upper={\stepcounter{popexr}\popboxhead{Exercice résolu \thepopexr}{poporange}{#1}},
  before lower={\poppill{Solution}{popink}{popcream}\par\vspace{6pt}}}
% Exercice (non résolu en place) — la correction est dans la partie Corrigés
\newtcolorbox{exercice}[1][]{popbase,colframe=popink,
  before upper={\stepcounter{popexo}\popboxhead{Exercice \thepopexo}{popink}{#1}}}
% Corrigé
\newtcolorbox{corrige}[1][]{popbase,colframe=popgreen,colback=popgreenL,
  before upper={\popboxhead{Corrigé}{popgreen}{#1}}}
% Objectifs / prérequis / bilan
\newtcolorbox{objectifs}{popbase,colframe=popink,colback=poporangeL,
  before upper={\popboxhead{Objectifs de la leçon}{popink}{}}}
\newtcolorbox{prerequis}{popbase,colframe=popink,colback=popblueL,
  before upper={\popboxhead{Prérequis}{popblue}{}}}
\newtcolorbox{bilan}{popbase,colframe=popink,colback=white,boxrule=2.8pt,
  before upper={\popboxhead{Fiche bilan}{poporange}{L'essentiel à connaître pour l'examen}}}
% Figure encadrée avec légende
\newtcolorbox{popfigure}[1][]{enhanced,breakable=false,colback=white,colframe=popline,boxrule=1.4pt,arc=8pt,
  left=10pt,right=10pt,top=10pt,bottom=8pt,before skip=10pt,after skip=12pt,halign=center,
  lower separated=false,halign lower=center,
  fontlower=\popsemi\fontsize{9}{11.5}\selectfont\color{popmuted},
  #1}
\newcommand{\legende}[1]{\par\vspace{6pt}{\popsemi\fontsize{9}{11.5}\selectfont\color{popmuted}#1}}

% ============================================================
% Signature OnBuch+
% ============================================================
\newcommand{\poplogoinline}[1]{{\poplogo\fontsize{#1}{#1}\selectfont\color{popink}OnBuch\color{poporange}+}}
\newcommand{\signatureonbuch}{%
  \begin{tcolorbox}[enhanced,colback=popink,colframe=popink,boxrule=2.2pt,arc=10pt,
      drop shadow=poporange,left=14pt,right=14pt,top=10pt,bottom=10pt,before skip=0pt,after skip=0pt]
    \tikz[baseline=-0.7ex]\node[circle,fill=popgreen,draw=popcream,line width=1.3pt,minimum size=22pt,inner sep=0]{\popcheck{white}};%
    \hspace{9pt}\begin{minipage}[c]{0.62\linewidth}
      {\popxbold\fontsize{12}{13}\selectfont\color{popcream}Signé\ \textbullet\ {\poplogo OnBuch\color{poporange}+}}\par\vspace{2pt}
      {\raggedright\popsemi\fontsize{8}{10.5}\selectfont\color{popline}\MakeUppercase{Équipe pédagogique OnBuch+}\\\MakeUppercase{Conforme au programme officiel MINESEC}\par}
    \end{minipage}\hfill
    {\poplogo\fontsize{22}{22}\selectfont\color{popcream}OnBuch\color{poporange}+}
  \end{tcolorbox}%
}

% ============================================================
% Page de garde
% #1 titre · #2 matière · #3 niveau · #4 module · #5 n° leçon · #6 durée ·
% #7 description / objectifs (texte ou itemize)
% ============================================================
\newcommand{\pagegarde}[7]{%
  \thispagestyle{empty}
  \newgeometry{a4paper,margin=1.7cm,top=1.6cm,bottom=1.4cm}%
  \begin{tikzpicture}[remember picture,overlay]
    \fill[poporange!18] ([xshift=-3.2cm,yshift=-3.6cm]current page.north east) circle (4.6cm);
    \fill[poppurple!14] ([xshift=2.4cm,yshift=7.6cm]current page.south west) circle (3.8cm);
  \end{tikzpicture}%
  {\poplogo\fontsize{30}{30}\selectfont\color{popink}OnBuch\color{poporange}+}\hfill
  \raisebox{0.6ex}{\poppill{Cours signé \textbullet\ Premium}{popink}{popcream}}\par
  \vspace{1pt}\popsquiggle{poppurple}\par
  \vspace{8pt}
  \begin{tcolorbox}[enhanced,colback=poporange,colframe=popink,boxrule=2.8pt,arc=13pt,
      drop shadow=popink,left=18pt,right=18pt,top=12pt,bottom=13pt,after skip=10pt]
    \poppill{#2 \textbullet\ #3}{popink}{popcream}\hspace{5pt}\poppill{#4}{white}{popink}\par\vspace{8pt}
    {\popdisplay\fontsize{14}{14}\selectfont\color{popink}LEÇON #5}\par\vspace{4pt}
    {\raggedright\hyphenpenalty=10000\exhyphenpenalty=10000\popdisplay\fontsize{25}{27}\selectfont\color{popcream}\MakeUppercase{#1}\par}
    \vspace{8pt}
    {\popxbold\fontsize{9.5}{9.5}\selectfont\color{popink}\MakeUppercase{Durée conseillée : #6}}
  \end{tcolorbox}
  \begin{tcolorbox}[enhanced,colback=white,colframe=popink,boxrule=2.2pt,arc=10pt,
      drop shadow=popink,left=16pt,right=16pt,top=10pt,bottom=10pt,after skip=8pt]
    \poppill{Dans cette leçon}{poppurple}{white}\par\vspace{5pt}
    \popsemi\fontsize{9.3}{12.6}\selectfont\color{popink}#7
  \end{tcolorbox}
  \noindent\begin{minipage}[t]{0.487\linewidth}
    \begin{tcolorbox}[enhanced,colback=popgreen,colframe=popink,boxrule=2.2pt,arc=10pt,
        drop shadow=popink,left=11pt,right=11pt,top=10pt,bottom=10pt,height=3.1cm,valign=center]
      \tcbox[colback=white,colframe=popink,boxrule=1.3pt,arc=5pt,boxsep=0pt,nobeforeafter,
        left=3pt,right=3pt,top=3pt,bottom=3pt,valign=center]{\qrcode[height=1.9cm]{https://play.google.com/store/apps/details?id=com.luvvix.onbuch&hl=fr}}%
      \hspace{8pt}\begin{minipage}[c]{0.5\linewidth}
        {\popdisplay\fontsize{11}{12.5}\selectfont\color{white}\MakeUppercase{Télécharge OnBuch}}\par\vspace{4pt}
        {\raggedright\popsemi\fontsize{8.4}{11}\selectfont\color{white}Gratuit \textbullet\ Google Play\\Tuteur IA, annales, concours, résultats\par}
      \end{minipage}
    \end{tcolorbox}
  \end{minipage}\hfill
  \begin{minipage}[t]{0.487\linewidth}
    \begin{tcolorbox}[enhanced,colback=poppurple,colframe=popink,boxrule=2.2pt,arc=10pt,
        drop shadow=popink,left=11pt,right=11pt,top=10pt,bottom=10pt,height=3.1cm,valign=center]
      \tikz[baseline=-0.7ex]\node[circle,fill=white,draw=popink,line width=1.3pt,minimum size=1.6cm,inner sep=0]{\popdisplay\fontsize{24}{24}\selectfont\color{poppurple}+};%
      \hspace{8pt}\begin{minipage}[c]{0.6\linewidth}
        {\popdisplay\fontsize{11}{12.5}\selectfont\color{white}\MakeUppercase{Passe à OnBuch+}}\par\vspace{4pt}
        {\raggedright\popsemi\fontsize{8.4}{11}\selectfont\color{white}Tous les cours signés, Tuteur IA illimité, fiches PDF — onglet OnBuch+ de l'app\par}
      \end{minipage}
    \end{tcolorbox}
  \end{minipage}
  \vspace{8pt}
  \popcontenu
  \vfill
  \signatureonbuch
  \restoregeometry
  \newpage
}

% Rangée « Ce cours contient » (page de garde)
\newcommand{\poptile}[3]{% #1 couleur #2 gros texte #3 légende
  \begin{tcolorbox}[enhanced,colback=#1,colframe=popink,boxrule=2pt,arc=9pt,shadow={2.5pt}{-2.5pt}{0pt}{popink},
      left=6pt,right=6pt,top=9pt,bottom=9pt,halign=center,valign=center,height=1.75cm,width=\linewidth,nobeforeafter]
    {\popdisplay\fontsize{12.5}{13}\selectfont\color{white}\MakeUppercase{#2}}\par\vspace{3pt}
    {\centering\popsemi\fontsize{7.8}{9.5}\selectfont\color{white}#3\par}
  \end{tcolorbox}}
\newcommand{\popcontenu}{%
  \par\noindent{\popxboldls\fontsize{7.5}{7.5}\selectfont\color{popmuted}CE COURS SIGNÉ CONTIENT}\par\vspace{7pt}
  \noindent\begin{minipage}[t]{0.235\linewidth}\poptile{popblue}{Cours}{complet, illustré, pas à pas}\end{minipage}\hfill
  \begin{minipage}[t]{0.235\linewidth}\poptile{poporange}{Méthodes}{et exercices résolus}\end{minipage}\hfill
  \begin{minipage}[t]{0.235\linewidth}\poptile{poppink}{Exercices}{type examen + corrigés}\end{minipage}\hfill
  \begin{minipage}[t]{0.235\linewidth}\poptile{popgreen}{Fiche bilan}{l'essentiel à retenir}\end{minipage}\par}

% Sommaire
\newtcolorbox{sommaire}{enhanced,colback=white,colframe=popink,boxrule=2.2pt,arc=10pt,
  drop shadow=popink,left=18pt,right=18pt,top=13pt,bottom=13pt,before skip=6pt,after skip=20pt,
  before upper={\poppill{Sommaire}{poppurple}{white}\par\vspace{9pt}\popsemi\fontsize{10.6}{14.5}\selectfont\color{popink}}}

% Signature de fin de cours — poussée tout en bas de la dernière page
\newcommand{\findecours}{%
  \par\vfill\nopagebreak
  \begin{center}\popsquiggle{poporange}\end{center}\vspace{4pt}
  \signatureonbuch
}
\AtBeginDocument{\def\llbracket{\mathopen{[\mkern-3.5mu[}}\def\rrbracket{\mathclose{]\mkern-3.5mu]}}} % intervalles d'entiers (glyphe ⟦ absent de la police)
% Glyphes absents de Fira Math : redéfinis à partir de glyphes disponibles
\AtBeginDocument{%
  \def\setminus{\mathbin{\backslash}}\let\smallsetminus\setminus
  \def\vdots{\mathord{\vbox{\baselineskip3.2pt\lineskiplimit0pt\kern4pt\hbox{.}\hbox{.}\hbox{.}}}}%
  \def\ddots{\mathinner{\mkern1mu\raise6.5pt\hbox{.}\mkern2mu\raise3.6pt\hbox{.}\mkern2mu\raise0.7pt\hbox{.}\mkern1mu}}%
  \def\triangle{\mathord{\scalebox{0.9}{$\symup{\Delta}$}}}%
  \def\nabla{\mathord{\raisebox{\depth}{\scalebox{1}[-1]{$\symup{\Delta}$}}}}%
  \def\checkmark{\popcheck{popdark}}%
  \def\star{\mathbin{\ast}}\def\bigstar{\mathord{\ast}}%
  \def\diamond{\mathbin{\scalebox{0.6}{$\Diamond$}}}%
  \def\bigcup{\mathop{\mathchoice{\raisebox{-0.3em}{\scalebox{1.7}{$\cup$}}}{\scalebox{1.25}{$\cup$}}{\cup}{\cup}}\displaylimits}%
  \def\bigcap{\mathop{\mathchoice{\raisebox{-0.3em}{\scalebox{1.7}{$\cap$}}}{\scalebox{1.25}{$\cap$}}{\cap}{\cap}}\displaylimits}%
  \def\lesseqqgtr{\mathrel{\lessgtr}}\def\bigoplus{\mathop{\oplus}\displaylimits}%
}
\providecommand{\ssi}{si et seulement si} % abréviation fréquente
\AtBeginDocument{\providecommand{\wideparen}{\overparen}} % arc de cercle

%% FIGFIX-BEGIN v3 — correctifs globaux des figures (docs/onbuch-generation/tools/figfix)
\usepackage{adjustbox}\usepackage{etoolbox}
% toute figure TikZ/pgfplots plus large que la ligne est réduite à la largeur disponible
\BeforeBeginEnvironment{tikzpicture}{\begin{adjustbox}{max width=\linewidth}}
\AfterEndEnvironment{tikzpicture}{\end{adjustbox}}
% légendes pgfplots lisibles par-dessus les courbes
\pgfplotsset{every axis/.append style={legend style={fill=white,fill opacity=0.92,text opacity=1,draw=black!15,inner sep=3pt}}}
% étiquettes d'arêtes (syntaxe edge["..."]) avec fond blanc
\tikzset{every edge quotes/.append style={fill=white,inner sep=1.2pt}}
% bulles de cartes mentales : texte à la ligne dans la bulle, bulle un peu plus grande
\tikzset{every concept/.append style={text width=2.0cm,align=center,minimum size=2.3cm,inner sep=2pt}}
%% FIGFIX-END
\begin{document}

\pagegarde{La pollution et le réchauffement climatique}{Allemand}{Tle A}{Chapitre 3 : Environnement, santé et bien-être}{AL13}{2 h}{Cette leçon de type « texte » propose la lecture et le commentaire d'un texte allemand original sur la pollution et le réchauffement climatique. Elle consolide le vocabulaire thématique de l'environnement et introduit deux points de grammaire essentiels : le subjonctif II (l'irréel) et les subordonnées de conséquence et de temps.
\vspace{4pt}
\begin{multicols}{2}\small
\begin{itemize}
\item À la fin de cette leçon, je sais comprendre globalement puis en détail un texte allemand sur la pollution et le réchauffement climatique
\item Je sais employer correctement le vocabulaire thématique de l'environnement (die Umweltverschmutzung, die Erderwärmung, das Klima, das CO₂, der Wald, die Dürre) avec articles et pluriels
\item Je sais former et employer le subjonctif II pour exprimer l'irréel, l'hypothèse et le souhait
\item Je sais construire des subordonnées de conséquence (so dass, sodass, so... dass) et de temps (als, wenn, nachdem, bevor, während, seitdem, bis)
\item Je sais répondre par écrit et à l'oral à des questions de compréhension sur le texte
\item Je sais résumer un texte en quelques phrases avec mes propres mots
\end{itemize}\end{multicols}}

\begin{sommaire}
\begin{multicols}{2}
\begin{enumerate}
\item Texte support : « Die Erde wird wärmer »
\item Compréhension du texte
\item Lexique thématique : Umwelt und Klima
\item Grammaire 1 : le subjonctif II (l'irréel)
\item Grammaire 2 : subordonnées de conséquence et de temps
\item Méthode du commentaire et production guidée
\item Activité d'intégration
\item Méthodes et exercices résolus
\item Exercices
\item Corrigés des exercices
\item Fiche bilan
\end{enumerate}
\end{multicols}
\end{sommaire}

\begin{objectifs}
\begin{itemize}
\item À la fin de cette leçon, je sais comprendre globalement puis en détail un texte allemand sur la pollution et le réchauffement climatique
\item Je sais employer correctement le vocabulaire thématique de l'environnement (die Umweltverschmutzung, die Erderwärmung, das Klima, das CO₂, der Wald, die Dürre) avec articles et pluriels
\item Je sais former et employer le subjonctif II pour exprimer l'irréel, l'hypothèse et le souhait
\item Je sais construire des subordonnées de conséquence (so dass, sodass, so... dass) et de temps (als, wenn, nachdem, bevor, während, seitdem, bis)
\item Je sais répondre par écrit et à l'oral à des questions de compréhension sur le texte
\item Je sais résumer un texte en quelques phrases avec mes propres mots
\item Je sais rédiger un court commentaire structuré (introduction, développement, conclusion) sur un sujet d'actualité environnementale
\item Je sais produire un court paragraphe argumenté sur la protection de l'environnement au Cameroun
\end{itemize}
\end{objectifs}

\begin{prerequis}
\begin{itemize}
\item Les temps principaux de l'indicatif : Präsens, Präteritum, Perfekt, Futur I (Première)
\item La place du verbe dans la proposition principale et la subordonnée (verbe conjugué en fin de subordonnée)
\item Le Präteritum des verbes forts et faibles (base de formation du subjonctif II)
\item Le vocabulaire général de la nature et de la météo vu en Seconde/Première (das Wetter, der Regen, die Hitze)
\item Les connecteurs de base : weil, denn, aber, und, deshalb
\end{itemize}
\end{prerequis}

\newpage

\coursec{Texte support : „Die Erde wird wärmer“}

\courssub{Situation d'accroche et problématique}

En classe de Terminale à Yaoundé, les élèves discutent des pluies de plus en plus irrégulières et des saisons sèches prolongées qui affectent les récoltes au Cameroun : le maïs souffre, le cacao craint la chaleur, les points d'eau tarissent plus tôt. En Allemagne, en Autriche et en Suisse, on parle aussi chaque jour de \all{Klimawandel} : canicules, inondations, fonte des glaciers alpins. Comment les jeunes germanophones débattent-ils de la pollution et du réchauffement climatique ? Quels mots utilisent-ils pour décrire le CO\textsubscript{2}, la déforestation ou la sécheresse ? Et comment expriment-ils ce qui se passerait SI l'on agissait autrement ? Cette leçon nous donne les outils linguistiques pour lire, comprendre et commenter un texte d'actualité sur ce sujet majeur du Bac.

\textbf{Problématique :} Comment un texte allemand d'actualité présente-t-il les causes, les conséquences et les solutions du réchauffement climatique, et quels outils grammaticaux (subjonctif II, subordonnées de conséquence et de temps) utilise-t-il ?

\courssub{Présentation du texte}

Le texte support est un article original rédigé dans le style d'un magazine jeunesse allemand (type \all{Spiegel} ou \all{Fluter} pour la jeunesse). Il compte environ 200 mots, soit 15 à 20 lignes, et il est construit selon un plan très clair : \textbf{problème} $\rightarrow$ \textbf{causes} $\rightarrow$ \textbf{conséquences} $\rightarrow$ \textbf{solutions}. Ce plan est typique des articles de presse allemands sur l'environnement : le lecteur est d'abord informé du phénomène, puis des responsabilités, puis des effets visibles, et enfin des pistes d'action.

Le texte contient volontairement des exemples des deux points de grammaire de la leçon : le \cle{subjonctif II} pour exprimer l'irréel (\all{Wenn wir weniger Energie verschwendeten, hätten wir eine sauberere Umwelt.}) et les \cle{subordonnées de conséquence et de temps} (\all{Die Erde wird wärmer, so dass die Gletscher schmelzen.} ; \all{Als der Regen ausblieb, vertrockneten die Felder.}). Nous les observerons ici, puis nous les étudierons en détail dans les sections 4 et 5.

\begin{methode}[Comment lire un texte au Bac]
\begin{enumerate}
\item Lire d'abord le \textbf{titre} et la \textbf{source} : ils annoncent le thème et le type de texte.
\item Faire une \textbf{première lecture globale}, sans dictionnaire, pour saisir le sens général.
\item \textbf{Souligner les mots transparents} (mots proches du français) : \all{die Emission, die Temperatur, global, die Energie, das Problem}.
\item Faire une \textbf{deuxième lecture détaillée} avec le glossaire : noter les mots-clés de chaque paragraphe.
\item \textbf{Repérer la structure du texte} : où commence le problème ? où sont les causes ? les conséquences ? les solutions ?
\end{enumerate}
\end{methode}

\begin{attention}[Piège de la première lecture]
Ne traduis \textbf{jamais mot à mot} pendant la première lecture ! Un texte allemand ne se décode pas comme un rébus : l'ordre des mots diffère du français (verbe en deuxième position, verbe à la fin dans les subordonnées). Cherche d'abord le \cle{sens global} : de quoi parle le texte ? qui fait quoi ? Les détails viendront à la deuxième lecture. Exemple : \all{Als der Regen letztes Jahr ausblieb, vertrockneten die Felder der Bauern.} — si tu traduis mot à mot, tu es perdu ; si tu repères \all{Regen} (pluie), \all{Felder} (champs), \all{Bauern} (paysans), tu comprends déjà : « pas de pluie $\rightarrow$ champs en danger ».
\end{attention}

\courssub{Le texte : lecture et observation}

\textbf{Consigne de lecture :} lis le texte une première fois en silence, sans t'arrêter sur les mots inconnus. Souligne les mots que tu reconnais. Puis écoute la lecture à voix haute par l'enseignant et vérifie tes hypothèses.

\begin{textebox}[Die Erde wird wärmer — article de magazine jeunesse]
Seit hundert Jahren steigt die Temperatur auf der Erde. Wissenschaftler aus aller Welt sind sich einig: der Klimawandel ist real, und der Mensch trägt die Hauptverantwortung dafür.

Die Ursachen sind bekannt. Die Fabriken und Autos produzieren viel CO\textsubscript{2}, und die Abgase vergiften die Luft in den großen Städten. Außerdem werden die Wälder immer kleiner, weil man zu viele Bäume fällt. Doch die Wälder sind die Lungen der Erde: sie nehmen das CO\textsubscript{2} auf und geben uns Sauerstoff.

Die Folgen sind schon sichtbar. Die Erde wird wärmer, so dass die Gletscher in den Alpen und an den Polen schmelzen. Das Meer steigt, und Inseln im Pazifik könnten eines Tages verschwinden. In vielen Ländern Afrikas gibt es lange Dürren. Als der Regen letztes Jahr ausblieb, vertrockneten die Felder der Bauern, und die Ernte war sehr schlecht.

Aber es gibt Hoffnung. Wenn wir weniger Energie verschwendeten, hätten wir eine sauberere Umwelt. Wenn wir weniger Auto fuhren, hätten wir sauberere Luft. Wir könnten mehr Solar- und Windenergie nutzen, den Müll trennen und öfter den Bus oder das Fahrrad nehmen. Viele junge Leute demonstrieren schon für den Klimaschutz.

Wir müssen handeln, bevor es zu spät ist! Die Erde ist unser einziges Zuhause.
\end{textebox}

\begin{definition}[Glossaire du texte]
\begin{itemize}
\item \all{die Dürre, -n} = la sécheresse
\item \all{schmelzen (schmilzt, schmolz, ist geschmolzen)} = fondre (verbe fort)
\item \all{das Abgas, -e} = le gaz d'échappement
\item \all{der Gletscher, -} = le glacier
\item \all{fällen (fällt, fällte, hat gefällt)} = abattre (un arbre) (verbe faible)
\item \all{die Lunge, -n} = le poumon
\item \all{der Sauerstoff} = l'oxygène
\item \all{aufnehmen (nimmt auf, nahm auf, hat aufgenommen)} = absorber (verbe à particule séparable)
\item \all{ausbleiben (blieb aus, ist ausgeblieben)} = ne pas venir, faire défaut (verbe fort, particule séparable)
\item \all{vertrocknen (vertrocknet, vertrocknete, ist vertrocknet)} = se dessécher (verbe fort)
\item \all{die Ernte, -n} = la récolte
\item \all{verschwenden} = gaspiller (verbe faible)
\item \all{der Müll} = les déchets, les ordures
\item \all{der Klimaschutz} = la protection du climat
\end{itemize}
\end{definition}

\courssub{Les mots transparents : un trésor pour le lecteur}

Une grande partie du vocabulaire scientifique allemand est proche du français ou de l'anglais. Ces \cle{mots transparents} permettent de comprendre l'essentiel d'un texte d'actualité sans dictionnaire.

\begin{center}
\begin{tabularx}{\linewidth}{|X|X|X|}\hline
\rowcolor{popblueL} \textbf{Allemand} & \textbf{Français} & \textbf{Remarque} \\ \hline
\all{die Temperatur, -en} & la température & presque identique \\ \hline
\all{die Emission, -en} & l'émission & identique \\ \hline
\all{global} & global & identique \\ \hline
\all{die Energie, -n} & l'énergie & identique \\ \hline
\all{das Problem, -e} & le problème & attention : neutre ! \\ \hline
\all{die Fabrik, -en} & l'usine & faux ami : pas « fabrique » artisanale \\ \hline
\all{der Pazifik} & le Pacifique & proche \\ \hline
\all{die Diskussion, -en} & la discussion & identique \\ \hline
\all{die Katastrophe, -n} & la catastrophe & identique \\ \hline
\all{der Klimawandel} & le changement climatique & composé transparent : Klima + Wandel \\ \hline
\end{tabularx}
\end{center}

\begin{attention}[Faux amis et pièges de genre]
Méfie-toi des \cle{faux amis} : \all{die Fabrik} n'est pas « la fabrique » au sens artisanal mais bien « l'usine » ; \all{das Gift} signifie « le poison » et non « le cadeau ». Et surtout : \textbf{le genre ne se devine pas} ! \all{das Problem} est neutre, \all{die Erde} est féminin, \all{der Wald} est masculin. Apprends toujours chaque nom \textbf{avec son article et son pluriel} : \all{der Wald, die Wälder}.
\end{attention}

\courssub{Repérage de la structure du texte}

Après la deuxième lecture, on découpe le texte en quatre mouvements, annoncés par des mots-charnières :

\begin{center}
\begin{tabularx}{\linewidth}{|X|X|X|}\hline
\rowcolor{popblueL} \textbf{Partie} & \textbf{Phrases} & \textbf{Mot-charnière} \\ \hline
1. Le problème & § 1 : la température monte & \all{der Klimawandel ist real} \\ \hline
2. Les causes & § 2 : usines, voitures, déforestation & \all{Die Ursachen sind bekannt.} \\ \hline
3. Les conséquences & § 3 : glaciers, mer, sécheresses & \all{Die Folgen sind schon sichtbar.} \\ \hline
4. Les solutions & § 4--5 : énergies propres, gestes quotidiens & \all{Aber es gibt Hoffnung.} \\ \hline
\end{tabularx}
\end{center}

\begin{exemplebox}[Les articulations du texte]
Observe les connecteurs qui structurent l'article :
\begin{itemize}
\item \all{Außerdem werden die Wälder immer kleiner.} „De plus, les forêts deviennent de plus en plus petites.“ (\all{außerdem} ajoute une cause)
\item \all{Doch die Wälder sind die Lungen der Erde.} „Pourtant, les forêts sont les poumons de la Terre.“ (\all{doch} marque l'opposition)
\item \all{Aber es gibt Hoffnung.} „Mais il y a de l'espoir.“ (\all{aber} introduit le tournant vers les solutions)
\item \all{Wir müssen handeln, bevor es zu spät ist!} „Nous devons agir avant qu'il ne soit trop tard !“ (conclusion appelant à l'action)
\end{itemize}
\end{exemplebox}

\begin{popfigure}
\begin{center}
\begin{tikzpicture}[mindmap, grow cyclic, every node/.style={concept, align=center, font=\small}, root concept/.append style={concept color=popblue, font=\bfseries}, level 1/.append style={level distance=3.2cm, sibling angle=90}]
\node[root concept]{Klimawandel}
  child[concept color=poporange]{ node[align=center]{Ursachen\\ Fabriken\\ Autos\\ Abholzung} }
  child[concept color=poppink]{ node[align=center]{Folgen\\ Gletscher\\ schmelzen\\ Dürren\\ Meer steigt} }
  child[concept color=popgreen]{ node[align=center]{Lösungen\\ Solarenergie\\ Müll trennen\\ Bus, Fahrrad} }
  child[concept color=poppurple]{ node[align=center]{Appell\\ handeln,\\ bevor es\\ zu spät ist!} };
\end{tikzpicture}
\end{center}
\legende{Carte mentale du texte : les quatre mouvements de l'article.}
\end{popfigure}

\begin{popfigure}
\begin{center}
\begin{tikzpicture}[node distance=1.6cm, every node/.style={draw, rounded corners, align=center, font=\small, inner sep=6pt}, every arrow/.style={-{Stealth[length=3mm]}, thick, popdark}]
\node[fill=poporangeL, align=center] (co2){CO\textsubscript{2}\\ (Fabriken, Autos)};
\node[fill=poppinkL, right=of co2, align=center] (warm){Erderwärmung\\ (Temperatur steigt)};
\node[fill=popblueL, right=of warm, yshift=1.1cm, align=center] (duerre){Dürre\\ (Felder vertrocknen)};
\node[fill=popgreenL, right=of warm, yshift=-1.1cm, align=center] (gletscher){Gletscherschmelze\\ (Meer steigt)};
\draw (co2) -- (warm);
\draw (warm) -- (duerre);
\draw (warm) -- (gletscher);
\end{tikzpicture}
\end{center}
\legende{La chaîne logique du texte : du CO\textsubscript{2} aux conséquences visibles.}
\end{popfigure}

\courssub{Observation grammaticale : l'irréel et les subordonnées dans le texte}

Avant l'étude détaillée des sections 4 et 5, repérons les formes grammaticales cibles dans le texte.

\begin{exemplebox}[Le subjonctif II dans le texte]
\begin{itemize}
\item \all{Wenn wir weniger Energie verschwendeten, hätten wir eine sauberere Umwelt.} „Si nous gaspillions moins d'énergie, nous aurions un environnement plus propre.“
\item \all{Wenn wir weniger Auto fuhren, hätten wir sauberere Luft.} „Si nous roulions moins en voiture, nous aurions un air plus pur.“
\item \all{Wir könnten mehr Solarenergie nutzen.} „Nous pourrions utiliser davantage d'énergie solaire.“
\end{itemize}
On observe : \all{verschwendeten}, \all{fuhren} et \all{hätten} ont des voyelles modifiées (\all{a} $\rightarrow$ \all{ä}, \all{u} $\rightarrow$ \all{ü}) et expriment une hypothèse \textbf{irréelle} : ce n'est pas le cas aujourd'hui. En français, on utilise le couple imparfait + conditionnel ; en allemand, on utilise le \cle{Konjunktiv II} dans les deux parties de la phrase (proposition \all{si} et proposition principale).
\end{exemplebox}

\begin{exemplebox}[Les subordonnées de conséquence et de temps dans le texte]
\begin{itemize}
\item \all{Die Erde wird wärmer, so dass die Gletscher schmelzen.} „La Terre se réchauffe, de sorte que les glaciers fondent.“ (conséquence : \all{so dass} + indicatif, fait réel)
\item \all{Als der Regen letztes Jahr ausblieb, vertrockneten die Felder.} „Quand la pluie a fait défaut l'an dernier, les champs se sont desséchés.“ (temps passé unique : \all{als} + indicatif prétérit)
\item \all{Wir müssen handeln, bevor es zu spät ist.} „Nous devons agir avant qu'il ne soit trop tard.“ (antériorité : \all{bevor} + indicatif présent)
\end{itemize}
Règle d'or déjà visible : dans toute subordonnée allemande, le \textbf{verbe conjugué va à la fin} : \all{..., so dass die Gletscher \textbf{schmelzen}} ; \all{..., bevor es zu spät \textbf{ist}} ; \all{..., als der Regen \textbf{ausblieb}}.
\end{exemplebox}

\begin{exoresolu}[Compréhension globale du texte]
Énoncé : réponds en français, en t'appuyant sur le texte. 1) Quel est le thème général de l'article ? 2) Cite deux causes du réchauffement mentionnées. 3) Quelle conséquence concrète a touché les paysans l'an dernier ? 4) Quelles solutions l'article propose-t-il ?
\tcblower
Solution détaillée : 1) Le thème général est le réchauffement climatique (\all{der Klimawandel}) : ses causes humaines, ses conséquences déjà visibles et les solutions possibles. 2) Deux causes : les usines et les voitures qui produisent beaucoup de CO\textsubscript{2} (\all{Die Fabriken und Autos produzieren viel CO\textsubscript{2}}), et la déforestation (\all{man fällt zu viele Bäume}). 3) La pluie a fait défaut et les champs des paysans se sont desséchés, ce qui a ruiné la récolte (\all{Als der Regen ausblieb, vertrockneten die Felder der Bauern, und die Ernte war sehr schlecht}). 4) Les solutions : gaspiller moins d'énergie, utiliser les énergies solaire et éolienne, trier les déchets, prendre le bus ou le vélo, et manifester pour la protection du climat.
\end{exoresolu}

\begin{exercice}[À toi de jouer]
\begin{enumerate}
\item Retrouve dans le texte trois mots transparents et traduis-les.
\item Recopie la phrase au subjonctif II qui commence par \all{Wenn wir weniger Auto...} et souligne les deux verbes conjugués.
\item Relevez la subordonnée introduite par \all{so dass} et encadre le verbe conjugué placé en fin de phrase.
\end{enumerate}
\end{exercice}

\begin{corrige}[Exercice « À toi de jouer »]
\begin{enumerate}
\item Par exemple : \all{die Temperatur} = la température ; \all{die Emission} = l'émission ; \all{global} = global ; \all{die Katastrophe} = la catastrophe.
\item \all{Wenn wir weniger Auto \textbf{fuhren}, \textbf{hätten} wir sauberere Luft.}
\item \all{Die Erde wird wärmer, so dass die Gletscher \textbf{schmelzen}.} — le verbe conjugué \all{schmelzen} occupe bien la dernière place de la subordonnée.
\end{enumerate}
\end{corrige}

\begin{savaistu}[L'« Energiewende » en Allemagne]
L'Allemagne a décidé l'\all{Energiewende} (la « transition énergétique ») : sortie progressive du nucléaire et développement massif des énergies renouvelables — éoliennes dans le nord, panneaux solaires sur les toits. Le mot \all{Klimaschutz} (protection du climat) fait partie du débat public quotidien : on le voit dans les journaux, à la télévision et sur les pancartes des jeunes manifestants du mouvement \all{Fridays for Future}. Le mot \all{Dürre} a même été élu « mot de l'année » en Allemagne après un été exceptionnellement sec : la langue allemande reflète directement les préoccupations écologiques de la société.
\end{savaistu}

\begin{aretenir}[L'essentiel de cette section]
\begin{itemize}
\item Un article allemand sur l'environnement suit souvent le plan : problème $\rightarrow$ causes $\rightarrow$ conséquences $\rightarrow$ solutions.
\item Première lecture : sens global, mots transparents, jamais de traduction mot à mot.
\item Vocabulaire-clé : \all{die Umweltverschmutzung} (la pollution), \all{die Erderwärmung} (le réchauffement), \all{das Klima}, \all{die Dürre}, \all{der Gletscher}, \all{der Müll}.
\item Le texte contient les deux cibles grammaticales de la leçon : le \cle{Konjunktiv II} (l'irréel) et les subordonnées avec \all{so dass}, \all{als}, \all{bevor} — verbe conjugué en fin de subordonnée.
\end{itemize}
\end{aretenir}

\coursec{Compréhension du texte}

Dans la section précédente, nous avons lu et écouté le texte \all{„Die Erde wird wärmer“}, un article de presse qui décrit les causes et les conséquences du réchauffement climatique, en Allemagne comme au Cameroun. Nous allons maintenant vérifier et approfondir notre compréhension : c'est exactement le travail demandé à l'examen, où les questions de compréhension (\all{Verständnisfragen}) représentent une part importante de l'épreuve.

\courssub{Compréhension globale}

Avant de répondre aux questions de détail, il faut toujours se poser trois questions générales. Elles constituent la grille de lecture de tout texte au Bac.

\begin{methode}[Les trois questions de la compréhension globale]
\begin{enumerate}
\item \textbf{De quoi parle le texte ?} (\all{Worum geht es?}) — Identifier le thème en une phrase.
\item \textbf{À qui s'adresse-t-il ?} (\all{An wen richtet sich der Text?}) — Le public visé : lecteurs d'un journal, jeunes, décideurs politiques.
\item \textbf{Quel est le but de l'auteur ?} (\all{Was will der Autor?}) — Informer (\all{informieren}), alerter (\all{warnen}), convaincre (\all{überzeugen}).
\end{enumerate}
\end{methode}

Appliquons cette grille à notre texte :

\begin{center}
\begin{tabularx}{\linewidth}{|X|X|}
\hline
\rowcolor{popblueL} \textbf{Question} & \textbf{Réponse} \\
\hline
\all{Worum geht es?} & Il s'agit du réchauffement climatique : ses causes (gaz à effet de serre, déforestation), ses conséquences (sécheresses, fonte des glaciers) et les solutions possibles. \\
\hline
\all{An wen richtet sich der Text?} & À un large public de lecteurs de presse, et en particulier aux jeunes, invités à agir. \\
\hline
\all{Was will der Autor?} & Informer d'abord, mais surtout alerter et convaincre : le ton devient urgent à la fin (\all{„Wir müssen handeln, bevor es zu spät ist.“}). \\
\hline
\end{tabularx}
\end{center}

\begin{aretenir}[L'intention de l'auteur]
Un texte journalistique sur l'environnement poursuit presque toujours un double but : \cle{informer} (donner des faits, des explications) et \cle{convaincre} (pousser le lecteur à changer de comportement). Repérez les indices : le vocabulaire de l'urgence (\all{dringend}, \all{sofort}, \all{bevor es zu spät ist}), les impératifs et les phrases avec \all{müssen} ou \all{sollen}.
\end{aretenir}

\courssub{Questions de compréhension détaillée}

Voici les questions classiques posées sur ce texte, avec leurs réponses rédigées en phrases complètes. Observez d'abord la technique, puis la règle.

\begin{exemplebox}[Questions et réponses modèles]
\textbf{1.} \all{Warum steigt die Temperatur auf der Erde?}\\
\all{Die Temperatur steigt, weil die Menschen zu viel CO$_2$ produzieren, zum Beispiel mit Autos, Fabriken und Flugzeugen.}\\
« La température augmente parce que les hommes produisent trop de CO$_2$, par exemple avec les voitures, les usines et les avions. »

\textbf{2.} \all{Welche Folgen der Erderwärmung nennt der Text?}\\
\all{Der Text nennt mehrere Folgen: die Gletscher schmelzen, der Meeresspiegel steigt, und in vielen Ländern gibt es Dürren, weil der Regen ausbleibt.}\\
« Le texte mentionne plusieurs conséquences : les glaciers fondent, le niveau de la mer monte, et dans beaucoup de pays il y a des sécheresses parce que la pluie fait défaut. »

\textbf{3.} \all{Was passierte, als der Regen ausblieb?}\\
\all{Als der Regen ausblieb, vertrockneten die Felder, und die Bauern konnten nichts mehr ernten.}\\
« Quand la pluie a fait défaut, les champs ont séché et les paysans n'ont plus rien pu récolter. »

\textbf{4.} \all{Was schlägt der Autor vor?}\\
\all{Der Autor schlägt vor, weniger Energie zu verbrauchen, mehr Bäume zu pflanzen und öffentliche Verkehrsmittel zu benutzen.}\\
« L'auteur propose de consommer moins d'énergie, de planter plus d'arbres et d'utiliser les transports en commun. »
\end{exemplebox}

\begin{methode}[Répondre à une question de compréhension en quatre étapes]
\begin{enumerate}
\item \textbf{Repérer} dans la question les mots-clés (\all{warum}, \all{welche Folgen}, \all{was passierte}) et localiser le passage correspondant dans le texte.
\item \textbf{Reprendre les mots de la question} dans le début de la réponse : \all{Warum schmelzen die Gletscher?} $\rightarrow$ \all{Die Gletscher schmelzen, weil ...}
\item \textbf{Reformuler} avec ses propres mots : ne pas recopier la phrase du texte, mais la dire autrement (synonymes, autre structure).
\item \textbf{Rédiger une phrase complète} : sujet + verbe conjugué en deuxième position, verbe à la fin dans la subordonnée.
\end{enumerate}
\end{methode}

\begin{popfigure}
\begin{tikzpicture}[
  node distance=0.55cm,
  box/.style={draw=popblue, fill=popblueL, rounded corners=3pt, align=center, text width=4.6cm, minimum height=1.1cm, font=\small},
  arr/.style={-{Stealth[length=3mm]}, thick, poporange}
]
\node[box, align=center] (q){\textbf{1. Question}\\ \all{Warum steigt die Temperatur?}};
\node[box, below=of q, align=center] (r){\textbf{2. Repérage}\\ trouver le passage du texte};
\node[box, below=of r, align=center] (f){\textbf{3. Reformulation}\\ dire autrement, avec ses mots};
\node[box, below=of f, align=center] (p){\textbf{4. Phrase complète}\\ \all{Die Temperatur steigt, weil ...}};
\draw[arr] (q) -- (r);
\draw[arr] (r) -- (f);
\draw[arr] (f) -- (p);
\end{tikzpicture}
\legende{Organigramme de la méthode de réponse à une question de compréhension}
\end{popfigure}

\begin{attention}[Ne recopiez pas le texte !]
À l'oral comme à l'écrit, évitez de recopier de longs passages du texte sans reformulation : l'examinateur valorise la \cle{reformulation}. Recopier prouve que vous avez trouvé le passage, mais pas que vous l'avez compris. Gardez les mots de la question (obligatoire pour construire la phrase), mais changez ceux du texte : \all{ausbleiben} $\rightarrow$ \all{nicht kommen} ; \all{vertrocknen} $\rightarrow$ \all{trocken werden}.
\end{attention}

\courssub{Exemples de reformulation : observation}

Observez ces deux phrases du texte et leur reformulation :

\begin{exemplebox}[Du texte à la reformulation]
Phrase du texte : \all{Als der Regen ausblieb, vertrockneten die Felder.}\\
« Quand la pluie a fait défaut, les champs ont séché. »\\
Reformulation : \all{Weil es nicht regnete, wurden die Felder trocken.}

Phrase du texte : \all{Wir müssen handeln, bevor es zu spät ist.}\\
« Nous devons agir avant qu'il ne soit trop tard. »\\
Reformulation : \all{Der Autor meint, dass wir sofort etwas tun sollen.}
\end{exemplebox}

Remarquez les transformations grammaticales : la subordonnée de temps avec \all{als} (passé, fait unique) peut devenir une subordonnée de cause avec \all{weil} ; une phrase principale avec \all{müssen} peut devenir une subordonnée avec \all{dass} + \all{sollen}. Dans tous les cas, le verbe conjugué passe à la fin de la subordonnée : c'est le point que le correcteur surveille en priorité.

\begin{propriete}[\all{als} ou \all{wenn} ? Ne les confondez plus]
\begin{itemize}
\item \all{als} : passé, événement \textbf{unique}. \all{Als der Regen ausblieb, vertrockneten die Felder.} « Quand la pluie a fait défaut (cette année-là), les champs ont séché. »
\item \all{wenn} : présent, futur, ou passé \textbf{répété} (habitude). \all{Wenn der Regen ausbleibt, vertrocknen die Felder.} « Quand la pluie fait défaut (chaque fois), les champs sèchent. »
\item Dans les deux cas, le verbe conjugué va à la fin de la subordonnée ; si la subordonnée est en tête, le verbe de la principale vient juste après la virgule (inversion) : \all{Als der Regen ausblieb, \textbf{vertrockneten} die Felder.}
\end{itemize}
\end{propriete}

\courssub{Vrai ou faux justifié : l'exercice type Bac}

Au Bac, l'exercice \all{Richtig oder Falsch?} exige non seulement de juger, mais de \textbf{justifier par une citation courte du texte} entre guillemets allemands.

\begin{exoresolu}[Richtig oder Falsch? Begründen Sie mit einem Zitat]
Énoncé. Dites si les affirmations suivantes sont vraies ou fausses et justifiez par une courte citation du texte.

a) \all{Die Erderwärmung hat nur Folgen für Europa.}

b) \all{Der Autor glaubt, dass man noch etwas ändern kann.}

c) \all{Die Bauern konnten trotz der Dürre gut ernten.}
\tcblower
Solution détaillée.

a) \textbf{Falsch.} Le texte parle aussi de l'Afrique : \all{„In vielen Ländern Afrikas bleibt der Regen aus.“} La fausseté est signalée par le mot \all{nur} (« seulement »), toujours suspect dans ce type d'exercice.

b) \textbf{Richtig.} La fin du texte est un appel à l'action : \all{„Wir müssen handeln, bevor es zu spät ist.“} Si l'auteur appelle à agir, c'est qu'il estime le changement encore possible.

c) \textbf{Falsch.} Le texte dit le contraire : \all{„Als der Regen ausblieb, vertrockneten die Felder, und die Bauern konnten nichts mehr ernten.“} Le mot \all{trotz} (« malgré ») inverse le sens : méfiez-vous des petits mots qui retournent la phrase.
\end{exoresolu}

\begin{attention}[Les pièges du vrai/faux]
Les affirmations fausses contiennent presque toujours un mot qui modifie le sens du texte : \all{nur} (seulement), \all{nie} (jamais), \all{alle} (tous), \all{trotz} (malgré), \all{ohne} (sans). Soulignez ces mots avant de juger, et citez toujours le passage exact qui prouve votre réponse.
\end{attention}

\courssub{Recherche dans le texte : familles de mots et connecteurs}

Deux exercices de repérage entraînent l'œil à naviguer dans le texte.

\textbf{1. La famille de \all{die Umwelt}.} Cherchez dans le texte tous les mots construits sur \all{Umwelt} ou appartenant au même champ lexical :

\begin{center}
\begin{tabularx}{\linewidth}{|X|X|X|}
\hline
\rowcolor{popblueL} \textbf{Mot du texte} & \textbf{Français} & \textbf{Formation} \\
\hline
\all{die Umwelt} & l'environnement & \all{um} (autour) + \all{Welt} (monde) \\
\hline
\all{die Umweltverschmutzung} & la pollution de l'environnement & \all{Umwelt} + \all{verschmutzen} (salir) \\
\hline
\all{der Umweltschutz} & la protection de l'environnement & \all{Umwelt} + \all{schützen} (protéger) \\
\hline
\all{umweltfreundlich} & écologique & \all{Umwelt} + \all{freundlich} (ami) \\
\hline
\end{tabularx}
\end{center}

\textbf{2. Les connecteurs de temps et de conséquence.} Relevez dans le texte les mots qui organisent le récit et l'argumentation :

\begin{center}
\begin{tabularx}{\linewidth}{|X|X|X|}
\hline
\rowcolor{popblueL} \textbf{Connecteur} & \textbf{Français} & \textbf{Exemple du texte} \\
\hline
\all{als} & quand (passé, fait unique) & \all{Als der Regen ausblieb, ...} \\
\hline
\all{bevor} & avant que & \all{..., bevor es zu spät ist.} \\
\hline
\all{seit} & depuis & \all{Seit hundert Jahren steigt die Temperatur.} \\
\hline
\all{deshalb} & c'est pourquoi & \all{Deshalb müssen wir handeln.} \\
\hline
\all{so dass} & si bien que & \all{Es wurde so trocken, so dass die Felder vertrockneten.} \\
\hline
\end{tabularx}
\end{center}

\begin{propriete}[Connecteurs : attention à la place du verbe]
\begin{itemize}
\item \all{als}, \all{bevor}, \all{weil}, \all{so dass} introduisent une \textbf{subordonnée} : le verbe conjugué va \textbf{à la fin}. \all{Als der Regen \textbf{ausblieb}, ...}
\item \all{deshalb}, \all{darum}, \all{trotzdem} sont des \textbf{adverbes} en tête de phrase principale : le verbe reste en \textbf{deuxième position} et le sujet suit le verbe (inversion). \all{\textbf{Deshalb müssen wir} handeln.} (et non \all{deshalb wir müssen}).
\item \all{seit} est ici une \textbf{préposition} (+ datif) : \all{seit hundert Jahren} ; elle n'introduit pas de verbe. Avec un verbe, on utilise la conjonction \all{seitdem} : \all{Seitdem die Temperatur steigt, schmelzen die Gletscher.}
\end{itemize}
\end{propriete}

\courssub{Texte complémentaire : entraînement à la compréhension}

Pour vous entraîner seul, voici un court texte sur le même thème, suivi de questions. Appliquez la méthode en quatre étapes.

\begin{textebox}[„Ein Dorf im Norden Kameruns“]
\all{In einem kleinen Dorf im Norden Kameruns war das Leben früher einfacher. Die Regenzeit kam jedes Jahr pünktlich, und die Bauern konnten Mais und Hirse anbauen. Doch seit einigen Jahren ist alles anders: Der Regen kommt später und bleibt manchmal ganz aus. Die Erde ist trocken, der Fluss hat wenig Wasser, und viele junge Leute verlassen das Dorf, weil sie keine Arbeit mehr finden.}

\all{Aminou, ein Bauer aus dem Dorf, erzählt: „Früher haben wir genug geerntet. Heute müssen wir weit laufen, um Wasser zu finden.“ Trotzdem gibt er die Hoffnung nicht auf. Mit einer Gruppe von Nachbarn hat er begonnen, Bäume zu pflanzen, denn Bäume halten den Boden feucht und geben Schatten. „Wenn jeder von uns nur einen Baum pro Jahr pflanzt“, sagt er, „wird unser Dorf in zehn Jahren wieder grün sein.“}

\all{Auch die Schule des Dorfes macht mit: Die Kinder lernen, wie man Wasser spart und den Müll richtig wegräumt. Vielleicht kann dieses kleine Dorf ein Beispiel für die ganze Region werden.}
\end{textebox}

\textbf{Glossaire :} \all{pünktlich} (ponctuel, à l'heure) ; \all{der Mais} (le maïs) ; \all{die Hirse, -n} (le mil) ; \all{anbauen} (cultiver) ; \all{die Ernte, -n / ernten} (la récolte / récolter) ; \all{feucht} (humide) ; \all{der Schatten, -} (l'ombre) ; \all{der Müll} (les ordures) ; \all{wegräumen} (enlever) ; \all{die Hoffnung aufgeben} (abandonner l'espoir).

\begin{exercice}[Fragen zum Text]
Répondez en allemand, par des phrases complètes, en reformulant.

1. \all{Warum verlassen viele junge Leute das Dorf?}

2. \all{Was macht Aminou mit seinen Nachbarn, und warum?}

3. \all{Was lernen die Kinder in der Schule?}

4. \all{Richtig oder falsch? Begründen Sie: „Aminou hat keine Hoffnung mehr.“}
\end{exercice}

\begin{corrige}[Exercice : Fragen zum Text]
1. \all{Viele junge Leute verlassen das Dorf, weil es wegen der Trockenheit keine Arbeit mehr gibt.} (Reformulation de \all{der Regen bleibt aus} par \all{Trockenheit}.)

2. \all{Aminou pflanzt mit seinen Nachbarn Bäume, weil Bäume den Boden feucht halten und Schatten geben.}

3. \all{In der Schule lernen die Kinder, Wasser zu sparen und den Müll richtig wegzuräumen.}

4. \textbf{Falsch.} \all{„Trotzdem gibt er die Hoffnung nicht auf.“} Le mot \all{trotzdem} montre qu'il garde espoir malgré les difficultés.
\end{corrige}

\courssub{Le résumé guidé}

Dernière étape de la compréhension : savoir résumer. On procède en deux temps.

\textbf{Étape 1 : résumé à trous.} Complétez avec les mots suivants : \all{Folgen — handeln — CO$_2$ — Regen — Bäume}.

\all{Die Temperatur der Erde steigt, weil die Menschen zu viel \dots\ produzieren. Der Text nennt viele \dots\ der Erderwärmung: Die Gletscher schmelzen, und wenn der \dots\ ausbleibt, vertrocknen die Felder. Der Autor schlägt vor, weniger Energie zu verbrauchen und mehr \dots\ zu pflanzen. Er meint, wir müssen sofort \dots, bevor es zu spät ist.}

(Corrigé : \all{CO$_2$} ; \all{Folgen} ; \all{Regen} ; \all{Bäume} ; \all{handeln}.)

\textbf{Étape 2 : résumé libre en 3 ou 4 phrases.} Un bon résumé suit toujours le plan du texte : cause $\rightarrow$ conséquences $\rightarrow$ solutions $\rightarrow$ appel.

\begin{exemplebox}[Modèle de résumé]
\all{Der Text erklärt, dass die Erde immer wärmer wird, weil die Menschen zu viel CO$_2$ produzieren. Als Folge schmelzen die Gletscher, und in vielen Ländern vertrocknen die Felder, wenn der Regen ausbleibt. Der Autor schlägt vor, Energie zu sparen und Bäume zu pflanzen. Er fordert alle Menschen auf, sofort zu handeln, bevor es zu spät ist.}
\end{exemplebox}

Observez le verbe à particule du modèle : \all{auffordern} (inviter, sommer) se conjugue \all{er fordert ... auf}, la particule \all{auf} allant à la fin de la phrase principale. C'est un excellent verbe de résumé pour exprimer l'appel de l'auteur.

\begin{savaistu}[Les trois niveaux des questions au Bac]
Au Bac, les questions de compréhension portent presque toujours sur trois niveaux : les \textbf{causes} (\all{Warum ...?}), les \textbf{conséquences} (\all{Welche Folgen ...?}) et l'\textbf{opinion de l'auteur} (\all{Was meint der Autor?}). Entraînez-vous, pour chaque texte lu, à repérer ces trois niveaux et à surligner les passages correspondants avec trois couleurs différentes : vous gagnerez un temps précieux le jour de l'examen.
\end{savaistu}

\begin{aretenir}[Ce qu'il faut retenir]
\begin{itemize}
\item Compréhension globale : thème, public, intention (informer, alerter, convaincre).
\item Réponse à une question : reprendre les mots de la question, reformuler, phrase complète, verbe à la fin dans les subordonnées.
\item \all{als} = passé unique ; \all{wenn} = présent, futur ou habitude ; \all{deshalb} = adverbe avec inversion.
\item Vrai/faux : toujours justifier par une courte citation ; se méfier de \all{nur}, \all{nie}, \all{alle}, \all{trotz}.
\item Résumé : suivre le plan du texte (cause, conséquences, solutions, appel) en 3-4 phrases.
\end{itemize}
\end{aretenir}

\coursec{Lexique thématique : Umwelt und Klima}

Après avoir lu et compris le texte \all{Die Erde wird wärmer}, nous avons rencontré de nombreux mots nouveaux : \all{die Umweltverschmutzung}, \all{die Erderwärmung}, \all{die Gletscher}, \all{die Dürre}... Cette section a pour but de les organiser, de les expliquer et de les mémoriser durablement. En allemand, le vocabulaire n'est jamais une simple liste : chaque nom possède un \cle{genre} (der, die, das) et un \cle{pluriel} qu'il faut apprendre en même temps que le mot. C'est la règle d'or du germaniste : on n'apprend jamais \all{Wald} tout seul, mais \all{der Wald, die Wälder}.

\courssub{Principes généraux : genre, pluriel et mots composés}

\begin{propriete}[Apprendre un nom allemand en trois temps]
Pour chaque nom allemand, on mémorise systématiquement :
\begin{enumerate}
\item \textbf{l'article} (le genre) : \all{der} (masculin), \all{die} (féminin), \all{das} (neutre) ;
\item \textbf{le pluriel} : \all{der Baum, die Bäume} ; \all{die Fabrik, die Fabriken} ; \all{das Abgas, die Abgase} ;
\item \textbf{un exemple en contexte} : \all{Die Fabriken produzieren viele Abgase.} « Les usines produisent beaucoup de gaz d'échappement. »
\end{enumerate}
Le genre allemand ne coïncide pas toujours avec le genre français : il faut donc l'apprendre mot par mot, sans se fier à l'intuition.
\end{propriete}

\begin{attention}[Le genre : ne pas copier le français]
\begin{itemize}
\item \all{das Klima} est \textbf{neutre}, alors que « le climat » est masculin en français. On dit : \all{Das Klima verändert sich.} « Le climat change. »
\item \all{die Umwelt} est féminin comme « l'environnement »... mais \all{der Umweltschutz} est masculin, car dans un mot composé, c'est toujours le \textbf{dernier élément} qui donne le genre : \all{die Umwelt} + \all{der Schutz} $\rightarrow$ \all{der Umweltschutz}.
\item \all{der Wald} a un pluriel avec tréma : \all{die Wälder} (et non \all{die Walds}).
\item Ne pas confondre \all{die Dürre (-n)} « la sécheresse » et \all{die Überschwemmung (-en)} « l'inondation » : ce sont deux catastrophes opposées, toutes deux féminines.
\end{itemize}
\end{attention}

\begin{savaistu}[L'allemand, langue des mots composés]
L'allemand adore coller les mots les uns aux autres pour créer des mots nouveaux. Ainsi \all{die Umweltverschmutzung} = \all{Umwelt} (environnement) + \all{Verschmutzung} (pollution). Autres exemples de notre thème : \all{die Luftverschmutzung} (Luft + Verschmutzung, « pollution de l'air »), \all{das Treibhausgas} (Treibhaus « serre » + Gas), \all{die Erderwärmung} (Erde « Terre » + Erwärmung « réchauffement »), \all{die Windkraft} (Wind + Kraft « force »). Savoir \cle{découper} un mot composé permet de deviner son sens au baccalauréat, même sans dictionnaire : cherchez toujours le dernier élément, c'est lui qui porte le sens principal et le genre.
\end{savaistu}

\courssub{Champ lexical 1 : die Verschmutzung (la pollution)}

\begin{definition}[Vocabulaire de la pollution]
\begin{itemize}
\item \all{die Umweltverschmutzung (-)} : la pollution de l'environnement (mot sans pluriel courant)
\item \all{die Luftverschmutzung (-)} : la pollution de l'air
\item \all{das Abgas (-e)} : le gaz d'échappement ; pluriel \all{die Abgase}
\item \all{der Müll (-s)} : les ordures, les déchets
\item \all{der Abfall (¨-e)} : le déchet ; pluriel \all{die Abfälle}
\item \all{die Fabrik (-en)} : l'usine
\item \all{der Rauch (-s)} : la fumée
\item \all{verschmutzen (verschmutzt, verschmutzte, hat verschmutzt)} : polluer (verbe faible)
\end{itemize}
\end{definition}

\begin{exemplebox}[La pollution en contexte]
\all{Die Fabriken in Douala produzieren viel Rauch und viele Abgase.} « Les usines de Douala produisent beaucoup de fumée et de gaz d'échappement. »

\all{Die Autos verschmutzen die Luft in den großen Städten.} « Les voitures polluent l'air dans les grandes villes. »

\all{Am Strand von Kribi liegt oft viel Müll.} « Sur la plage de Kribi, il y a souvent beaucoup de déchets. »
\end{exemplebox}

\courssub{Champ lexical 2 : das Klima und die Erderwärmung}

\begin{definition}[die Erderwärmung]
\all{Die Erderwärmung (-)} désigne le réchauffement de la planète causé surtout par les émissions de \all{CO₂} (dioxyde de carbone). En français : \emph{le réchauffement climatique}. On dit aussi \all{der Klimawandel} « le changement climatique ».
\end{definition}

\begin{definition}[Vocabulaire du climat]
\begin{itemize}
\item \all{das Klima (-s / Klimata)} : le climat
\item \all{der Klimawandel (-)} : le changement climatique
\item \all{die Temperatur (-en)} : la température
\item \all{das CO₂} : le dioxyde de carbone (on prononce « tse-o-tsvaï »)
\item \all{das Treibhausgas (-e)} : le gaz à effet de serre
\item \all{der Treibhauseffekt (-e)} : l'effet de serre
\item \all{steigen (steigt, stieg, ist gestiegen)} : monter, augmenter — verbe fort, auxiliaire \all{sein}
\end{itemize}
\end{definition}

\begin{exemplebox}[Le climat en contexte]
\all{Die Temperaturen steigen jedes Jahr.} « Les températures augmentent chaque année. »

\all{Die Temperatur ist im letzten Jahrhundert gestiegen.} « La température a augmenté au siècle dernier. » (Perfekt avec \all{sein} : changement d'état.)

\all{CO₂ ist ein Treibhausgas.} « Le CO₂ est un gaz à effet de serre. »
\end{exemplebox}

\begin{attention}[steigen : verbe fort avec sein]
\all{steigen} est un verbe fort irrégulier : \all{steigt} (présent), \all{stieg} (prétérit), \all{ist gestiegen} (parfait). Il exprime un changement d'état, donc il se conjugue avec \all{sein} au Perfekt : \all{Die Temperatur \textbf{ist} gestiegen.} et non \all{hat gestiegen}. Même famille : \all{die Erde wird wärmer} « la Terre devient plus chaude ».
\end{attention}

\courssub{Champ lexical 3 : die Natur in Gefahr (la nature menacée)}

\begin{definition}[Vocabulaire de la nature menacée]
\begin{itemize}
\item \all{der Wald (¨-er)} : la forêt ; pluriel \all{die Wälder}
\item \all{der Baum (¨-e)} : l'arbre ; pluriel \all{die Bäume}
\item \all{der Gletscher (-)} : le glacier (pluriel identique : \all{die Gletscher})
\item \all{das Meer (-e)} : la mer
\item \all{die Dürre (-n)} : la sécheresse
\item \all{die Überschwemmung (-en)} : l'inondation
\item \all{schmelzen (schmilzt, schmolz, ist geschmolzen)} : fondre (verbe fort, auxiliaire \all{sein})
\item \all{aussterben (stirbt aus, starb aus, ist ausgestorben)} : s'éteindre, disparaître (espèces ; verbe fort, auxiliaire \all{sein})
\end{itemize}
\end{definition}

\begin{exemplebox}[Exemples traduits]
\all{Die Gletscher schmelzen, weil die Erde wärmer wird.} « Les glaciers fondent parce que la Terre se réchauffe. »

\all{Im Norden Kameruns gibt es oft Dürren.} « Dans le Nord du Cameroun, il y a souvent des sécheresses. »

\all{Viele Tiere sterben aus, weil die Wälder kleiner werden.} « Beaucoup d'animaux s'éteignent parce que les forêts deviennent plus petites. »

\all{Nach starkem Regen gibt es in Douala oft Überschwemmungen.} « Après de fortes pluies, il y a souvent des inondations à Douala. »
\end{exemplebox}

\courssub{Champ lexical 4 : die Lösungen (les solutions)}

\begin{definition}[Vocabulaire des solutions]
\begin{itemize}
\item \all{die erneuerbaren Energien (pl.)} : les énergies renouvelables
\item \all{die Solarenergie (-n)} : l'énergie solaire
\item \all{die Windkraft (-)} : l'énergie éolienne
\item \all{recyceln (recycelt, recycelte, hat recycelt)} : recycler
\item \all{den Müll trennen} : trier les déchets
\item \all{sparen (spart, sparte, hat gespart)} : économiser
\item \all{schützen (schützt, schützte, hat geschützt)} : protéger
\item \all{der Umweltschutz (-)} : la protection de l'environnement
\item \all{die öffentlichen Verkehrsmittel (pl.)} : les transports en commun
\end{itemize}
\end{definition}

\begin{exemplebox}[Expressions utiles à réutiliser à l'écrit et à l'oral]
\all{Wir müssen die Umwelt schützen.} « Nous devons protéger l'environnement. »

\all{Alle Länder müssen gegen den Klimawandel kämpfen.} « Tous les pays doivent lutter contre le changement climatique. »

\all{Wir sollen Energie sparen und den Müll trennen.} « Nous devrions économiser l'énergie et trier les déchets. »

\all{In Yaoundé benutzen viele Schüler öffentliche Verkehrsmittel.} « À Yaoundé, beaucoup d'élèves utilisent les transports en commun. »
\end{exemplebox}

\courssub{Tableau récapitulatif du vocabulaire}

\begin{center}
\begin{tabularx}{\linewidth}{|l|l|X|}
\hline
\rowcolor{popblueL} \textbf{Deutsch} & \textbf{Français} & \textbf{Exemple} \\
\hline
die Umweltverschmutzung (-) & la pollution de l'environnement & \all{Die Umweltverschmutzung ist ein großes Problem.} \\
\hline
das Abgas (-e) & le gaz d'échappement & \all{Abgase verschmutzen die Luft.} \\
\hline
der Müll (-s) & les déchets & \all{Wir müssen den Müll trennen.} \\
\hline
das Klima & le climat & \all{Das Klima verändert sich.} \\
\hline
die Erderwärmung (-) & le réchauffement climatique & \all{Die Erderwärmung ist gefährlich.} \\
\hline
das Treibhausgas (-e) & le gaz à effet de serre & \all{CO₂ ist ein Treibhausgas.} \\
\hline
der Wald (¨-er) & la forêt & \all{Die Wälder werden kleiner.} \\
\hline
der Gletscher (-) & le glacier & \all{Die Gletscher schmelzen.} \\
\hline
die Dürre (-n) & la sécheresse & \all{Die Dürre zerstört die Ernte.} \\
\hline
die Überschwemmung (-en) & l'inondation & \all{Die Überschwemmung zerstört Häuser.} \\
\hline
die Solarenergie (-n) & l'énergie solaire & \all{Solarenergie ist sauber.} \\
\hline
der Umweltschutz (-) & la protection de l'environnement & \all{Umweltschutz ist wichtig.} \\
\hline
\end{tabularx}
\end{center}

\courssub{Carte mentale : le champ lexical de l'environnement}

\begin{popfigure}
\begin{tikzpicture}[mindmap, grow cyclic,
  every node/.style={concept, text=white, font=\small, align=center},
  root concept/.append style={concept color=popblue, font=\bfseries},
  level 1/.append style={level distance=3.6cm, sibling angle=90},
  level 1 concept/.append style={font=\small\bfseries}]
\node[root concept]{Umwelt}
  child[concept color=poporange]{ node {Verschmutzung}
    child[level distance=2.6cm, sibling angle=60, concept color=poporange!70]{ node[concept, font=\scriptsize, align=center]{das Abgas\\der Müll\\der Rauch} }
  }
  child[concept color=popgreen]{ node {Klima}
    child[level distance=2.6cm, sibling angle=60, concept color=popgreen!70]{ node[concept, font=\scriptsize, align=center]{die Erderwärmung\\das CO₂\\die Temperatur} }
  }
  child[concept color=poppurple]{ node {Natur in Gefahr}
    child[level distance=2.6cm, sibling angle=60, concept color=poppurple!70]{ node[concept, font=\scriptsize, align=center]{der Wald\\der Gletscher\\die Dürre} }
  }
  child[concept color=poppink]{ node {Lösungen}
    child[level distance=2.6cm, sibling angle=60, concept color=poppink!70]{ node[concept, font=\scriptsize, align=center]{recyceln\\sparen\\die Solarenergie} }
  };
\end{tikzpicture}
\legende{Carte mentale du champ lexical : quatre familles autour du mot \all{Umwelt}.}
\end{popfigure}

\courssub{Familles de mots : construire son vocabulaire}

En allemand, un même radical donne naissance à des verbes, des noms et des adjectifs. Connaître ces familles multiplie votre vocabulaire sans effort.

\begin{center}
\begin{tabularx}{\linewidth}{|X|X|X|}
\hline
\rowcolor{popblueL} \textbf{Verbe} & \textbf{Nomen} & \textbf{Adjektiv / Beispiel} \\
\hline
\all{schützen} (protéger) & \all{der Schutz} / \all{der Umweltschutz} & \all{Wir schützen die Umwelt.} \\
\hline
\all{verschmutzen} (polluer) & \all{die Verschmutzung} & \all{verschmutzt} : \all{die verschmutzte Luft} „l'air pollué“ \\
\hline
\all{wärmen} (chauffer) & \all{die Wärme} / \all{die Erderwärmung} & \all{warm} : \all{Die Erde wird wärmer.} \\
\hline
\all{steigen} (monter) & \all{der Anstieg} (la hausse) & \all{Die Temperaturen steigen.} \\
\hline
\all{schmelzen} (fondre) & \all{das Schmelzen} & \all{Die Gletscher schmelzen.} \\
\hline
\all{recyceln} (recycler) & \all{das Recycling} & \all{recycelbar} : \all{recycelbares Material} \\
\hline
\all{anbauen} (cultiver) & \all{der Anbau} & \all{Der Anbau von Bäumen hilft.} \\
\hline
\end{tabularx}
\end{center}

\begin{propriete}[Les verbes à particule séparable du thème]
Certains verbes de ce champ lexical sont \cle{séparables} : à l'indicatif présent, la particule se détache du verbe et va \textbf{en fin de proposition}.
\begin{itemize}
\item \all{aussterben} (s'éteindre) : \all{Viele Tierarten \textbf{sterben} heute \textbf{aus}.} „Beaucoup d'espèces animales s'éteignent aujourd'hui.“
\item \all{ausbleiben} (faire défaut) : \all{Der Regen \textbf{bleibt} dieses Jahr \textbf{aus}.} „La pluie fait défaut cette année.“
\item \all{wegwerfen} (jeter) : \all{Man \textbf{wirft} zu viel Plastik \textbf{weg}.} „On jette trop de plastique.“
\item \all{anbauen} (cultiver) : \all{Die Bauern \textbf{bauen} Mais \textbf{an}.} „Les paysans cultivent du maïs.“
\end{itemize}
Au Perfekt, le \all{ge-} se place entre la particule et le radical : \all{ist ausgestorben}, \all{ist ausgeblieben}, \all{hat weggeworfen}, \all{hat angebaut}.
\end{propriete}

\courssub{Texte de consolidation avec glossaire}

\begin{textebox}[Umweltschutz in Kamerun — ein Beispiel aus Yaoundé]
In Yaoundé sehen die Schüler jeden Tag die Probleme der Umweltverschmutzung. Auf den Straßen liegt oft Müll, und die alten Autos und Motorräder produzieren viele Abgase. Die Luft in der Stadt ist manchmal schlecht. Aber es gibt auch Hoffnung: in vielen Schulen lernen die Jugendlichen jetzt, wie man die Umwelt schützen kann. Sie pflanzen Bäume im Schulhof, sie trennen den Müll und sie sparen Wasser und Energie.

Auch der Klimawandel ist in Kamerun zu sehen. Im Norden des Landes gibt es lange Dürren: der Regen bleibt aus, und die Bauern können kein Gemüse mehr anbauen. In anderen Regionen, zum Beispiel in Douala, gibt es dagegen oft Überschwemmungen nach starkem Regen. Die Temperaturen steigen, und viele Menschen haben Angst vor der Zukunft.

Die Regierung und die Jugend wollen handeln. Solarenergie ist für Kamerun sehr interessant, denn die Sonne scheint fast das ganze Jahr. Wenn alle Menschen den Müll trennen, Energie sparen und öffentliche Verkehrsmittel benutzen, wird die Luft sauberer. Wir müssen die Umwelt schützen — für uns und für unsere Kinder. Die Erde ist unser einziges Zuhause.
\end{textebox}

\textbf{Glossaire :} \all{der Schulhof (¨-e)} la cour d'école ; \all{pflanzen} planter ; \all{der Bauer (-n)} le paysan ; \all{das Gemüse (-)} les légumes ; \all{anbauen} cultiver (séparable) ; \all{handeln} agir ; \all{scheinen (scheint, schien, hat geschienen)} briller ; \all{das Zuhause (-)} le foyer, la maison.

\begin{exercice}[Questions de compréhension du texte de consolidation]
\begin{enumerate}
\item \all{Welche Umweltprobleme sieht man in Yaoundé?}
\item \all{Was lernen die Schüler in den Schulen?}
\item \all{Warum ist Solarenergie für Kamerun interessant?}
\item \all{Was passiert im Norden Kameruns mit dem Regen?}
\item \all{Was schlägt der Text vor, damit die Luft sauberer wird?}
\end{enumerate}
\end{exercice}

\courssub{Exercice résolu et pièges}

\begin{exoresolu}[Complétez avec le bon article, le bon pluriel ou la bonne forme verbale]
Énoncé. Complétez :
1. \all{... Gletscher schmelzen, weil ... Erde wärmer wird.}
2. \all{... (das Klima / die Klima) verändert sich schnell.}
3. \all{Viele ... (Wald) sterben wegen der Dürre.}
4. \all{Wir müssen ... Müll trennen.}
5. \all{Viele Tierarten ... (aussterben, Präsens).}
6. \all{Der Regen ... (ausbleiben, Präteritum) diesen Sommer.}
\tcblower
Solution.
1. \all{\textbf{Die} Gletscher schmelzen, weil \textbf{die} Erde wärmer wird.} (\all{der Gletscher} $\rightarrow$ pluriel identique \all{die Gletscher} ; \all{die Erde} féminin ; \all{weil} renvoie le verbe à la fin).
2. \all{\textbf{Das} Klima verändert sich schnell.} (\all{das Klima} est neutre !).
3. \all{Viele \textbf{Wälder} sterben wegen der Dürre.} (Pluriel avec tréma : \all{der Wald, die Wälder} ; \all{wegen} + génitif \all{der Dürre}).
4. \all{Wir müssen \textbf{den} Müll trennen.} (Accusatif masculin : \all{der} $\rightarrow$ \all{den}).
5. \all{Viele Tierarten \textbf{sterben aus}.} (Verbe à particule séparable : la particule \all{aus} va en fin de phrase).
6. \all{Der Regen \textbf{blieb aus} diesen Sommer.} (Prétérit de \all{ausbleiben} : \all{blieb aus} ; auxiliaire \all{sein} au Perfekt : \all{ist ausgeblieben}).
\end{exoresolu}

\begin{attention}[Erreurs fréquentes des francophones]
\begin{itemize}
\item Genre : ne dites jamais \all{die Klima} ; c'est \all{das Klima}. De même \all{das Meer}, \all{das CO₂}, \all{das Abgas}, \all{das Treibhausgas} sont neutres.
\item Pluriels à tréma : \all{der Wald, die Wälder} ; \all{der Baum, die Bäume} ; \all{der Abfall, die Abfälle} ; \all{der Anbau, die Anbauten} (pas de tréma ici). Le tréma fait partie intégrante du pluriel pour ces mots.
\item Faux ami : \all{die Dürre} signifie « la sécheresse », pas « la dureté ». Et \all{die Überschwemmung} = « l'inondation », à ne pas confondre.
\item Particule séparable : à l'oral comme à l'écrit, n'oubliez jamais la particule finale : \all{Die Arten sterben aus.} et non \all{Die Arten aussterben.}
\item Auxiliaire Perfekt : \all{steigen}, \all{schmelzen}, \all{aussterben}, \all{ausbleiben} se conjuguent avec \all{sein} (\all{sind gestiegen, sind geschmolzen, sind ausgestorben, sind ausgeblieben}). \all{wegwerfen}, \all{recyceln}, \all{schützen} prennent \all{haben}.
\item Préposition \all{wegen} : toujours suivi du \textbf{génitif} : \all{wegen der Dürre}, \all{wegen des Klimawandels}.
\end{itemize}
\end{attention}

\begin{methode}[Méthode : Enrichir et réemployer le vocabulaire thématique]
Pour réussir l'expression écrite et orale sur ce thème :
\begin{enumerate}
\item \textbf{Apprenez par familles} : regroupez les mots par champ (pollution, climat, nature, solutions) et par famille morphologique (verbe / nom / adjectif).
\item \textbf{Systématisez article + pluriel} : écrivez \all{der Wald, die Wälder} sur vos fiches, jamais \all{Wald} seul.
\item \textbf{Repérez les verbes forts et leur auxiliaire} : notez \all{steigen – stieg – ist gestiegen} (sein), \all{schmelzen – schmolz – ist geschmolzen} (sein).
\item \textbf{Entraînez-vous aux particules séparables} : écrivez 3 phrases avec \all{aussterben}, \all{anbauen}, \all{wegwerfen} en plaçant la particule à la fin.
\item \textbf{Réutilisez les connecteurs logiques} : \all{weil} (verbe à la fin), \all{deshalb} / \all{deswegen} (verbe en 2e position), \all{wegen} + génitif, \all{wenn} (verbe à la fin) pour les conséquences et causes.
\end{enumerate}
\end{methode}

\begin{aretenir}[L'essentiel du lexique]
Quatre familles de mots : la \cle{pollution} (\all{die Umweltverschmutzung, das Abgas, der Müll}), le \cle{climat} (\all{das Klima, die Erderwärmung, das Treibhausgas}), la \cle{nature menacée} (\all{der Wald, der Gletscher, die Dürre, die Überschwemmung}) et les \cle{solutions} (\all{die Solarenergie, recyceln, der Umweltschutz}). Pour chaque nom : article + pluriel + exemple. Pour chaque verbe fort : les trois formes de base (Präsens, Präteritum, Partizip II) et l'auxiliaire (sein/haben). Ce vocabulaire est la boîte à outils indispensable pour le commentaire de texte et la production écrite de la fin du chapitre.
\end{aretenir}

\coursec{Grammaire 1 : le subjonctif II (l'irréel)}

Dans la section précédente, nous avons enrichi notre vocabulaire sur le thème \all{Umwelt und Klima} : nous savons maintenant parler de \all{die Umweltverschmutzung} (la pollution de l'environnement), de \all{die Erderwärmung} (le réchauffement climatique) et de \all{die Dürre} (la sécheresse). Mais parler de l'environnement, c'est aussi imaginer ce qui \emph{pourrait} être différent : « Si nous gaspillions moins d'énergie, nous aurions un environnement plus propre. » En français, nous utilisons l'imparfait et le conditionnel. En allemand, c'est le \cle{subjonctif II} (\all{der Konjunktiv II}) qui joue ce rôle. C'est l'une des formes les plus fréquentes de la langue parlée et écrite, et elle tombe régulièrement au baccalauréat, aussi bien en thème qu'en commentaire de texte. Observons-la d'abord dans notre texte.

\courssub{Observation : que remarque-t-on dans le texte ?}

Reprenons la phrase clé du texte support \all{„Die Erde wird wärmer“} :

\begin{exemplebox}[Observation guidée]
\all{Wenn wir weniger Energie verschwendeten, hätten wir eine sauberere Umwelt.}\\
« Si nous gaspillions moins d'énergie, nous aurions un environnement plus propre. »
\end{exemplebox}

Posons-nous trois questions :

\begin{enumerate}
\item \textbf{Quel temps utilise le français ?} L'imparfait dans la subordonnée (« si nous gaspillions ») et le conditionnel dans la principale (« nous aurions »).
\item \textbf{Quelle est la forme des verbes allemands ?} \all{verschwendeten} ressemble au Präteritum de \all{verschwenden} (verbe faible : \all{verschwendete}) ; \all{hätten} ressemble au Präteritum de \all{haben} (\all{hatten}), mais avec un \textbf{tréma} sur le \emph{a}.
\item \textbf{L'action est-elle réelle ?} Non ! Nous gaspillons toujours de l'énergie : c'est une \textbf{hypothèse irréelle}, un rêve, un « si seulement... ».
\end{enumerate}

\begin{aretenir}[Ce qu'il faut retenir de l'observation]
Le subjonctif II exprime l'\cle{irréel} : une hypothèse contraire à la réalité, un souhait, un conseil ou une demande polie. Sa forme ressemble au Präteritum, souvent avec un tréma pour les verbes forts. Le français lui oppose l'imparfait (après « si ») ou le conditionnel.
\end{aretenir}

\courssub{Les emplois du subjonctif II}

Le subjonctif II a quatre grands emplois qu'il faut savoir reconnaître et produire.

\textbf{1. L'hypothèse irréelle} (la réalité est différente) :

\all{Wenn ich reich wäre, würde ich einen Wald pflanzen.} « Si j'étais riche, je planteraient une forêt. » (Mais je ne suis pas riche !)

\all{Wenn wir mehr recycelten, gäbe es weniger Müll.} « Si nous recyclions davantage, il y aurait moins de déchets. »

\textbf{2. Le souhait irréel}, souvent introduit par \all{wenn ... doch} ou \all{wenn ... nur} (« si seulement... ») :

\all{Wenn es doch nicht so heiß wäre !} « Si seulement il ne faisait pas si chaud ! »

\all{Wenn es doch regnen würde !} « Si seulement il pleuvait ! »

\all{Hätten wir nur mehr Wasser !} « Si seulement nous avions plus d'eau ! » (souhait sans \all{wenn} : le verbe est en tête de phrase)

\textbf{3. Le conseil poli}, avec \all{sollte} ou \all{könnte} :

\all{Du solltest weniger Auto fahren.} « Tu devrais moins prendre la voiture. »

\all{Du solltest den Müll trennen.} « Tu devrais trier les déchets. »

\all{Wir könnten öfter das Fahrrad nehmen.} « Nous pourrions prendre plus souvent le vélo. »

\textbf{4. La demande polie}, dans les magasins, les restaurants, la vie quotidienne :

\all{Ich hätte gern ein Glas Wasser.} « J'aimerais un verre d'eau. »

\all{Könnten Sie mir bitte helfen ?} « Pourriez-vous m'aider, s'il vous plaît ? »

\begin{savaistu}[La politesse allemande]
\all{Ich hätte gern...} est la formule de politesse la plus utilisée dans les magasins et restaurants allemands : le subjonctif II sert aussi à être poli ! Dire \all{Ich will ein Brot} (« je veux un pain ») est correct grammaticalement mais très impoli ; un Allemand dira toujours \all{Ich hätte gern ein Brot}. À Yaoundé comme à Berlin, la politesse ouvre toutes les portes !
\end{savaistu}

\courssub{Formation du subjonctif II}

\begin{propriete}[Règle de formation]
Le subjonctif II se forme sur le \textbf{radical du Präteritum} auquel on ajoute les terminaisons \textbf{-e, -est, -e, -en, -et, -en}.

\begin{itemize}
\item \textbf{Verbes forts} : la voyelle du Präteritum \emph{a, o, u} prend un \textbf{tréma} : \all{haben : hatte → hätte} ; \all{sein : war → wäre} ; \all{kommen : kam → käme} ; \all{geben : gab → gäbe}.
\item \textbf{Verbes faibles} : la forme est identique au Präteritum (\all{machte, verschwendete, kaufte}) — c'est pourquoi on leur préfère la forme avec \all{würde} (voir plus bas).
\item \textbf{Verbes modaux} : tréma pour \all{können → könnte}, \all{müssen → müsste}, \all{dürfen → dürfte} ; pas de tréma pour \all{sollen → sollte} et \all{wollen → wollte}.
\end{itemize}
\end{propriete}

\begin{popfigure}
\begin{center}
\begin{tikzpicture}[node distance=0.5cm and 0.7cm]
\node[draw=popblue, fill=popblueL, rounded corners, align=center, font=\bfseries] (praet) {Präteritum\\ \all{kam}};
\node[draw=poporange, fill=poporangeL, rounded corners, align=center, right=of praet] (umlaut) {+ tréma\\ \all{käm-}};
\node[draw=popgreen, fill=popgreenL, rounded corners, align=center, right=of umlaut] (end) {+ terminaison\\ \all{käm-e}};
\node[draw=poppurple, fill=poppurpleL, rounded corners, align=center, right=of end, font=\bfseries] (res) {Subjonctif II\\ \all{ich käme}};
\draw[-{Stealth}, very thick, popdark] (praet) -- (umlaut);
\draw[-{Stealth}, very thick, popdark] (umlaut) -- (end);
\draw[-{Stealth}, very thick, popdark] (end) -- (res);
\end{tikzpicture}
\legende{Schéma de formation : du Präteritum au subjonctif II (verbe fort \all{kommen})}
\end{center}
\end{popfigure}

\begin{exemplebox}[Autres transformations]
\begin{itemize}
\item \all{geben : gab → gäbe} : \all{Es gäbe keine Dürre mehr.} « Il n'y aurait plus de sécheresse. »
\item \all{sein : war → wäre} : \all{Die Luft wäre sauberer.} « L'air serait plus propre. »
\item \all{haben : hatte → hätte} : \all{Wir hätten mehr Zeit.} « Nous aurions plus de temps. »
\item \all{fahren : fuhr → führe} : la forme simple existe (\all{wenn du führest}), mais la langue courante préfère \all{würde fahren} : \all{Wenn du das Fahrrad nehmen würdest, wäre das besser.} « Si tu prenais le vélo, ce serait mieux. »
\end{itemize}
\end{exemplebox}

\courssub{Les trois formes à connaître par cœur et les modaux}

Trois formes reviennent dans presque toutes les phrases au subjonctif II : \all{wäre} (\all{sein}), \all{hätte} (\all{haben}) et \all{gäbe} (\all{geben}, dans la tournure \all{es gäbe} = « il y aurait »). Apprenez-les par cœur, ainsi que les modaux.

\begin{popfigure}
\begin{center}
\begin{tikzpicture}[node distance=0.4cm]
\node[draw=popblue, fill=popblueL, rounded corners, minimum width=3.6cm, align=center] (t1) {\textbf{sein → wäre}\\ ich wäre\\ du wär(e)st\\ er/sie/es wäre\\ wir wären\\ ihr wär(e)t\\ sie/Sie wären};
\node[draw=poporange, fill=poporangeL, rounded corners, minimum width=3.6cm, align=center, right=0.6cm of t1] (t2) {\textbf{haben → hätte}\\ ich hätte\\ du hättest\\ er/sie/es hätte\\ wir hätten\\ ihr hättet\\ sie/Sie hätten};
\node[draw=popgreen, fill=popgreenL, rounded corners, minimum width=3.6cm, align=center, right=0.6cm of t2] (t3) {\textbf{können → könnte}\\ ich könnte\\ du könntest\\ er/sie/es könnte\\ wir könnten\\ ihr könntet\\ sie/Sie könnten};
\end{tikzpicture}
\legende{Conjugaison complète de \all{wäre}, \all{hätte} et \all{könnte} aux six personnes}
\end{center}
\end{popfigure}

\begin{center}
\begin{tabular}{|l|l|l|}
\hline
\rowcolor{popblueL} \textbf{Infinitif} & \textbf{Präteritum} & \textbf{Subjonctif II (ich)} \\
\hline
sein & war & wäre \\
\hline
haben & hatte & hätte \\
\hline
geben & gab & gäbe \\
\hline
kommen & kam & käme \\
\hline
können & konnte & könnte \\
\hline
müssen & musste & müsste \\
\hline
sollen & sollte & sollte \\
\hline
dürfen & durfte & dürfte \\
\hline
wollen & wollte & wollte \\
\hline
\end{tabular}
\end{center}

\all{Du müsstest weniger Plastik kaufen.} « Tu devrais acheter moins de plastique. » (au sens de : il faudrait que tu achètes moins de plastique — obligation non respectée)

\all{Das dürfte nicht so teuer sein.} « Cela ne devrait pas être si cher. »

\courssub{Le conditionnel avec \all{würde} + infinitif}

Pour les verbes faibles et la grande majorité des verbes forts, la langue courante préfère la forme composée : \textbf{\all{würde} + infinitif} (à la fin de la phrase). \all{würde} est le subjonctif II de \all{werden} (\all{wurde → würde}).

\begin{center}
\begin{tabular}{|l|l|}
\hline
\rowcolor{poporangeL} \textbf{Personne} & \textbf{würde} \\
\hline
ich & würde \\
\hline
du & würdest \\
\hline
er/sie/es & würde \\
\hline
wir & würden \\
\hline
ihr & würdet \\
\hline
sie/Sie & würden \\
\hline
\end{tabular}
\end{center}

\all{Ich würde mehr recyceln.} « Je recyclerais davantage. »

\all{Wenn ich Zeit hätte, würde ich mehr für die Umwelt tun.} « Si j'avais le temps, je ferais plus pour l'environnement. »

\all{Ohne Autos würden wir besser atmen.} « Sans voitures, nous respirerions mieux. »

\begin{methode}[Subjonctif II simple ou \all{würde} + infinitif ?]
\begin{enumerate}
\item Le verbe est-il \all{sein}, \all{haben}, \all{geben} (\all{es gibt}) ou un \textbf{modal} ? → utilise la \textbf{forme simple} : \all{wäre, hätte, gäbe, könnte, sollte...}
\item Sinon (tous les autres verbes, surtout les verbes faibles) → utilise \textbf{\all{würde} + infinitif} : \all{ich würde kaufen, du würdest fahren...}
\item Vérifie la place des mots : \all{würde} en position 2 dans la principale (l'infinitif à la fin) ; dans une subordonnée introduite par \all{wenn}, le verbe conjugué va à la fin : \all{Wenn ich Zeit hätte, ...}
\end{enumerate}
\end{methode}

\courssub{Le subjonctif II du passé (complément Bac)}

Pour exprimer un irréel \textbf{dans le passé} (« si nous avions agi plus tôt... », ce qui n'a pas eu lieu), on utilise \all{hätte} ou \all{wäre} + participe II, selon l'auxiliaire du verbe :

\all{Wenn wir früher gehandelt hätten, wäre die Lage besser.} « Si nous avions agi plus tôt, la situation serait meilleure. »

\all{Wenn die Regenzeit früher gekommen wäre, hätten die Bauern mehr geerntet.} « Si la saison des pluies était arrivée plus tôt, les paysans auraient plus récolté. »

Règle : verbes conjugués avec \all{haben} → \all{hätte} + participe II ; verbes de mouvement ou de changement conjugués avec \all{sein} → \all{wäre} + participe II. Le français emploie ici le plus-que-parfait après « si » et le conditionnel passé dans la principale.

\courssub{Comparaison français / allemand}

\begin{center}
\begin{tabularx}{\linewidth}{|X|X|X|}
\hline
\rowcolor{popblueL} \textbf{Français} & \textbf{Deutsch} & \textbf{Remarque} \\
\hline
Si j'avais le temps... & Wenn ich Zeit hätte... & « si + imparfait » → \all{wenn} + subj. II \\
\hline
je ferais plus & würde ich mehr tun & conditionnel → \all{würde} + infinitif \\
\hline
Si seulement il pleuvait ! & Wenn es doch regnen würde ! & \all{doch} / \all{nur} renforcent le souhait \\
\hline
Tu devrais trier. & Du solltest den Müll trennen. & conseil = \all{sollte} \\
\hline
J'aimerais un café. & Ich hätte gern einen Kaffee. & politesse = \all{hätte gern} \\
\hline
\end{tabularx}
\end{center}

\begin{attention}[Erreurs fréquentes des francophones]
\begin{itemize}
\item \textbf{Indicatif après \all{wenn}} pour une hypothèse irréelle : ✗ \all{Wenn ich Zeit habe, mache ich mehr.} (cela signifie « quand j'ai le temps », réel !) → ✓ \all{Wenn ich Zeit hätte, würde ich mehr tun.}
\item \textbf{Mélange des formes} : ✗ \all{ich wäre habe} ou ✗ \all{ich hätte bin} — impensable ! Une seule forme : \all{ich hätte} ou \all{ich wäre}.
\item \textbf{Oubli du tréma} : ✗ \all{ich hatte gern} (indicatif : « j'avais envie ») au lieu de ✓ \all{ich hätte gern} (« j'aimerais »).
\item \textbf{Place de l'infinitif} : ✗ \all{Ich würde recyceln mehr.} → ✓ \all{Ich würde mehr recyceln.} (l'infinitif se place en toute fin de phrase).
\item \textbf{Double marque de l'irréel} : ✗ \all{ich würde kämen} — on ne combine jamais \all{würde} avec une forme de subjonctif II : ✓ \all{ich käme} ou \all{ich würde kommen}.
\end{itemize}
\end{attention}

\begin{textebox}[Mini-texte : „Wenn ich Umweltminister wäre...“]
\all{Wenn ich Umweltminister wäre, würde ich zuerst den Müll in allen Städten trennen lassen. Ich hätte gern sauberere Straßen in Yaoundé und Douala. Die Bürger sollten weniger Plastik kaufen und mehr recyceln. Wenn jeder nur einen Baum pflanzen würde, gäbe es bald neue Wälder! Ich würde auch mehr Busse und Fahrradwege bauen, damit die Luft sauberer wäre. Wenn wir früher gehandelt hätten, wäre die Erderwärmung heute nicht so stark. Aber es ist nicht zu spät: Wenn wir jetzt zusammenarbeiteten, könnten wir die Welt retten. Wenn das doch alle verstehen würden!}

\textbf{Glossaire} : der Umweltminister, - (ministre de l'Environnement) ; trennen lassen (faire trier) ; der Bürger, - (le citoyen) ; der Baum, die Bäume (l'arbre) ; pflanzen (planter) ; der Fahrradweg, -e (la piste cyclable) ; damit (afin que) ; retten (sauver).

\textbf{Questions} : 1. Que ferait le narrateur en premier ? 2. Quel moyen de transport veut-il développer ? 3. Relevez trois formes de subjonctif II et identifiez leur infinitif.
\end{textebox}

\begin{exoresolu}[Mets à la forme correcte du subjonctif II]
Énoncé : Conjugue les verbes entre parenthèses au subjonctif II.\\
1. Wenn ich reich (sein), (kaufen) ich viele Bäume.\\
2. Du (sollen) das Licht ausschalten.\\
3. Wenn es (geben) weniger Autos, (sein) die Luft sauberer.\\
4. Ich (haben) gern eine saubere Stadt.
\tcblower
Solution détaillée :\\
1. \all{Wenn ich reich wäre, würde ich viele Bäume kaufen.} — \all{sein} → forme simple \all{wäre} ; \all{kaufen} est un verbe faible → \all{würde kaufen}, infinitif en fin de phrase.\\
2. \all{Du solltest das Licht ausschalten.} — modal \all{sollen} → forme simple \all{solltest} ; attention, \all{ausschalten} est un verbe à particule séparable : \all{aus-} va en fin de phrase.\\
3. \all{Wenn es weniger Autos gäbe, wäre die Luft sauberer.} — \all{geben} → \all{gäbe} (tréma sur \all{gab}) ; \all{sein} → \all{wäre}.\\
4. \all{Ich hätte gern eine saubere Stadt.} — formule de politesse : \all{hätte gern}, avec tréma obligatoire.
\end{exoresolu}

\begin{exercice}[À toi de jouer]
Traduis en allemand : 1. « Si j'avais plus de temps, je planteraient des arbres. » 2. « Tu devrais moins gaspiller l'eau. » 3. « Si seulement il ne faisait pas si chaud ! » 4. « J'aimerais un monde sans pollution. » 5. « Si nous avions agi plus tôt, la situation serait meilleure. »
\end{exercice}

\begin{corrige}[Exercice « À toi de jouer »]
1. \all{Wenn ich mehr Zeit hätte, würde ich Bäume pflanzen.} 2. \all{Du solltest weniger Wasser verschwenden.} 3. \all{Wenn es doch nicht so heiß wäre !} 4. \all{Ich hätte gern eine Welt ohne Umweltverschmutzung.} 5. \all{Wenn wir früher gehandelt hätten, wäre die Lage besser.}
\end{corrige}

\begin{aretenir}[L'essentiel du subjonctif II]
\begin{itemize}
\item Emplois : hypothèse irréelle, souhait (\all{doch}/\all{nur}), conseil (\all{sollte}), politesse (\all{hätte gern}, \all{könnten Sie...}).
\item Formation : radical du Präteritum + tréma (verbes forts) + terminaisons -e, -est, -e, -en, -et, -en.
\item À connaître par cœur : \all{wäre, hätte, gäbe, könnte, sollte, müsste, dürfte, wollte}.
\item Tous les autres verbes : \all{würde} + infinitif en fin de phrase.
\item Passé : \all{hätte}/\all{wäre} + participe II.
\end{itemize}
\end{aretenir}

\coursec{Grammaire 2 : subordonnées de conséquence et de temps}

Dans la section précédente, nous avons appris à exprimer l'irréel et l'hypothèse grâce au \cle{Konjunktiv II} : \all{Wenn wir weniger CO₂ produzierten, wäre die Luft sauberer.} (« Si nous produisions moins de CO₂, l'air serait plus propre. »). Vous avez remarqué que cette phrase contient une subordonnée introduite par \all{wenn}, avec le verbe conjugué \all{produzierten} rejeté à la fin. Nous allons maintenant approfondir cette mécanique fondamentale de la phrase allemande en étudiant deux grandes familles de subordonnées : les subordonnées de \cle{conséquence} (\all{so dass}...) et les subordonnées de \cle{temps} (\all{als}, \all{wenn}, \all{nachdem}, \all{bevor}...). Ce sont des points de grammaire centraux pour la version, le thème et l'expression écrite au baccalauréat.

\courssub{Observation : ce que le texte nous montre}

Reprenons deux phrases tirées de notre texte support « Die Erde wird wärmer » :

\begin{exemplebox}[Phrases relevées dans le texte]
\all{Die Erde wird wärmer, so dass die Gletscher schmelzen.}\\
« La Terre se réchauffe, si bien que les glaciers fondent. »

\all{Als der Regen ausblieb, vertrockneten die Felder.}\\
« Lorsque la pluie fit défaut, les champs se desséchèrent. »
\end{exemplebox}

Observons attentivement :

\begin{itemize}
\item Dans la première phrase, la conjonction \all{so dass} introduit la \emph{conséquence} du réchauffement ; le verbe conjugué \all{schmelzen} se trouve \textbf{en toute dernière position} de la subordonnée.
\item Dans la seconde, la conjonction \all{als} situe un événement \emph{unique au passé} ; le verbe \all{ausblieb} est rejeté à la fin de la subordonnée, et comme celle-ci ouvre la phrase, le verbe de la principale (\all{vertrockneten}) vient \textbf{immédiatement après la virgule} : c'est l'inversion.
\end{itemize}

Toute cette section repose sur ces deux constats. Formalisons-les.

\courssub{Les subordonnées de conséquence : so dass / sodass}

\begin{propriete}[La subordonnée de conséquence avec « so dass »]
\all{so dass} (orthographe réformée : \all{sodass}) introduit une subordonnée exprimant la \textbf{conséquence} de ce qui est dit dans la principale. Le verbe conjugué se place \textbf{en dernière position} de la subordonnée.

\all{Es regnete so stark, dass die Straßen überflutet wurden.}\\
« Il a plu si fort que les rues ont été inondées. »

Structure : [Principale] , \all{so dass} + sujet + ... + \textbf{verbe conjugué en fin}.
\end{propriete}

Très souvent, la principale contient un adverbe ou un adjectif intensifié par \all{so} (« si, tellement »), ce qui correspond au français « si... que », « tellement... que » :

\begin{exemplebox}[La construction « so + adjectif/adverbe, dass... »]
\all{Es war so heiß, dass die Felder vertrockneten.}\\
« Il faisait si chaud que les champs se desséchèrent. »

\all{Die Dürre dauerte so lange, dass viele Bauern ihre Ernte verloren.}\\
« La sécheresse dura si longtemps que beaucoup de paysans perdirent leur récolte. »

\all{Der Smog war so dicht, dass man die Berge nicht mehr sehen konnte.}\\
« Le smog était si épais qu'on ne pouvait plus voir les montagnes. »
\end{exemplebox}

\begin{attention}[« so dass » n'est pas « also »]
\begin{itemize}
\item \all{so dass} / \all{sodass} : les deux graphies sont correctes depuis la réforme orthographique. C'est une \textbf{conjonction de subordination} → verbe à la fin.
\item \all{also} (« donc ») est un \textbf{connecteur} à l'intérieur d'une proposition principale → le verbe reste en 2\textsuperscript{e} position : \all{Es regnete stark, also blieben wir zu Hause.} (« Il pleuvait fort, donc nous sommes restés à la maison. »).
\item Ne traduisez jamais « si... que » par \all{wenn... dass} : le français « si » d'intensité se rend par \all{so}, jamais par \all{wenn}.
\end{itemize}
\end{attention}

\begin{propriete}[Conséquence en proposition principale : deshalb / darum]
Pour exprimer la conséquence sans subordonnée, l'allemand utilise les connecteurs \all{deshalb} et \all{darum} (« c'est pourquoi »). Ce sont des \textbf{adverbes} : ils occupent la 1\textsuperscript{re} position et provoquent l'\textbf{inversion} (verbe en 2\textsuperscript{e} position, sujet après le verbe).

\all{Der Regen blieb aus. Deshalb vertrockneten die Felder.}\\
« La pluie fit défaut. C'est pourquoi les champs se desséchèrent. »

\all{Die Luft ist verschmutzt, darum tragen viele Menschen Masken.}\\
« L'air est pollué, c'est pourquoi beaucoup de gens portent des masques. »
\end{propriete}

\begin{center}
\begin{tabularx}{\linewidth}{|X|X|X|}
\hline
\rowcolor{popblueL} Français & Deutsch & Remarque \\
\hline
si chaud que... & so heiß, dass... & verbe à la fin \\
\hline
si bien que & so dass / sodass & conjonction de subordination \\
\hline
c'est pourquoi, donc & deshalb, darum, also & connecteur : inversion, verbe en 2\textsuperscript{e} position \\
\hline
Il a plu, donc les rues ont été inondées. & Es hat geregnet, deshalb wurden die Straßen überflutet. & deux propositions principales \\
\hline
Il a si plu que les rues ont été inondées. & Es hat so stark geregnet, dass die Straßen überflutet wurden. & principale + subordonnée \\
\hline
\end{tabularx}
\end{center}

\courssub{Les subordonnées de temps}

\begin{propriete}[« als » ou « wenn » ? La distinction capitale]
\begin{itemize}
\item \all{als} s'emploie pour un événement \textbf{unique et accompli au passé} (Präteritum ou Plusquamperfekt) : « quand, lorsque ».
\item \all{wenn} s'emploie pour le \textbf{présent}, le \textbf{futur} et toute \textbf{répétition} (même au passé : « chaque fois que »).
\end{itemize}

\all{Als ich ein Kind war, war der Wald noch größer.}\\
« Quand j'étais enfant, la forêt était encore plus grande. » (période unique du passé)

\all{Wenn es heiß ist, trinken wir mehr Wasser.}\\
« Quand il fait chaud, nous buvons plus d'eau. » (présent, vérité générale)

\all{Wenn der Sommer kam, gingen wir früher immer an den Fluss.}\\
« Quand l'été arrivait, nous allions toujours à la rivière. » (répétition au passé)
\end{propriete}

\begin{attention}[Le piège n° 1 des francophones]
En français, « quand » sert à tout. En allemand, il faut choisir :
\begin{itemize}
\item « Quand j'ai visité Yaoundé en 2019... » → \all{Als ich 2019 Yaoundé besuchte, ...} (événement unique au passé).
\item « Quand il pleut, le marché est vide. » → \all{Wenn es regnet, ist der Markt leer.} (présent).
\item « Quand tu viendras, nous planterons des arbres. » → \all{Wenn du kommst, werden wir Bäume pflanzen.} (futur : l'allemand met le présent après \all{wenn}).
\end{itemize}
\textbf{Jamais} \all{als} pour le présent ou le futur ; \textbf{jamais} \all{wenn} pour un fait unique du passé.
\end{attention}

\begin{propriete}[« nachdem » et la concordance des temps]
\all{nachdem} (« après que ») impose une \textbf{concordance des temps} stricte :
\begin{itemize}
\item \all{nachdem} + \textbf{Plusquamperfekt} → \textbf{Präteritum} dans la principale (antériorité au passé) ;
\item \all{nachdem} + \textbf{Perfekt} → \textbf{Präsens} dans la principale (antériorité au présent).
\end{itemize}

\all{Nachdem der Regen ausgeblieben war, vertrockneten die Felder.}\\
« Après que la pluie eut fait défaut, les champs se desséchèrent. »

\all{Nachdem die Fabrik geschlossen hatte, wurde die Luft sauberer.}\\
« Après que l'usine eut fermé, l'air devint plus propre. »

\all{Nachdem wir den Müll gesammelt haben, pflanzen wir Bäume.}\\
« Après avoir ramassé les déchets, nous plantons des arbres. »
\end{propriete}

\begin{propriete}[Les autres conjonctions temporelles]
\begin{itemize}
\item \all{bevor} (« avant que / avant de ») : \all{Bevor man neue Fabriken baut, muss man die Umwelt schützen.} « Avant de construire de nouvelles usines, il faut protéger l'environnement. »
\item \all{während} (« pendant que ») : \all{Während die Politiker diskutierten, stieg die Temperatur weiter.} « Pendant que les politiciens discutaient, la température continuait de monter. »
\item \all{seitdem} (« depuis que ») : \all{Seitdem der Fluss verschmutzt ist, gibt es keine Fische mehr.} « Depuis que la rivière est polluée, il n'y a plus de poissons. »
\item \all{bis} (« jusqu'à ce que ») : \all{Wir warteten, bis der Bus kam.} « Nous attendîmes jusqu'à ce que le bus arrive. »
\item \all{sobald} (« dès que ») : \all{Sobald die Schule endet, beginnt unser Umweltprojekt.} « Dès que l'école finit, notre projet environnemental commence. »
\end{itemize}
\end{propriete}

\begin{attention}[« bevor » ou « nachdem » ?]
Ne confondez pas les deux : \all{bevor} = \textbf{avant} (l'action de la subordonnée est \emph{postérieure} à celle de la principale), \all{nachdem} = \textbf{après} (l'action de la subordonnée est \emph{antérieure}). Erreur typique : *\all{Nachdem man Fabriken baut, muss man die Umwelt schützen} pour dire « avant de construire... ». Et après \all{nachdem}, respectez toujours la concordance : Plusquamperfekt dans la subordonnée, Präteritum dans la principale.
\end{attention}

\begin{center}
\begin{tabularx}{\linewidth}{|l|l|X|}
\hline
\rowcolor{popblueL} Conjonction & Français & Exemple traduit \\
\hline
als & quand (passé, unique) & \all{Als der Sturm kam, fielen viele Bäume.} « Quand la tempête arriva, de nombreux arbres tombèrent. » \\
\hline
wenn & quand (présent, futur, répétition) & \all{Wenn die Sonne scheint, trocknet die Erde.} « Quand le soleil brille, la terre sèche. » \\
\hline
nachdem & après que & \all{Nachdem es geregnet hatte, wurde es kühler.} « Après qu'il eut plu, il fit plus frais. » \\
\hline
bevor & avant que / avant de & \all{Bevor wir weggehen, löschen wir das Feuer.} « Avant de partir, nous éteignons le feu. » \\
\hline
während & pendant que & \all{Während wir schliefen, regnete es.} « Pendant que nous dormions, il a plu. » \\
\hline
seitdem & depuis que & \all{Seitdem der Wald abgeholzt wurde, gibt es Dürren.} « Depuis que la forêt a été abattue, il y a des sécheresses. » \\
\hline
bis & jusqu'à ce que & \all{Er wartete, bis der Bus kam.} « Il attendit jusqu'à ce que le bus arrive. » \\
\hline
sobald & dès que & \all{Sobald die Gletscher schmelzen, steigt das Meer.} « Dès que les glaciers fondent, la mer monte. » \\
\hline
so dass & si bien que & \all{Es war so trocken, dass der Brunnen leer war.} « Il faisait si sec que le puits était vide. » \\
\hline
\end{tabularx}
\end{center}

\begin{popfigure}
\begin{tikzpicture}[scale=1]
\draw[->, thick, popdark] (0,0) -- (10.5,0) node[right] {\small Zeit};
\draw[thick, popdark] (1,-0.15) -- (1,0.15) node[above=0.02cm] at (1,0.15) {\small Vergangenheit};
\draw[thick, popdark] (5.5,-0.15) -- (5.5,0.15) node[above=0.02cm] at (5.5,0.15) {\small Gegenwart};
\draw[thick, popdark] (9.5,-0.15) -- (9.5,0.15) node[above=0.02cm] at (9.5,0.15) {\small Zukunft};
\fill[poporange] (2.5,0) circle (0.14cm);
\node[align=center, poporange, font=\small\bfseries] at (2.5,-0.9) {als\\ einmaliges Ereignis};
\foreach \x in {1.6,2.6,3.6} \fill[popblue] (\x,1.1) circle (0.11cm);
\node[align=center, popblue, font=\small\bfseries] at (2.6,1.9) {wenn\\ Wiederholung};
\fill[popgreen] (6.5,0) circle (0.14cm);
\fill[popgreen] (8.5,0) circle (0.14cm);
\node[align=center, popgreen, font=\small\bfseries] at (7.5,-0.9) {wenn\\ Präsens / Futur};
\end{tikzpicture}
\legende{Frise des temps : \all{als} marque un point unique au passé ; \all{wenn} couvre la répétition au passé, le présent et le futur.}
\end{popfigure}

\courssub{La règle d'or : la place du verbe}

\begin{propriete}[Verbe en fin de subordonnée et inversion]
\begin{enumerate}
\item Dans \textbf{toute} subordonnée introduite par une conjonction de subordination (\all{als}, \all{wenn}, \all{dass}, \all{weil}, \all{so dass}, \all{bevor}, \all{nachdem}, \all{während}, \all{bis}, \all{sobald}...), le \textbf{verbe conjugué va en dernière position} (verbe plein au Präteritum/Präsens ; auxiliaire \all{haben}/\all{sein} en toute fin aux temps composés).
\item Si la subordonnée \textbf{précède} la principale, elle occupe la 1\textsuperscript{re} position : le verbe de la principale vient \textbf{immédiatement après la virgule} (inversion sujet-verbe).
\end{enumerate}

\all{Als der Regen ausblieb, vertrockneten die Felder.}\\
(subordonnée : \all{ausblieb} en fin ; principale : \all{vertrockneten} juste après la virgule, puis le sujet \all{die Felder}.)

\all{Die Felder vertrockneten, als der Regen ausblieb.}\\
(ordre inverse possible : principale d'abord, sans inversion.)
\end{propriete}

\begin{popfigure}
\begin{tikzpicture}[node distance=0.35cm]
\node[draw=popblue, fill=popblueL, rounded corners, align=center, font=\small] (sub) {Subordonnée\\ \all{Als der Regen \textbf{ausblieb}}\\ verbe en fin};
\node[draw=poporange, fill=poporangeL, rounded corners, align=center, font=\small, right=0.8cm of sub] (haupt) {Principale\\ \all{\textbf{vertrockneten} die Felder}\\ verbe en 1\textsuperscript{re} position};
\draw[->, very thick, popdark] (sub) -- (haupt) node[midway, above, font=\small] {,};
\node[below=0.25cm of sub, font=\small\itshape, align=center] {1\textsuperscript{re} position};
\node[below=0.25cm of haupt, font=\small\itshape, align=center] {inversion : verbe puis sujet};
\end{tikzpicture}
\legende{Schéma de structure : [Subordonnée + verbe en fin] , [verbe + sujet ...].}
\end{popfigure}

\begin{methode}[Placer le verbe correctement en quatre étapes]
\begin{enumerate}
\item \textbf{Repérer} la conjonction qui ouvre la proposition.
\item \textbf{Classer} : est-ce une conjonction de subordination (\all{als}, \all{wenn}, \all{dass}, \all{weil}, \all{so dass}, \all{bevor}, \all{nachdem}, \all{während}, \all{seitdem}, \all{bis}, \all{sobald}) ou un connecteur (\all{deshalb}, \all{darum}, \all{also}, \all{aber}, \all{denn}) ?
\item Si c'est une \textbf{conjonction de subordination} → envoyer le \textbf{verbe conjugué à la fin} de la subordonnée (aux temps composés : l'auxiliaire \all{haben}/\all{sein} en toute fin, ex. \all{..., dass die Straßen überflutet \textbf{worden sind}}).
\item Si la subordonnée \textbf{ouvre la phrase} → placer le verbe de la principale \textbf{juste après la virgule}, puis le sujet.
\end{enumerate}
\end{methode}

\begin{textebox}[Ein heißer Sommer in Douala]
Als der Sommer kam, wurde es in Douala so heiß, dass viele Menschen kaum schlafen konnten. Während die Sonne brannte, vertrockneten die Felder rund um die Stadt, und seitdem der Regen ausgeblieben war, stieg der Preis für Gemüse auf den Märkten. Die Bauern warteten bis in den Abend, bevor sie ihre Pflanzen gossen, weil das Wasser am Tag zu schnell verdunstete. Nachdem der erste Regen endlich gefallen war, atmete die ganze Stadt auf. Sobald die Erde wieder weich war, begannen die Familien, Mais und Maniok zu pflanzen. Ein alter Bauer sagte: „Wenn ich jung war, regnete es jedes Jahr zur gleichen Zeit. Heute ist das Klima so verändert, dass wir nichts mehr vorhersagen können.“ Deshalb beschloss die Schule des Viertels, einen Umweltclub zu gründen. Die Schüler sammelten Plastikmüll, und nachdem sie den Schulhof sauber gemacht hatten, pflanzten sie zwanzig junge Bäume. „Wenn jeder nur einen Baum pflanzt“, erklärte die Lehrerin, „wird unsere Stadt in zehn Jahren grüner und kühler sein.“ Bevor die Schüler nach Hause gingen, versprachen sie, auch ihre Eltern zu überzeugen. So entstand aus einem heißen Sommer eine kleine Hoffnung für die Zukunft.

\textbf{Glossaire} : der Regen (--) la pluie ; das Feld, die Felder le champ ; der Bauer, die Bauern le paysan ; gießen (goss, gegossen) arroser ; verdunsten s'évaporer ; aufatmen (atmete auf, aufgeatmet) pousser un soupir de soulagement ; weich mou, tendre ; der Mais le maïs ; der Maniok le manioc ; vorhersagen prédire ; der Plastikmüll les déchets plastiques ; der Schulhof, die Schulhöfe la cour de l'école ; überzeugen convaincre ; die Hoffnung, die Hoffnungen l'espoir.
\end{textebox}

Questions de compréhension : 1) Pourquoi les paysans arrosaient-ils leurs plantes le soir ? 2) Quelle décision l'école a-t-elle prise, et pourquoi ? 3) Relevez dans le texte deux subordonnées avec \all{nachdem} et indiquez les temps employés. 4) Relevez une subordonnée de conséquence et traduisez-la.

\begin{exoresolu}[Exercice résolu : traduction et choix de la conjonction]
\textbf{Énoncé.} Traduisez en allemand : 1) « Quand j'étais petit, la rivière était encore propre. » 2) « Quand il fait chaud, nous buvons beaucoup d'eau. » 3) « Après que la pluie eut cessé, le soleil revint. » 4) « Il faisait si chaud que les champs se desséchèrent. »
\tcblower
\textbf{Solution détaillée.}

1) Événement unique au passé → \all{als} ; verbe à la fin ; inversion dans la principale :\\
\all{Als ich klein war, war der Fluss noch sauber.}

2) Présent, vérité générale → \all{wenn} :\\
\all{Wenn es heiß ist, trinken wir viel Wasser.}

3) Antériorité au passé → \all{nachdem} + Plusquamperfekt, Präteritum dans la principale ; auxiliaire \all{hatte} en toute fin :\\
\all{Nachdem der Regen aufgehört hatte, kam die Sonne zurück.}

4) « si... que » → \all{so ... dass} ; verbe à la fin :\\
\all{Es war so heiß, dass die Felder vertrockneten.}
\end{exoresolu}

\begin{exercice}[À vous de jouer]
1) Complétez avec \all{als} ou \all{wenn} : a) \all{... ich in Yaoundé ankam, regnete es.} b) \all{... es regnet, bleiben die Bendskins stehen.} c) \all{... wir früher den Wald besuchten, sangen immer viele Vögel.}
2) Remettez les mots dans l'ordre : \all{die Gletscher / so dass / wärmer / wird / schmelzen / die Erde}.
3) Traduisez : « Après que l'usine eut fermé, l'air devint plus propre. » puis « Avant de partir, nous éteignons le feu. »
4) Transformez avec \all{deshalb} : \all{Es regnete so stark, dass wir zu Hause blieben.}
\end{exercice}

\begin{corrige}[Exercice « À vous de jouer »]
1) a) \all{Als ich in Yaoundé ankam, regnete es.} (événement unique au passé) ; b) \all{Wenn es regnet, bleiben die Bendskins stehen.} (présent, généralité) ; c) \all{Wenn wir früher den Wald besuchten, sangen immer viele Vögel.} (répétition au passé, signalée par \all{immer}).
2) \all{Die Erde wird wärmer, so dass die Gletscher schmelzen.}
3) \all{Nachdem die Fabrik geschlossen hatte, wurde die Luft sauberer.} ; \all{Bevor wir weggehen, löschen wir das Feuer.}
4) \all{Es regnete stark, deshalb blieben wir zu Hause.} (connecteur : verbe \all{blieben} en 2\textsuperscript{e} position, sujet après le verbe).
\end{corrige}

\begin{savaistu}[Un classique du baccalauréat]
La distinction \all{als}/\all{wenn} et la place du verbe dans les subordonnées font partie des points les plus fréquemment évalués au Bac, aussi bien en \textbf{version} (repérer la conjonction pour comprendre la chronologie du texte) qu'en \textbf{thème} et en \textbf{expression écrite} (une faute de choix entre \all{als} et \all{wenn}, ou un verbe mal placé, y est systématiquement sanctionnée). Petit truc de correcteur : si vous pouvez remplacer « quand » par « chaque fois que » en français, c'est \all{wenn} ; si l'événement n'est arrivé qu'une seule fois dans le passé, c'est \all{als}.
\end{savaistu}

\begin{aretenir}[L'essentiel de la section]
\begin{itemize}
\item \all{so dass} / \all{sodass} introduit une conséquence, verbe en fin ; souvent avec \all{so + adjectif} dans la principale (« si... que »).
\item \all{deshalb}, \all{darum}, \all{also} sont des connecteurs : proposition principale, inversion.
\item \all{als} = passé unique ; \all{wenn} = présent, futur, répétition.
\item \all{nachdem} + Plusquamperfekt → Präteritum dans la principale.
\item Dans toute subordonnée, le verbe conjugué va à la fin ; si la subordonnée précède, le verbe de la principale suit immédiatement la virgule.
\end{itemize}
\end{aretenir}

\coursec{Méthode du commentaire et production guidée}

Dans les sections précédentes, tu as appris à construire le \cle{subjonctif II} pour exprimer l'irréel (\all{Wenn jeder den Müll trennen würde, wäre die Stadt sauberer.}) et à manier les subordonnées de conséquence (\all{so dass}, \all{so... dass}) et de temps (\all{wenn}, \all{als}, \all{bevor}, \all{nachdem}). Ces outils grammaticaux ne sont pas des fins en soi : au baccalauréat, ils servent à \textbf{commenter un texte} et à \textbf{produire un paragraphe d'opinion}. Cette dernière section te montre comment les mobiliser dans un écrit structuré et dans un débat oral.

\courssub{La méthode du commentaire de texte au Bac}

\begin{methode}[Structurer un commentaire de texte en trois étapes]
Un commentaire de texte au baccalauréat suit toujours le même plan en trois parties :
\begin{enumerate}
\item \textbf{Einleitung (introduction)} : présente le texte en une ou deux phrases — son type (article, reportage, interview), son thème et sa source. Formules-clés : \all{Der Text „...“ handelt von...} « Le texte traite de... » ; \all{Der Text ist ein Artikel aus einer Zeitung.} « Le texte est un article de journal. »
\item \textbf{Hauptteil (développement)} : analyse les idées principales, le vocabulaire du thème et les procédés de l'auteur. Formules-clés : \all{Der Autor erklärt, dass...} « L'auteur explique que... » ; \all{Er nennt die Folgen...} « Il cite les conséquences... » ; \all{Zuerst..., dann..., außerdem...} « D'abord..., ensuite..., de plus... ».
\item \textbf{Schluss (conclusion)} : fais un bilan bref puis donne ton \textbf{opinion personnelle justifiée}. Formule-clé : \all{Meiner Meinung nach ist es wichtig, dass...} « À mon avis, il est important que... », toujours suivie d'une justification avec \all{weil} ou \all{denn}.
\end{enumerate}
\end{methode}

\begin{popfigure}
\begin{center}
\begin{tikzpicture}[
node distance=0.9cm,
box/.style={rectangle, rounded corners=4pt, draw=popblue, fill=popblueL, text width=5.6cm, align=center, font=\small, inner sep=6pt},
lbl/.style={font=\small\itshape, text=popdark},
arr/.style={-{Stealth[length=3mm]}, thick, poporange}]
\node[box, align=center] (ein){\textbf{Einleitung}\\ \all{Der Text handelt von...}\\ (type, thème, source)};
\node[box, below=of ein, align=center] (haupt){\textbf{Hauptteil — Analyse}\\ \all{Der Autor erklärt, dass...}\\ \all{Er nennt die Folgen...}\\ (idées, vocabulaire, procédés)};
\node[box, below=of haupt, draw=poppurple, fill=poppurpleL, align=center] (schluss){\textbf{Schluss — Meinung}\\ \all{Meiner Meinung nach...}\\ \all{Zusammenfassend...}\\ (bilan + opinion justifiée)};
\draw[arr] (ein) -- node[lbl, right]{présenter} (haupt);
\draw[arr] (haupt) -- node[lbl, right]{juger} (schluss);
\end{tikzpicture}
\end{center}
\legende{Organigramme du commentaire de texte : Einleitung → Hauptteil → Schluss}
\end{popfigure}

\courssub{Commentaire modèle sur « Die Erde wird wärmer »}

Observe d'abord ce commentaire rédigé selon la méthode. Les \textbf{connecteurs} sont mis en italique : repère-les et note leur fonction.

\begin{textebox}[Commentaire modèle — Die Erde wird wärmer]
Der Text „Die Erde wird wärmer“ handelt von der Umweltverschmutzung und ihren Folgen. \emph{Zuerst} erklärt der Autor die Ursachen : Fabriken und Autos produzieren viel CO$_2$. \emph{Dann} nennt er die Folgen : die Gletscher schmelzen und es gibt Dürren. \emph{Außerdem} zeigt er, dass der Klimawandel auch Afrika betrifft. \emph{Deshalb} fordert er, dass alle Länder zusammenarbeiten. \emph{Meiner Meinung nach} ist der Text sehr aktuell, denn auch in Kamerun sehen wir den Klimawandel. Wir sollten sofort handeln, bevor es zu spät ist. \emph{Zusammenfassend} kann man sagen : der Text informiert gut und überzeugt den Leser.
\end{textebox}

\textbf{Traduction :} « Le texte „La Terre se réchauffe“ traite de la pollution de l'environnement et de ses conséquences. D'abord, l'auteur explique les causes : les usines et les voitures produisent beaucoup de CO$_2$. Ensuite, il cite les conséquences : les glaciers fondent et il y a des sécheresses. De plus, il montre que le changement climatique touche aussi l'Afrique. C'est pourquoi il demande que tous les pays coopèrent. À mon avis, le texte est très actuel, car au Cameroun aussi nous voyons le changement climatique. Nous devrions agir immédiatement, avant qu'il ne soit trop tard. En résumé, on peut dire : le texte informe bien et convainc le lecteur. »

\textbf{Glossaire :} \all{die Ursache, -n} « la cause » ; \all{die Folge, -n} « la conséquence » ; \all{der Gletscher, -} « le glacier » ; \all{schmelzen (schmilzt, schmolz, ist geschmolzen)} « fondre » ; \all{betreffen (betrifft, betraf, hat betroffen)} « toucher, concerner » ; \all{fordern} « exiger, demander » ; \all{zusammenarbeiten} « coopérer » ; \all{handeln} « agir » ; \all{überzeugen} « convaincre ».

\courssub{Les connecteurs d'opinion et d'argumentation}

Pour exprimer ton point de vue et organiser ton argumentation, tu disposes d'une boîte à outils de connecteurs. Apprends-les par cœur : ce sont eux qui donnent sa structure à ton texte.

\begin{center}
\begin{tabularx}{\linewidth}{|X|X|X|}
\hline
\rowcolor{popblueL} \textbf{Deutsch} & \textbf{Français} & \textbf{Emploi / remarque} \\
\hline
\all{ich denke, dass...} & je pense que... & opinion + subordonnée (verbe à la fin) \\
\hline
\all{ich meine, dass...} & j'estime que... & synonyme de \all{ich denke} \\
\hline
\all{meiner Meinung nach} & à mon avis & postposition : \all{nach} suit le nom \\
\hline
\all{einerseits... andererseits...} & d'une part... d'autre part... & peser le pour et le contre \\
\hline
\all{zuerst / dann / außerdem} & d'abord / ensuite / de plus & ordonner le développement \\
\hline
\all{deshalb / deswegen} & c'est pourquoi & conséquence (verbe en 2\textsuperscript{e} position) \\
\hline
\all{denn / weil} & car / parce que & justifier une opinion \\
\hline
\all{zum Schluss / zusammenfassend} & pour conclure / en résumé & annoncer la conclusion \\
\hline
\end{tabularx}
\end{center}

\begin{exemplebox}[Exemples traduits]
\all{Meiner Meinung nach sollten wir mehr Bäume pflanzen.} « À mon avis, nous devrions planter plus d'arbres. »

\all{Wenn jeder den Müll trennen würde, wäre die Stadt sauberer.} « Si chacun triait les déchets, la ville serait plus propre. »

\all{Einerseits ist Plastik praktisch, andererseits verschmutzt es die Umwelt.} « D'une part le plastique est pratique, d'autre part il pollue l'environnement. »

\all{Ich denke, dass der Staat mehr für den Wald tun muss.} « Je pense que l'État doit faire plus pour la forêt. »

\all{Zum Schluss möchte ich sagen : Jeder kann helfen.} « Pour conclure, je voudrais dire : chacun peut aider. »
\end{exemplebox}

\begin{attention}[Pièges fréquents des francophones]
\begin{itemize}
\item Dans \all{meiner Meinung nach}, la préposition \all{nach} est une \textbf{postposition} : elle se place \emph{après} le nom. Ne dis pas ✗ \all{nach meiner Meinung} dans ce sens ; la forme correcte est ✓ \all{meiner Meinung nach}.
\item Évite les opinions « nues » sans justification. Une opinion au Bac doit toujours être suivie de \all{weil} ou \all{denn} : ✗ \all{Der Text ist gut.} → ✓ \all{Der Text ist gut, denn er gibt konkrete Beispiele.} « Le texte est bon, car il donne des exemples concrets. »
\item Après \all{deshalb} en tête de phrase, le verbe conjugué reste en deuxième position : ✓ \all{Deshalb sollten wir handeln.} (et non ✗ \all{Deshalb wir sollten handeln.}).
\item Attention à la différence entre \all{denn} (conjonction de coordination : le verbe reste en deuxième position) et \all{weil} (conjonction de subordination : le verbe conjugué va à la fin) : \all{..., denn er gibt Beispiele.} / \all{..., weil er Beispiele gibt.}
\item N'oublie pas la majuscule aux noms : \all{die Meinung}, \all{der Müll}, \all{die Umwelt}.
\end{itemize}
\end{attention}

\begin{exoresolu}[Transformer des notes en phrases de commentaire]
\textbf{Énoncé :} À partir de ces notes, rédige deux phrases de commentaire avec les connecteurs indiqués. Notes : \all{Autor / Ursachen nennen / CO$_2$ und Fabriken} ; opinion : \all{Text / aktuell / Klimawandel in Kamerun sichtbar}. Connecteurs : \all{zuerst}, \all{meiner Meinung nach... denn}.
\tcblower
\textbf{Solution détaillée :}
\begin{enumerate}
\item \all{Zuerst nennt der Autor die Ursachen : Fabriken und Autos produzieren viel CO$_2$.} « D'abord, l'auteur cite les causes : les usines et les voitures produisent beaucoup de CO$_2$. » — Le connecteur \all{zuerst} ouvre le développement ; le verbe \all{nennt} reste en deuxième position.
\item \all{Meiner Meinung nach ist der Text sehr aktuell, denn der Klimawandel ist auch in Kamerun sichtbar.} « À mon avis, le texte est très actuel, car le changement climatique est visible aussi au Cameroun. » — \all{meiner Meinung nach} en tête, verbe \all{ist} en deuxième position ; la justification est introduite par \all{denn} (conjonction de coordination, sans rejet du verbe à la fin).
\end{enumerate}
\end{exoresolu}

\courssub{Production guidée écrite : « Was können wir in Kamerun für die Umwelt tun ? »}

\begin{exercice}[Rédiger un paragraphe argumenté (6 à 8 phrases)]
Rédige un court paragraphe en allemand (6 à 8 phrases) sur le sujet : \all{Was können wir in Kamerun für die Umwelt tun ?} « Que pouvons-nous faire au Cameroun pour l'environnement ? »

\textbf{Contraintes obligatoires :}
\begin{itemize}
\item au moins \textbf{deux verbes au subjonctif II} (\all{wir sollten...}, \all{wenn wir... würden, wäre...}) ;
\item au moins \textbf{une subordonnée de conséquence} (\all{so dass...} ou \all{so..., dass...}) ;
\item au moins \textbf{deux subordonnées de temps} (\all{wenn}, \all{bevor}, \all{nachdem}) ;
\item au moins \textbf{trois mots du vocabulaire du thème} (\all{die Umweltverschmutzung}, \all{der Müll}, \all{der Wald}, \all{die Dürre}, \all{das Klima}...) ;
\item des connecteurs : \all{zuerst}, \all{außerdem}, \all{meiner Meinung nach}, \all{zum Schluss}.
\end{itemize}
\end{exercice}

\begin{corrige}[Production guidée — modèle de rédaction]
\all{In Kamerun können wir viel für die Umwelt tun. Zuerst sollten wir in den Schulen und Städten mehr Bäume pflanzen, denn der Wald schützt das Klima. Wenn jeder Bürger den Müll trennen würde, wären unsere Straßen sauberer. Außerdem müssen wir Wasser sparen, bevor die Dürre kommt. Die Regierung wirbt so viel für den Umweltschutz, dass immer mehr Jugendliche mitmachen. Nachdem wir den Plastikmüll reduziert haben, können wir ihn recyceln. Meiner Meinung nach sollten wir sofort handeln, denn die Erderwärmung betrifft auch uns. Zum Schluss : Wenn wir zusammenarbeiten, wird Kamerun grüner und gesünder.}

« Au Cameroun, nous pouvons beaucoup faire pour l'environnement. D'abord, nous devrions planter plus d'arbres dans les écoles et les villes, car la forêt protège le climat. Si chaque citoyen triait les déchets, nos rues seraient plus propres. De plus, nous devons économiser l'eau avant que la sécheresse n'arrive. Le gouvernement fait tellement de publicité pour la protection de l'environnement que de plus en plus de jeunes participent. Après avoir réduit les déchets plastiques, nous pouvons les recycler. À mon avis, nous devrions agir immédiatement, car le réchauffement climatique nous concerne aussi. Pour conclure : si nous coopérons, le Cameroun deviendra plus vert et plus sain. »

\textbf{Vérification des contraintes :} subjonctif II : \all{sollten}, \all{würde... trennen}, \all{wären} ; conséquence : \all{so viel..., dass...} ; temps : \all{bevor}, \all{nachdem}, \all{wenn} ; vocabulaire : \all{die Umwelt}, \all{der Wald}, \all{der Müll}, \all{die Dürre}, \all{die Erderwärmung}.

\textbf{Remarque grammaticale :} après \all{nachdem}, la subordonnée est au \emph{Perfekt} quand la principale est au présent (\all{Nachdem wir den Müll reduziert haben, können wir ihn recyceln.}) : l'action de la subordonnée est toujours \emph{antérieure} à celle de la principale.
\end{corrige}

\courssub{Production orale guidée : débat en binômes}

\begin{experience}[Débat : « Sollen wir Plastik verbieten ? »]
Par binômes, organisez un mini-débat sur la question : \all{Sollen wir Plastik verbieten ?} « Faut-il interdire le plastique ? » L'élève A est \textbf{pour} (\all{dafür}), l'élève B est \textbf{contre} (\all{dagegen}). Déroulement :
\begin{enumerate}
\item Chacun prépare deux arguments avec les structures de la leçon : un subjonctif II (\all{Wenn wir Plastik verbieten würden, hätten wir weniger Müll.}) et une subordonnée de temps ou de conséquence.
\item Utilisez \all{einerseits... andererseits} pour reconnaître l'argument adverse : \all{Einerseits ist Plastik billig und praktisch, andererseits verschmutzt es Flüsse und Meere.}
\item Concluez chacun avec \all{meiner Meinung nach... weil / denn...}.
\item La classe vote : quel binôme a le mieux argumenté ?
\end{enumerate}
\textbf{Exemples de répliques :} \all{Ich meine, dass wir Plastik verbieten sollten, denn es tötet die Tiere im Meer.} « J'estime que nous devrions interdire le plastique, car il tue les animaux de la mer. » — \all{Ich denke, dass ein Verbot schwierig ist, weil viele Händler auf den Märkten von Douala Plastiktüten benutzen.} « Je pense qu'une interdiction est difficile, parce que beaucoup de vendeurs des marchés de Douala utilisent des sachets plastiques. »
\end{experience}

\courssub{Grille d'auto-évaluation}

Après ta production, évalue-toi honnêtement avec cette grille. Coche chaque critère atteint.

\begin{center}
\begin{tabularx}{\linewidth}{|X|c|c|}
\hline
\rowcolor{popgreenL} \textbf{Critère} & \textbf{Oui} & \textbf{À revoir} \\
\hline
J'ai utilisé au moins 3 mots du vocabulaire du thème (\all{die Umweltverschmutzung}, \all{das Klima}, \all{der Wald}...) & $\square$ & $\square$ \\
\hline
J'ai employé au moins 2 subjonctifs II corrects (\all{sollten}, \all{würden}, \all{wäre}) & $\square$ & $\square$ \\
\hline
J'ai écrit 1 subordonnée de conséquence et 2 subordonnées de temps avec le verbe à la fin & $\square$ & $\square$ \\
\hline
Mon paragraphe a une structure claire : introduction, arguments, conclusion & $\square$ & $\square$ \\
\hline
J'ai justifié mon opinion avec \all{weil} ou \all{denn} & $\square$ & $\square$ \\
\hline
Orthographe : majuscules aux noms, Umlaute, ponctuation & $\square$ & $\square$ \\
\hline
\end{tabularx}
\end{center}

\begin{savaistu}[La qualité avant la quantité]
Au baccalauréat, un paragraphe d'opinion bien structuré — avec des connecteurs (\all{zuerst}, \all{außerdem}, \all{meiner Meinung nach}, \all{zum Schluss}), une grammaire correcte (verbe en deuxième position, verbe rejeté à la fin des subordonnées) et un vocabulaire précis — vaut plus de points qu'un long texte plein d'erreurs. Les correcteurs récompensent la maîtrise, pas la longueur. Privilégie donc toujours la qualité : six phrases irréprochables valent mieux que quinze phrases approximatives. Relis-toi en vérifiant trois choses : la place du verbe, la majuscule des noms et la justification de ton opinion.
\end{savaistu}

\newpage

\coursec{Activité d'intégration}

\courssub{Objectifs de l'activité}

À l'issue de cette activité d'intégration, l'élève sera capable de :

\begin{itemize}
\item réutiliser dans une production authentique le \cle{vocabulaire thématique} de l'environnement et du climat (\all{die Umweltverschmutzung}, \all{die Erderwärmung}, \all{das Klima}, \all{das CO$_2$}, \all{der Wald}, \all{die Dürre}, \all{der Müll}, \all{recyceln}, \all{sparen}) ;
\item employer correctement le \cle{subjonctif II} pour formuler des propositions, des souhaits et des hypothèses irréelles (\all{Wir könnten...}, \all{Wir sollten...}, \all{Wenn wir..., hätten wir...}) ;
\item construire des \cle{subordonnées de temps} (\all{als}, \all{wenn}, \all{nachdem}, \all{bevor}) et de \cle{conséquence} (\all{so dass} / \all{sodass}) avec le verbe conjugué en fin de proposition ;
\item travailler en groupe, rédiger collectivement une affiche-article et la présenter à l'oral en deux minutes ;
\item évaluer la production des autres groupes selon des critères simples (vocabulaire, grammaire, clarté).
\end{itemize}

\courssub{Situation}

Le proviseur de ton lycée à Yaoundé a lancé, à l'occasion de la Journée mondiale de l'environnement, un concours intitulé « \all{Unser Aktionsplan gegen die Umweltverschmutzung} ». Chaque classe de Terminale doit présenter un plan d'action concret pour rendre l'établissement plus propre et plus écologique : trop de plastique dans la cour, gaspillage de l'eau et de l'électricité, absence de tri des déchets. La classe d'allemand participe... en allemand ! Le meilleur plan sera affiché à l'entrée du lycée et envoyé à l'association partenaire allemande du lycée. Voici le message officiel du concours :

\begin{textebox}[Appel à projets du lycée — Wettbewerb]
\textbf{Liebe Schülerinnen und Schüler!}\par
Unsere Schule hat ein Problem: Es gibt zu viel Plastikmüll auf dem Schulhof, wir verschwenden Wasser und Strom, und wir trennen den Müll nicht. Die Erde wird wärmer, und auch in Kamerun sehen wir die Folgen: Die Dürren werden länger, und die Regenzeit kommt oft zu spät.\par
Deshalb fragen wir euch: Was könnten wir an unserer Schule ändern? Schreibt einen Aktionsplan! Erklärt zuerst kurz das Problem, dann die Folgen, wenn wir nichts tun, und macht am Ende drei konkrete Vorschläge.\par
Die beste Gruppe gewinnt einen Preis, und ihr Plan wird an der Schule ausgehängt. Viel Erfolg!\par
\emph{Die Schulleitung}
\end{textebox}

\textbf{Glossaire :} \all{der Plastikmüll} (les déchets plastiques) ; \all{verschwenden} (gaspiller) ; \all{trennen} (trier) ; \all{die Folge, -n} (la conséquence) ; \all{die Dürre, -n} (la sécheresse) ; \all{die Regenzeit, -en} (la saison des pluies) ; \all{der Vorschlag, „-e} (la proposition, la suggestion) ; \all{aushängen} (afficher) ; \all{der Preis, -e} (le prix, la récompense).

\courssub{Consignes}

Par groupes de quatre, vous rédigez une affiche-article intitulée « \all{Unser Aktionsplan gegen die Umweltverschmutzung} ». Suivez les étapes dans l'ordre.

\begin{enumerate}
\item \textbf{Relecture du texte support.} Relisez « Die Erde wird wärmer » et le message du concours ci-dessus. Soulignez cinq mots du champ lexical de l'environnement que vous pourrez réutiliser dans votre affiche. Vérifiez l'article et le pluriel de chaque nom dans votre lexique.

\item \textbf{Le problème (3 phrases).} Résumez le problème de votre lycée en trois phrases. Chaque phrase doit contenir une subordonnée de temps introduite par \all{als}, \all{wenn} ou \all{nachdem}. Attention : le verbe conjugué va à la fin de la subordonnée, et la principale qui suit commence par le verbe. Exemple de départ : \all{Als wir gestern den Schulhof besuchten, ...}

\item \textbf{Les conséquences (1 phrase).} Écrivez une phrase avec \all{so dass} (ou \all{sodass}) qui décrit ce qui arrivera si l'on ne fait rien. Exemple de structure : \all{Wenn wir nichts ändern, wird der Müll immer mehr, so dass ...}

\item \textbf{Les propositions (3 phrases au subjonctif II).} Formulez trois propositions concrètes pour votre lycée :
\begin{itemize}
\item une avec \all{Wir könnten...} (une possibilité) ;
\item une avec \all{Wir sollten...} (un devoir, un conseil) ;
\item une hypothèse complète : \all{Wenn wir ..., hätten wir ...} ou \all{..., würden wir ...}.
\end{itemize}

\item \textbf{Mise en page de l'affiche.} Donnez un titre visible, écrivez vos sept phrases lisiblement, et ajoutez deux ou trois mots-clés en gras (\all{Wasser sparen!}, \all{Müll trennen!}, \all{Bäume pflanzen!}). Relisez-vous : majuscules aux noms, place du verbe, formes du subjonctif II.

\item \textbf{Présentation orale.} Chaque groupe présente son affiche en deux minutes maximum. Répartissez la parole : un élève présente le problème, un autre les conséquences, deux autres les propositions. Parlez lentement et regardez la classe.

\item \textbf{Vote et évaluation.} La classe vote pour le meilleur plan. Critères : exactitude du vocabulaire thématique (5 points), correction du subjonctif II et des subordonnées (10 points), clarté de la présentation orale (5 points).
\end{enumerate}

\courssub{Ressources à mobiliser}

\begin{aretenir}[Le subjonctif II : les formes indispensables]
\begin{itemize}
\item \all{haben} $\to$ \all{hätten} : \all{wir hätten} « nous aurions » ;
\item \all{sein} $\to$ \all{wären} : \all{das wäre} « ce serait » ;
\item \all{können} $\to$ \all{könnten} : \all{wir könnten} « nous pourrions » ;
\item \all{sollen} $\to$ \all{sollten} : \all{wir sollten} « nous devrions » ;
\item \all{müssen} $\to$ \all{müssten} « nous devrions (obligation) » ; \all{dürfen} $\to$ \all{dürften} « nous aurions le droit » ;
\item pour les autres verbes : \all{würden} + infinitif : \all{wir würden mehr Bäume pflanzen} « nous planterions plus d'arbres ».
\end{itemize}
Dans la phrase hypothétique : \all{Wenn wir den Müll trennten, hätten wir einen sauberen Schulhof.} « Si nous triions les déchets, nous aurions une cour propre. » (subordonnée avec verbe conjugué à la fin, principale avec le verbe en première position après la virgule, sujet après le verbe).
\end{aretenir}

\begin{aretenir}[Les subordonnées de temps et de conséquence]
\begin{center}
\begin{tabularx}{\linewidth}{|l|X|X|}
\hline
\rowcolor{popblueL} \textbf{Conjonction} & \textbf{Sens} & \textbf{Exemple} \\
\hline
\all{als} & quand, lorsque (passé, action unique) & \all{Als wir ankamen, sahen wir den Müll.} \\
\hline
\all{wenn} & quand, si (présent, répétition, hypothèse) & \all{Wenn es regnet, gibt es weniger Staub.} \\
\hline
\all{nachdem} & après que & \all{Nachdem wir gegessen hatten, warfen wir den Müll weg.} \\
\hline
\all{bevor} & avant que & \all{Bevor wir gehen, löschen wir das Licht.} \\
\hline
\all{so dass} & si bien que, de sorte que & \all{Es gibt zu viel Plastik, so dass die Tiere krank werden.} \\
\hline
\end{tabularx}
\end{center}
Dans toutes ces subordonnées, le verbe conjugué se place \textbf{en dernière position}. Remarque orthographique : \all{so dass} peut aussi s'écrire en un seul mot, \all{sodass} ; les deux formes sont correctes.
\end{aretenir}

\begin{definition}[Lexique utile pour l'affiche]
\all{die Umweltverschmutzung} (la pollution de l'environnement) ; \all{die Erderwärmung} (le réchauffement climatique) ; \all{das Klima} (le climat) ; \all{das CO$_2$} (le CO$_2$) ; \all{der Wald, „-er} (la forêt) ; \all{die Dürre, -n} (la sécheresse) ; \all{der Müll} (les déchets) ; \all{die Mülldeponie, -n} (la décharge) ; \all{der Mülleimer, -} (la poubelle) ; \all{die Plastikflasche, -n} (la bouteille en plastique) ; \all{der Schulhof, „-e} (la cour de récréation) ; \all{Wasser sparen} (économiser l'eau) ; \all{den Müll trennen} (trier les déchets) ; \all{recyceln} (recycler) ; \all{Bäume pflanzen} (planter des arbres) ; \all{die Umwelt schützen} (protéger l'environnement) ; \all{sauber} (propre) ; \all{verschmutzt} (pollué).
\end{definition}

\begin{attention}[Pièges à éviter dans votre production]
\begin{itemize}
\item Ne dites jamais \all{wir würden können} : on utilise directement \all{wir könnten}. De même, évitez \all{wir würden haben} : dites \all{wir hätten}.
\item Après \all{so dass}, le verbe va à la fin : \all{..., so dass die Schule sauberer wird.} (et non \all{so dass die Schule wird sauberer}).
\item \all{nachdem} exige souvent le plus-que-parfait dans la subordonnée quand la principale est au passé : \all{Nachdem wir den Müll gesammelt hatten, gingen wir nach Hause.}
\item N'oubliez pas la majuscule à tous les noms : \all{der Wald}, \all{die Dürre}, \all{das Klima}.
\item Faux ami : \all{die Dürre} signifie « la sécheresse », pas « la dureté ». Autre piège : \all{das Klima} est neutre, pas féminin.
\item Verbes à particule dans une subordonnée : la particule se \textbf{resoude} au verbe en fin de proposition : \all{Wenn die Schüler weggehen, ...} (et non \all{gehen weg}).
\end{itemize}
\end{attention}

\begin{methode}[Réussir la présentation orale en quatre étapes]
\begin{enumerate}
\item \textbf{Préparer} : chaque membre du groupe apprend sa partie par cœur ou avec des notes très courtes ; chronométrez-vous (deux minutes maximum).
\item \textbf{Articuler} : prononcez clairement les sons difficiles : \all{Umweltverschmutzung} « oum-vèlt-fer-chmou-tsoung », \all{Erderwärmung} « ér-deu-vèr-moung », \all{Dürre} « du-ré ».
\item \textbf{Regarder} : ne lisez pas l'affiche ; regardez la classe et désignez les mots-clés en gras au bon moment.
\item \textbf{Conclure} : terminez par une phrase forte, par exemple : \all{Gemeinsam können wir unsere Schule sauber machen!} « Ensemble, nous pouvons rendre notre école propre ! »
\end{enumerate}
\end{methode}

\courssub{Proposition de production et corrigé commenté}

\begin{exoresolu}[Modèle d'affiche-article : « Unser Aktionsplan gegen die Umweltverschmutzung »]
Rédigez l'affiche complète de votre groupe (résumé du problème, conséquences, trois propositions) selon les consignes.
\tcblower
\textbf{Unser Aktionsplan gegen die Umweltverschmutzung}\par
\emph{Das Problem.} Als wir letzte Woche den Schulhof besuchten, sahen wir überall Plastikflaschen und alte Tüten. Wenn die Schüler nach der Pause weggehen, bleibt der Müll auf dem Boden liegen. Nachdem der Wind gestern stark geweht hatte, lag der Plastikmüll sogar vor den Klassenzimmern.\par
\emph{Die Folgen.} Wenn wir nichts ändern, wird der Müll immer mehr, so dass unsere Schule bald wie eine Mülldeponie aussieht und die Erderwärmung noch schlimmer wird.\par
\emph{Unsere Vorschläge.} Erstens könnten wir in jedem Klassenzimmer zwei Mülleimer aufstellen: einen für Plastik und einen für Papier. Zweitens sollten wir jeden Freitag eine „Grüne Stunde“ organisieren, in der alle Schüler den Schulhof sauber machen und Bäume pflanzen. Drittens: Wenn wir weniger Wasser und Strom verschwendeten, hätten wir mehr Geld für Bücher und Sport.\par
\textbf{Wasser sparen! Müll trennen! Bäume pflanzen!}\par\medskip
\textbf{Commentaire des choix de langue.} (1) \all{Als wir ... besuchten} : subordonnée de temps au passé avec \all{als} (action unique), verbe conjugué à la fin ; la principale commence par le verbe (\all{sahen wir}). (2) \all{Wenn die Schüler ... weggehen} : \all{wenn} pour une situation répétée au présent ; le verbe à particule \all{weggehen} reste soudé en fin de subordonnée. (3) \all{Nachdem der Wind ... geweht hatte} : plus-que-parfait correctement employé après \all{nachdem}, la principale étant au prétérit (\all{lag}). (4) \all{so dass unsere Schule ... aussieht} : conséquence avec verbe final ; \all{aussehen} (ressembler) est un verbe à particule séparable, resoudé ici en fin de subordonnée. (5) \all{könnten wir ... aufstellen} : subjonctif II de \all{können} + infinitif du verbe à particule en fin de phrase. (6) \all{sollten wir ... organisieren} : subjonctif II de \all{sollen} pour exprimer le conseil. (7) \all{Wenn wir ... verschwendeten, hätten wir ...} : hypothèse irréelle complète au subjonctif II, avec inversion du sujet dans la principale. (8) Le lexique thématique est bien présent : \all{der Plastikmüll}, \all{die Erderwärmung}, \all{die Mülldeponie}, \all{Bäume pflanzen}, \all{Wasser sparen}, \all{Müll trennen}.
\end{exoresolu}

\begin{exercice}[À vous de jouer : entraînement rapide avant l'affiche]
\begin{enumerate}
\item Mettez les phrases suivantes au subjonctif II : \all{Wir können den Müll trennen.} / \all{Wir sollen Wasser sparen.} / \all{Wir haben mehr Zeit für Sport.}
\item Complétez avec \all{als} ou \all{wenn} : \all{... ich gestern nach Hause ging, regnete es.} / \all{... es im März regnet, sind die Bauern froh.}
\item Remettez les mots dans l'ordre : \all{der Müll / so dass / ist / überall / die Tiere / krank / werden}.
\end{enumerate}
\end{exercice}

\begin{corrige}[Exercice : entraînement rapide]
\begin{enumerate}
\item \all{Wir könnten den Müll trennen.} / \all{Wir sollten Wasser sparen.} / \all{Wir hätten mehr Zeit für Sport.}
\item \all{Als ich gestern nach Hause ging, regnete es.} (action unique au passé) / \all{Wenn es im März regnet, sind die Bauern froh.} (répétition, présent).
\item \all{Der Müll ist überall, so dass die Tiere krank werden.} (verbe conjugué \all{werden} en dernière position de la subordonnée de conséquence).
\end{enumerate}
\end{corrige}

\courssub{Bilan de l'activité}

Cette activité a permis de réinvestir l'ensemble des acquis de la leçon dans une tâche de communication réelle et motivante. Sur le plan de la \cle{compréhension}, les élèves ont relu le texte support et le message du concours pour y repérer les informations et le lexique réutilisables. Sur le plan du \cle{lexique}, le vocabulaire de l'environnement a été employé dans un contexte concret et proche de la vie des élèves : la cour du lycée, le plastique, l'eau, les arbres. Sur le plan de la \cle{grammaire}, les deux points du programme ont été systématiquement mobilisés : le subjonctif II pour les propositions (\all{könnten}, \all{sollten}, \all{hätten}, \all{würden}) et les subordonnées de temps (\all{als}, \all{wenn}, \all{nachdem}) et de conséquence (\all{so dass}) pour articuler le raisonnement. Enfin, la présentation orale et le vote final ont développé l'\cle{expression orale} et l'esprit de coopération. L'enseignant veillera à corriger collectivement les erreurs les plus fréquentes relevées pendant les présentations (place du verbe, formes du subjonctif II, majuscules des noms) et pourra afficher le plan gagnant dans la salle de classe comme trace écrite de la leçon.

\newpage

\coursec{Méthodes et exercices résolus}

\coursec{Leçon AL13 : La pollution et le réchauffement climatique}

\courssub{Texte support : „Die Erde fiebert – Auch Kamerun spürt die Hitze“}

\begin{textebox}[Article simulé, Yaoundé, mars 2024]
„Die Erde fiebert.“ Dieser Satz der Klimaforscher gilt auch für Kamerun. Seit hundert Jahren steigt die globale Durchschnittstemperatur stetig an. Der Hauptgrund ist der vom Menschen gemachte Treibhauseffekt: Durch die Verbrennung von Kohle, Öl und Gas gelangt zu viel CO$_2$ in die Atmosphäre. In Kamerun sind die Folgen bereits sichtbar. Im Norden führen längere Dürreperioden zu Ernteausfällen und Wasserknappheit. In Douala verursachen extreme Regenfälle immer häufiger Überschwemmungen. Die Abholzung der Regenwälder für Kakaoplantagen und Brennholz verschärft das Problem, weil die Wälder als CO$_2$-Speicher verloren gehen. Gleichzeitig verschmutzt Plastikmüll die Flüsse und das Meer. Doch es gibt Hoffnung: Junge Kamerunerinnen und Kameruner pflanzen Bäume, fördern Solarenergie und sammeln Müll. „Wenn wir jetzt handeln, können wir das Schlimmste verhindern“, sagt eine Schülerin in Yaoundé.
\end{textebox}

\begin{aretenir}[Glossaire du texte support]
\begin{center}
\begin{tabularx}{\linewidth}{|>{\bfseries}l|X|X|X|}
\hline
\rowcolor{popblueL} \textbf{Deutsch (Artikel, Plural)} & \textbf{Français} & \textbf{Famille de mots / Collocations} & \textbf{Remarque} \\
\hline
die Erderwärmung, -en & le réchauffement climatique & die globale Erderwärmung, sich erwärmen (V) & Nom composé : Erde + Erwärmung \\
\hline
der Klimawandel & le changement climatique & der menschengemachte Klimawandel, klimaneutral (Adj) & Kein Plural üblich \\
\hline
der Treibhauseffekt, -e & l'effet de serre & der anthropogene Treibhauseffekt & Treibhaus = serre \\
\hline
das CO$_2$ (Kohlendioxid) & le dioxyde de carbone & der CO$_2$-Ausstoß, CO$_2$-neutral & Symbole sont neutres (das) \\
\hline
die Dürre, -n & la sécheresse & eine schwere Dürre, dürreperioden & \emph{Nicht} die Trockenheit (terme générique) \\
\hline
die Überschwemmung, -en & l'inondation & eine Überschwemmung verursachen, überschwemmen (V) & Verbe : überschwemmen (faible) \\
\hline
die Abholzung, -en & la déforestation / le déboisement & abholzen (V, stark: hieb, gehauen), der Kahlschlag & Particule séparable : abholzen \\
\hline
der Wald, die Wälder & la forêt & der Regenwald, der Mischwald, walden (V) & Umlaut au pluriel ! \\
\hline
das Plastik / der Kunststoff & le plastique & Plastikmüll, plastikfrei, recyceln & Kein Plural (Masse) \\
\hline
der Müll & les déchets / ordures & Müll trennen, Müll sammeln, der Mülleimer & Kein Plural (collectif) \\
\hline
verbrannte, verbrannt (verbrennen) & brûler (tr.) & Kohle verbrennen, die Verbrennung & Verbe fort : a/o/u $\to$ ä/o/ü (Présent: er verbrennt) \\
\hline
ausbleiben, blieb aus, ausgeblieben & faire défaut / ne pas venir & die Regenzeit bleibt aus & Verbe fort + particule séparable \\
\hline
bedrohen & menacer & die bedrohten Arten, die Bedrohung & Faible \\
\hline
verhindern & empêcher / prévenir & das Schlimmste verhindern, die Prävention & Faible \\
\hline
\end{tabularx}
\end{center}
\end{aretenir}

\courssub{Savoir-faire 1 : Comprendre un texte sur l'environnement (global puis détaillé)}

\begin{methode}[Lire et exploiter un texte au Bac — 5 étapes]
\begin{enumerate}
\item \textbf{Lecture globale (2 min).} Lis le titre, la source, la date et le premier/dernier paragraphe. Repère les \cle{mots-clés} du champ lexical (die Erderwärmung, der Treibhauseffekt, die Dürre, die Überschwemmung, die Abholzung, Solarenergie). Formule l'idée centrale : \all{Der Text thematisiert die Ursachen und Folgen des Klimawandels in Kamerun und zeigt Lösungsansätze.}
\item \textbf{Lecture détaillée (5 min).} Souligne : \textbf{Chiffres/Dates} (seit hundert Jahren), \textbf{Causes} (introduites par \all{weil, da, durch + Akkusativ}), \textbf{Conséquences} (\all{deshalb, so dass, führen zu}), \textbf{Solutions/Opinions} (\all{sollen, würden, meiner Meinung nach}).
\item \textbf{Analyse des questions.} Identifie le type : \emph{Global} (Worum geht es?), \emph{Détail} (Was sagt der Text zu...?), \emph{Inférence} (Warum...?), \emph{Vocabulaire} (Erklären Sie...).
\item \textbf{Rédaction des réponses.} \textbf{Ne citez pas}, reformulez ! Utilisez : \all{Laut dem Text...} (D'après le texte), \all{Der Autor schreibt, dass...} (L'auteur écrit que...), \all{Es wird deutlich, dass...} (Il apparaît que...). Verbe conjugué en 2\textsuperscript{e} place, noms en majuscule.
\item \textbf{Relecture (1 min).} Vérifiez : majuscules (noms), Umlaute (ä, ö, ü), ß, place du verbe (principal / subordonnée), accords (cas, genre, nombre).
\end{enumerate}
\end{methode}

\begin{exoresolu}[Compréhension détaillée — Texte support „Die Erde fiebert“]
\textbf{Questions :}
1. Worum geht es in dem Text? (Sujet global)
2. Nennen Sie zwei Ursachen für die Erderwärmung, die im Text genannt werden. (Causes)
3. Welche konkreten Folgen des Klimawandels gibt es in Kamerun? (Conséquences locales)
4. Was bedeutet „die Erde fiebert“ im übertragenen Sinn? (Métaphore / Vocabulaire)
5. Nennen Sie zwei Lösungsansätze, die am Ende des Textes erwähnt werden. (Solutions)
\tcblower
\textbf{Raisonnement :} Q1 = synthèse. Q2 = repérer \emph{weil / durch}. Q3 = repérer \emph{In Kamerun... führen zu / verursachen}. Q4 = expliquer l'image (fièvre = température anormale). Q5 = dernier paragraphe (Jeunes, Solarenergie, Müll sammeln).\par
\textbf{Réponses modèles :}\par
1. \all{Der Text handelt von den Ursachen und Folgen des Klimawandels speziell in Kamerun sowie von möglichen Gegenmaßnahmen.}\par
2. \all{Die Hauptursachen sind die Verbrennung fossiler Brennstoffe (Kohle, Öl, Gas) und die Abholzung der Regenwälder.}\par
3. \all{Im Norden gibt es längere Dürren und Wasserknappheit, in Douala häufigere Überschwemmungen durch extreme Regenfälle.}\par
4. \all{Die Metapher bedeutet, dass die Erde „krank“ ist: ihre Temperatur steigt pathologisch an, genau wie bei Fieber.}\par
5. \all{Junge Menschen pflanzen Bäume, setzen auf Solarenergie und organisieren Müllsammelaktionen.}\par
\textbf{Auto-évaluation :} Mes réponses sont-elles complètes ? Les verbes sont-ils conjugués ? Les noms ont-ils une majuscule ?
\end{exoresolu}

\courssub{Savoir-faire 2 : Maîtriser le subjonctif II (l'irréel) — Konjunktiv II}

\begin{propriete}[Formation du Konjunktiv II — Règle générale et exceptions]
\begin{center}
\begin{tabularx}{\linewidth}{|>{\bfseries}l|X|X|X|}
\hline
\rowcolor{popblueL} \textbf{Type de verbe} & \textbf{Präteritum} & \textbf{Konjunktiv II (Présent)} & \textbf{Règle / Remarque} \\
\hline
\all{sein} (être) & \all{war} & \all{wäre} & \cle{Forme de base irrégulière} — \textbf{Jamais} *würde sein* \\
\hline
\all{haben} (avoir) & \all{hatte} & \all{hätte} & \cle{Forme de base irrégulière} — \textbf{Jamais} *würde haben* \\
\hline
\all{werden} (devenir) & \all{wurde} & \all{würde} & Auxiliaire pour la forme \emph{würde + infinitif} \\
\hline
\textbf{Modaux} & \all{kannte, musste, wollte...} & \all{könnte, müsste, wollte, dürfte, möchte} & Umlaut systématique + terminaisons \emph{-e, -est, -e, -en, -et, -en} \\
\hline
\textbf{Verbes forts} (ex. \all{geben, kommen, verbrennen}) & \all{gab, kam, verbrannte} & \all{gäbe, käme, verbrennte} & \textbf{Radical Präteritum + Umlaut (a/o/u $\to$ ä/ö/ü) + terminaisons} \\
\hline
\textbf{Verbes faibles} (ex. \all{machen, schützen, pflanzen}) & \all{machte, schützte, pflanzte} & \all{machte, schützte, pflanzte} & \textbf{Identique au Präteritum} $\to$ \textbf{On préfère \all{würde + infinitif}} \\
\hline
\end{tabularx}
\end{center}
\end{propriete}

\begin{methode}[Former et employer le Konjunktiv II au Bac]
\begin{enumerate}
\item \textbf{Forme de secours universelle.} Pour \textbf{tous les verbes} (sauf \all{sein, haben, modaux}) : \all{würde} (conjugué) + \textbf{infinitif à la fin}. \all{Ich würde Bäume pflanzen.} (Plus sûr que *Ich pflanzte*).
\item \textbf{Trois emplois majeurs :}
\begin{itemize}
\item \textbf{Irréel du présent (condition) :} \all{Wenn ich Zeit \textbf{hätte}, würde ich die Umwelt \textbf{schützen}.} (\all{wenn} + KII à la fin ; principale : verbe en 1\textsuperscript{re} place).
\item \textbf{Souhait irréalisable :} \all{Wenn das Klima \textbf{nicht so heiß wäre}!} (Souvent sans principale, ton exclamatif).
\item \textbf{Conseil / Politesse :} \all{Du \textbf{solltest} weniger Plastik kaufen.} / \all{Ich \textbf{würde} an deiner Stelle recyceln.}
\end{itemize}
\item \textbf{Ordre des mots (Satzklammer) :}
\begin{itemize}
\item Subordonnée \all{wenn} : \all{Wenn ich Minister \textbf{wäre} ...} (verbe tout à la fin).
\item Principale qui suit : \all{... \textbf{würde} ich Solarparks \textbf{bauen}.} (verbe auxiliaire en 1\textsuperscript{re} place, infinitif à la fin).
\end{itemize}
\end{enumerate}
\end{methode}

\begin{popfigure}
\begin{tikzpicture}[node distance=1.2cm, >=Stealth,
  box/.style={draw=popblue, thick, rounded corners, fill=popblueL, text width=3.5cm, align=center, font=\small},
  arrow/.style={->, thick, popdark}]
\node[box, align=center] (si){\textbf{Si-clause (Condition)}\\ \all{Wenn} + Sujet + ... + \textbf{KII (fin)}};
\node[box, right=of si, align=center] (alors){\textbf{Hauptsatz (Conséquence)}\\ \textbf{KII (1\textsuperscript{re} place)} + Sujet + ... + Infinitif};
\draw[arrow] (si.east) -- node[above, font=\tiny] {ordre inversé si \all{wenn} en tête} (alors.west);
\end{tikzpicture}
\legende{Structure de la phrase conditionnelle irréelle (Konjunktiv II).}
\end{popfigure}

\begin{exoresolu}[Thème — Konjunktiv II (Français $\to$ Allemand)]
Traduisez : « Si nous protégions les forêts, le climat serait meilleur. Malheureusement, les gens brûlent trop de bois. Si j'étais ministre, je construirais plus de parcs solaires. »
\tcblower
\textbf{Analyse grammaticale :}
\begin{itemize}
\item Phrase 1 : Condition irréelle présente. \all{schützen} (faible) $\to$ \all{würden schützen}. \all{sein} $\to$ \all{wäre}.
\item Phrase 2 : Réel (constat) $\to$ \all{Präsens} : \all{verbrennen} (fort, présent \all{er verbrennt}) $\to$ \all{verbrennen}.
\item Phrase 3 : Condition irréelle. \all{sein} $\to$ \all{wäre}. \all{bauen} (faible) $\to$ \all{würde bauen}.
\end{itemize}
\textbf{Solution :}\par
\all{Wenn wir die Wälder schützen \textbf{würden}, \textbf{wäre} das Klima besser. Leider verbrennen die Menschen zu viel Holz. Wenn ich Minister \textbf{wäre}, \textbf{würde} ich mehr Solarparks \textbf{bauen}.}\par
\textbf{Justifications clés :} \all{die Wälder} (pluriel de \all{der Wald}, Umlaut) ; \all{wären} (KII \all{sein}) ; \all{würden ... schützen} (forme de secours, infinitif final) ; \all{würde ... bauen} (principal commence par \all{würde} car subordonnée \all{wenn} en tête).
\end{exoresolu}

\begin{exercice}[Entraînement Konjunktiv II]
Mettez les verbes au Konjunktiv II (forme appropriée : \all{wäre/hätte/könnte} ou \all{würde + infinitif}).
1. Wenn ich du \_\_\_\_\_\_\_\_\_\_ (sein), ich \_\_\_\_\_\_\_\_\_\_ (nicht / Müll / wegwerfen).
2. Die Bauern \_\_\_\_\_\_\_\_\_\_ (können) mehr anbauen, wenn es mehr \_\_\_\_\_\_\_\_\_\_ (regnen).
3. \_\_\_\_\_\_\_\_\_\_ (haben) wir mehr Solarenergie, \_\_\_\_\_\_\_\_\_\_ (sein) die Luft sauberer.
\end{exercice}

\begin{corrige}[Exercice Konjunktiv II]
1. \all{wäre / würde ... nicht wegwerfen} (sein $\to$ wäre ; wegwerfen faible $\to$ würde wegwerfen).\par
2. \all{könnten / regnete} (können modal $\to$ könnten ; regnen fort $\to$ regnete [ou würde regnen]).\par
3. \all{Hätten / wäre} (haben $\to$ hätten ; sein $\to$ wären). \emph{Note :} \all{Hätten wir...} = inversion sujet/verbe (style soutenu) = \all{Wenn wir mehr Solarenergie hätten...}.
\end{corrige}

\courssub{Savoir-faire 3 : Construire des subordonnées de conséquence et de temps}

\begin{propriete}[Conjonctions de conséquence et de temps — Syntaxe et emploi]
\begin{center}
\begin{tabularx}{\linewidth}{|>{\bfseries}l|l|X|X|}
\hline
\rowcolor{popblueL} \textbf{Fonction} & \textbf{Connecteur} & \textbf{Structure / Temps} & \textbf{Exemple (Environnement)} \\
\hline
\textbf{Conséquence} (Conj.) & \all{so dass / sodass} & Subordonnée : verbe à la fin. & \all{Es regnet stark, \textbf{so dass} die Flüsse \textbf{überlaufen}.} \\
\hline
\textbf{Conséquence} (Adv.) & \all{deshalb / deswegen / darum} & \textbf{Principale} : Verbe en 2\textsuperscript{e} place (V2). & \all{Es regnet stark. \textbf{Deshalb} \textbf{überlaufen} die Flüsse.} \\
\hline
\textbf{Conséquence négative} & \all{zu ... als dass} & \all{zu} + Adj/Adv + \all{als dass} + KII (verbe fin). & \all{Es ist \textbf{zu heiß}, \textbf{als dass} man \textbf{draußen arbeiten könnte}.} \\
\hline
\textbf{Temps : Passé unique} & \all{als} & Subordonnée : verbe à la fin (Präteritum/Perfekt). & \all{\textbf{Als} die Dürre \textbf{kam}, \textbf{verließen} viele Bauern ihre Dörfer.} \\
\hline
\textbf{Temps : Présent / Futur / Répétition passé} & \all{wenn} & Subordonnée : verbe à la fin. & \all{\textbf{Wenn} es heiß \textbf{ist}, trinken wir viel. / \textbf{Wenn} ich Kind \textbf{war}, gab es mehr Regen.} \\
\hline
\textbf{Simultanéité} & \all{während} & Subordonnée : verbe à la fin. & \all{\textbf{Während} die Sonne \textbf{scheint}, arbeiten die Solarzellen.}\\
\hline
\textbf{Antériorité} & \all{bevor / nachdem} & \all{bevor} + Infinitiv (sujet identique) OU \all{nachdem} + Plusquamperfekt. & \all{\textbf{Bevor} wir \textbf{gehen}, pflanzen wir einen Baum. / \textbf{Nachdem} der Regen \textbf{aufgehört hatte}, kamen wir raus.} \\
\hline
\end{tabularx}
\end{center}
\end{propriete}

\begin{attention}[Piège n°1 : \all{deshalb} $\neq$ \all{weil/da/so dass}]
\all{deshalb} est un \textbf{adverbe} (Konnektor), pas une conjonction. Il ne renvoie pas le verbe à la fin.
\begin{center}
\textbf{Faux :} \all{Es regnet, deshalb die Straßen überflutet sind.}\par
\textbf{Juste :} \all{Es regnet, deshalb \textbf{sind} die Straßen überflutet.} (V2)\par
\textbf{Juste :} \all{Es regnet, \textbf{so dass} die Straßen überflutet \textbf{sind}.} (Verbe fin)
\end{center}
\end{attention}

\begin{attention}[Piège n°2 : \all{als} vs \all{wenn} au passé]
\all{als} = \textbf{un seul événement ponctuel} dans le passé (storytelling). \all{wenn} = \textbf{répétition} (chaque fois que) au passé, ou présent/futur.
\begin{center}
\textbf{Faux :} \all{Wenn die Dürre 2022 kam, ...} (événement unique)\par
\textbf{Juste :} \all{\textbf{Als} die Dürre 2022 \textbf{kam}, ...}\par
\textbf{Juste :} \all{\textbf{Wenn} es im Sommer \textbf{heiß war}, \textbf{blieben} wir zu Hause.} (répétition habituelle)
\end{center}
\end{attention}

\begin{exoresolu}[Transformation et complétion — Conséquence \& Temps]
1. Reliez avec \all{so dass} : \all{Die Temperatur stieg sehr stark. Viele Tiere starben.}\par
2. Complétez : a) \all{\_\_\_\_\_\_ ich ein Kind war, gab es mehr Regen.} b) \all{\_\_\_\_\_\_ es im Sommer heiß ist, trinken wir viel Wasser.}\par
3. Traduisez : « Quand la sécheresse est arrivée, les récoltes ont été détruites. »\par
4. Reformulez avec \all{zu ... als dass} : \all{Der Müll ist so viel, dass man ihn nicht verbrennen kann.}
\tcblower
\textbf{Solutions :}\par
1. \all{Die Temperatur stieg \textbf{so stark, dass} viele Tiere \textbf{starben}.} (Präteritum, verbe \all{starben} à la fin).\par
2. a) \all{\textbf{Als} ich ein Kind \textbf{war}, gab es mehr Regen.} (Passé unique).\quad b) \all{\textbf{Wenn} es im Sommer heiß \textbf{ist}, trinken wir viel Wasser.} (Présent / Habitude).\par
3. \all{\textbf{Als} die Dürre \textbf{kam}, \textbf{wurden} die Ernten \textbf{zerstört}.} (\all{die Dürre}, \all{die Ernte, -n} ; Passif Präteritum \all{wurden zerstört} ; \all{kam} fin subordonnée $\to$ \all{wurden} tête principale).\par
4. \all{Der Müll ist \textbf{zu viel, als dass} man ihn \textbf{verbrennen könnte}.} (KII \all{könnte} obligatoire après \all{als dass}).
\end{exoresolu}

\courssub{Savoir-faire 4 : Traduire un texte sur le climat (Version \& Thème)}

\begin{methode}[Réussir la Version (Allemand $\to$ Français) et le Thème (Français $\to$ Allemand)]
\begin{enumerate}
\item \textbf{Version :} Lisez 2x. Repérez le temps dominant (Présens/Perfekt presse ; Präteritum récit). Segmentez : Sujet (Nominatif) / Verbe / Compléments (Accusatif/Datif/Génitif). Traduisez le \emph{sens}, pas mot à mot. \cle{Faux amis} : \all{bekommen} = recevoir (\all{werden} = devenir) ; \all{die Aktualität} $\neq$ actualité (journal) ; \all{die Umgebung} = les alentours (pas \all{die Umwelt} = l'environnement).
\item \textbf{Thème :} Pensez \emph{allemand} : Cas (qui ? quoi ? à qui ?), Prépositions (Wechselpräpositionen : \all{in} + Acc = mouvement / + Dat = lieu), Verbes à particule (\all{ausbleiben, wegwerfen, anbauen, aufstellen}), Place du verbe (V2 / Fin / Infinitif).
\item \textbf{Relecture croisée :} Version $\to$ français naturel ? Thème $\to$ majuscules noms, Umlaute, ß, verbe 2\textsuperscript{e} place / fin ?
\end{enumerate}
\end{methode}

\begin{exoresolu}[Version — Type Bac]
Traduisez en français : \all{Wenn die Regenzeit ausbleibt, können die Bauern in der Nordregion Kameruns kein Gemüse mehr anbauen. Deshalb müssen viele Familien ihre Dörfer verlassen und in die Städte ziehen.}
\tcblower
\textbf{Analyse :} \all{wenn} + \all{ausbleibt} (particule \all{aus} à la fin) = « Quand la saison des pluies fait défaut ». \all{können ... anbauen} (particule \all{an} à la fin) = « peuvent cultiver ». \all{kein ... mehr} = « plus aucun / plus de ». \all{deshalb} = adverbe $\to$ V2 (\all{müssen}). \all{verlassen} (Acc) / \all{ziehen} (in + Acc = mouvement vers) = « quitter / s'installer ».\par
\textbf{Traduction :} « Quand la saison des pluies fait défaut, les paysans de la région du Nord du Cameroun ne peuvent plus cultiver de légumes. C'est pourquoi de nombreuses familles doivent quitter leurs villages et s'installer dans les villes. »
\end{exoresolu}

\begin{exoresolu}[Thème — Type Bac]
Traduisez en allemand : « Parce que le climat change, les sécheresses deviennent plus longues. Si nous ne protégions pas la nature, nous détruirions notre avenir. C'est pourquoi nous devrions planter des arbres et utiliser l'énergie solaire. »
\tcblower
\textbf{Analyse :} P1 Cause : \all{weil/da} + verbe fin. \all{sich ändern} (réfléchi). P2 Condition irréelle : \all{wenn} + KII (\all{würden ... schützen}), principale \all{würden ... zerstören}. P3 Conséquence/Conseil : \all{deshalb} + V2 (\all{sollten}), \all{und} coordination (\all{nutzen}).\par
\textbf{Solution :}\par
\all{Weil sich das Klima \textbf{ändert}, \textbf{werden} die Dürren länger. \textbf{Wenn} wir die Natur \textbf{nicht schützen würden}, \textbf{würden} wir unsere Zukunft \textbf{zerstören}. \textbf{Deshalb} \textbf{sollten} wir Bäume \textbf{pflanzen} und Solarenergie \textbf{nutzen}.}\par
\textbf{Points clés :} \all{sich ändern} (réfléchi) ; \all{die Dürre, -n} ; \all{nicht} avant l'infinitif groupé ; \all{deshalb} $\to$ V2 (\all{sollten}) ; \all{Bäume} (pluriel \all{der Baum}) ; \all{Solarenergie} (f, pas de pluriel ici).
\end{exoresolu}

\courssub{Savoir-faire 5 : Produire un court paragraphe argumenté (Écrit \& Oral)}

\begin{methode}[Rédiger 6–8 phrases cohérentes — Plan « Constat $\to$ Causes $\to$ Conséquences $\to$ Solutions (KII) »]
\begin{enumerate}
\item \textbf{Constat (1 phr.) :} \all{In Kamerun ist die Umweltverschmutzung ein großes Problem.}
\item \textbf{Causes (1–2 phr. \all{weil/da}) :} \all{Weil viele Menschen Plastik in die Flüsse werfen, ist das Wasser verschmutzt. Da die Wälder abgeholzt werden, fehlt der CO$_2$-Speicher.}
\item \textbf{Conséquences (1–2 phr. \all{deshalb/so dass}) :} \all{Deshalb sind die Fische und die Gesundheit der Menschen bedroht. Die Erderwärmung führt dazu, dass die Ernten ausbleiben.}
\item \textbf{Solutions personnelles (2 phr. KII / Modaux) :} \all{Meiner Meinung nach sollten wir weniger Plastik kaufen und mehr recyceln. Wenn ich Bürgermeister wäre, würde ich mehr Mülleimer aufstellen und Solarenergie fördern.}
\item \textbf{Conclusion (1 phr.) :} \all{Nur gemeinsam können wir unsere Umwelt schützen.}
\end{enumerate}
\textbf{Check-list finale :} Majuscules (Noms) ? Verbe 2\textsuperscript{e} place (principal) / Fin (subordonnée) ? Umlaute (ä, ö, ü) ? ß (après voyelle longue/diphtongue) ? Particules séparables à la fin ? Cas (Acc/Dat après prépositions) ?
\end{methode}

\begin{exoresolu}[Production écrite guidée — Sujet d'examen]
\textbf{Sujet :} \all{Was können wir in Kamerun gegen die Umweltverschmutzung tun?} (6–8 phrases)
\tcblower
\textbf{Modèle de réponse (niveau Bac A2/B1) :}\par
\all{In vielen Städten Kameruns gibt es zu viel Müll auf den Straßen und in den Flüssen. \textbf{Weil} viele Menschen Plastikflaschen und Tüten einfach \textbf{wegwerfen}, \textbf{werden} die Gewässer stark \textbf{verschmutzt}. \textbf{Deshalb} \textbf{sterben} Fische, und die Gesundheit der Menschen \textbf{ist} gefährdet. \textbf{Meiner Meinung nach} \textbf{sollten} wir alle \textbf{weniger Plastik kaufen} und Müll \textbf{trennen}. \textbf{Wenn} die Regierung \textbf{mehr Recyclingzentren bauen würde}, \textbf{könnte} man den Müll besser \textbf{verwerten}. \textbf{Wenn ich} Umweltminister \textbf{wäre}, \textbf{würde} ich ein Gesetz gegen Plastiktüten \textbf{erlassen}. \textbf{Zum Schluss} : \textbf{Nur} wenn alle mitmachen, \textbf{können} wir unsere Erde \textbf{retten}.}\par
\textbf{Traduction :} « Dans de nombreuses villes du Cameroun, il y a trop de déchets dans les rues et les rivières. Parce que beaucoup de gens jettent simplement bouteilles et sacs en plastique, les cours d'eau sont fortement pollués. C'est pourquoi les meurent et la santé des gens est menacée. À mon avis, nous devrions tous acheter moins de plastique et trier les déchets. Si le gouvernement construisait plus de centres de recyclage, on pourrait mieux valoriser les déchets. Si j'étais ministre de l'Environnement, j'édicterais une loi contre les sacs plastiques. En conclusion : seulement si tous participent, nous pourrons sauver notre Terre. »\par
\textbf{Validation grammaticale :} \all{wegwerfen} (séparable, fin subordonnée) ; Passif \all{werden verschmutzt} ; \all{deswegen/deshalb} + V2 ; \all{sollten} (KII modal) + infinitifs fin ; \all{wenn} + KII (\all{würde bauen / wäre / würde erlassen}) ; \all{Nur wenn}... inversion (verbe 1\textsuperscript{re} place).
\end{exoresolu}

\begin{experience}[Expression orale — Débat / Jeu de rôle (10–15 min)]
\textbf{Consigne :} En binôme. L'un est un \textbf{journaliste} (Radio Environnement), l'autre un \textbf{expert local} (agriculteur, ONG, maire, élève).
\textbf{Thème :} \all{„Klimawandel in unserer Region – Was tun?“}
\textbf{Structure :}
1. Journaliste : \all{„Guten Tag. Was beobachten Sie bei uns vor Ort?“}
2. Expert : \all{„Ich beobachte, dass... (constat + chiffres approximatifs). Die Ursache ist... (weil/da).“}
3. Journaliste : \all{„Und die Folgen für die Bevölkerung?“}
4. Expert : \all{„Die Folgen sind... (deshalb/so dass). Besonders betroffen sind...“}
5. Journaliste : \all{„Was schlagen Sie vor? Was würden Sie an meiner Stelle tun?“}
6. Expert : \all{„Meiner Meinung nach sollte man... (sollen). Wenn ich die Macht hätte, würde ich... (KII).“}
\textbf{Évaluation par les pairs :} Fluidité, vocabulaire du chapitre, KII correct, place du verbe, prononciation (\all{ü, ö, ä, ch, sch, z, v/w}).
\end{experience}

\begin{attention}[Les 5 erreurs fatales au Bac — Récapitulatif]
\begin{enumerate}
\item \textbf{Majuscules oubliées :} \all{die erderwärmung, der wald, das plastik} $\to$ \textbf{Zéro tolérance}. \cle{Tout nom = Majuscule}. (\all{die Erderwärmung, der Wald, das Plastik}).
\item \textbf{Verbe pas à la fin dans la subordonnée :} \all{weil die Menschen verbrennen viel Öl} $\to$ \all{weil die Menschen viel Öl \textbf{verbrennen}}. (\all{weil, da, dass, wenn, als, so dass, bevor, nachdem...} = verbe \textbf{final}).
\item \textbf{Calque « si j'aurais » / Double conditionnel :} \all{Wenn ich reich werden würde, würde ich...} $\to$ \all{Wenn ich reich \textbf{wäre}, \textbf{würde} ich...}. Subordonnée \all{wenn} = KII (\all{wäre/hätte/könnte} ou \all{verb+Umlaut}), Principale = KII (\all{würde + inf}).
\item \textbf{Confusion \all{als} / \all{wenn} (passé) :} Événement unique passé = \all{Als} (\all{Als der Sturm kam...}). Répétition passée = \all{Wenn} (\all{Wenn es regnete, blieben wir drin}).
\item \textbf{Faux amis \& Lexique approximatif :} \all{bekommen} $\neq$ devenir (\all{werden}) ; \all{die Umgebung} $\neq$ l'environnement (\all{die Umwelt}) ; \all{die Trockenheit} $\neq$ la sécheresse climatique (\all{die Dürre}) ; \all{die Aktualität} n'existe pas (actualité = \all{die Aktualität} n'est pas utilisé, dire \all{die Nachrichten, die Lage}).
\end{enumerate}
\end{attention}

\begin{savaistu}[Cameroun \& Allemagne : Coopération verte]
Le Cameroun et l'Allemagne coopèrent dans le cadre de la \all{Entwicklungszusammenarbeit} (coopération au développement). La \all{GIZ} (Deutsche Gesellschaft für Internationale Zusammenarbeit) appuie des projets d'\all{Agroforstwirtschaft} (agroforesterie) dans le Nord et l'Est, de protection de la biodiversité (ex. Parc national de la Bénoué) et d'\all{Erneuerbare Energien} (énergies renouvelables : solaire, hydroélectricité). L'initiative \all{„Grüner Sahel“} vise à reverdir les zones arides. Vocabulaire utile : \all{der Partner, die Förderung, nachhaltig, die Wiederaufforstung, der Klimaschutz}.
\end{savaistu}

\newpage

\coursec{Exercices}

\courssub{Je vérifie mes connaissances}

\begin{exercice}[QCM : Umwelt und Klima]
Entoure la bonne réponse.
\begin{enumerate}
\item \all{die Erderwärmung} signifie :
\begin{enumerate}
\item la déforestation \item le réchauffement climatique \item la sécheresse
\end{enumerate}
\item Quel est le pluriel de \all{der Wald} ?
\begin{enumerate}
\item \all{die Walds} \item \all{die Wälder} \item \all{die Walden}
\end{enumerate}
\item \all{Die Gletscher schmelzen} veut dire :
\begin{enumerate}
\item Les glaciers fondent. \item Les glaciers grandissent. \item Les glaciers se forment.
\end{enumerate}
\item Quelle conjonction introduit une subordonnée de conséquence ?
\begin{enumerate}
\item \all{als} \item \all{so dass} \item \all{bevor}
\end{enumerate}
\end{enumerate}
\end{exercice}

\begin{exercice}[Vrai ou faux ? Justifie ta réponse]
Réponds par \all{richtig} ou \all{falsch} et justifie en une phrase complète en allemand.
\begin{enumerate}
\item \all{CO$_2$ ist ein Gas, das das Klima erwärmt.}
\item \all{Eine Dürre bedeutet zu viel Regen.}
\item \all{Der Konjunktiv II drückt die Wirklichkeit aus.}
\item \all{Nach „so dass“ steht das Verb am Ende des Satzes.}
\end{enumerate}
\end{exercice}

\begin{exercice}[Vocabulaire : articles et pluriels]
Complète le tableau avec l'article et le pluriel de chaque nom.

\begin{center}
\begin{tabular}{|l|l|l|}
\hline
\rowcolor{popblueL} Nom & Article & Pluriel \\
\hline
Umweltverschmutzung & \ldots & \ldots \\
\hline
Klima & \ldots & \ldots \\
\hline
Dürre & \ldots & \ldots \\
\hline
Gletscher & \ldots & \ldots \\
\hline
Müll & \ldots & \ldots \\
\hline
Ernte & \ldots & \ldots \\
\hline
\end{tabular}
\end{center}
\end{exercice}

\begin{exercice}[Conjugaison : mets les verbes au subjonctif II]
\begin{enumerate}
\item \all{Wenn wir weniger Plastik (kaufen), (haben) wir weniger Müll.}
\item \all{Ich (sein) froh, wenn es (regnen).}
\item \all{Wenn die Menschen mehr Bäume (pflanzen), (werden) die Luft sauberer.}
\item \all{Ohne Autos (können) wir besser atmen.}
\end{enumerate}
\end{exercice}

\courssub{Je m'entraîne}

\begin{exercice}[Transformation : du réel à l'irréel]
Transforme les hypothèses réelles en hypothèses irréelles au subjonctif II.

Exemple : \all{Wenn es regnet, bleibe ich zu Hause.} $\rightarrow$ \all{Wenn es regnen würde, bliebe ich zu Hause.}
\begin{enumerate}
\item \all{Wenn die Temperatur steigt, schmelzen die Gletscher.}
\item \all{Wenn wir den Müll trennen, schützen wir die Umwelt.}
\item \all{Wenn es nicht regnet, vertrocknen die Felder.}
\item \all{Wenn ich Zeit habe, pflanze ich einen Baum.}
\end{enumerate}
\end{exercice}

\begin{exercice}[Relie les phrases avec la conjonction indiquée]
Attention à la place du verbe conjugué !
\begin{enumerate}
\item \all{Die Temperatur steigt. Die Gletscher schmelzen.} (\all{so dass})
\item \all{Der Regen blieb aus. Die Felder vertrockneten.} (\all{als})
\item \all{Die Bauern ernten. Der Regen kommt.} (\all{bevor})
\item \all{Es war sehr heiß. Die Flüsse trockneten aus.} (\all{so dass})
\end{enumerate}
\end{exercice}

\begin{exercice}[Als, wenn, nachdem ou bevor ?]
Complète avec la conjonction qui convient.
\begin{enumerate}
\item \all{\ldots{} ich ein Kind war, gab es mehr Wald.}
\item \all{\ldots{} der Regen ausgeblieben war, begann die Dürre.}
\item \all{\ldots{} es im Frühling regnet, freuen sich die Bauern.}
\item \all{Wir müssen handeln, \ldots{} es zu spät ist.}
\end{enumerate}
\end{exercice}

\begin{exercice}[Texte à trous : Le climat du Cameroun]
Complète le texte avec les mots suivants : \all{die Dürre -- der Wald -- das Klima -- die Erderwärmung -- die Ernten -- regnen}.

\begin{textebox}[Das Klima in Kamerun]
\all{In Kamerun verändert sich \ldots{} immer schneller. Wegen \ldots{} wird es jedes Jahr wärmer. Im Norden des Landes gibt es oft \ldots{} : es will monatelang nicht \ldots{} . Die Bauern verlieren ihre \ldots{} und die Tiere finden kein Wasser. Auch \ldots{} wird kleiner, weil die Menschen viele Bäume fällen. Wenn wir nichts tun würden, hätten unsere Kinder eine schwere Zukunft.}
\end{textebox}
\end{exercice}

\courssub{Je me prépare au Bac}

\begin{exercice}[Texte de compréhension : la déforestation en Afrique centrale]
Lis le texte suivant, puis réponds aux questions \textbf{en allemand, par des phrases complètes}.

\begin{textebox}[Der Wald verschwindet]
\all{Der Regenwald in Zentralafrika ist einer der größten Wälder der Welt. Er ist wichtig für das Klima der ganzen Erde, denn die Bäume nehmen CO$_2$ aus der Luft. Aber jedes Jahr verschwinden große Flächen. Die Menschen fällen Bäume, weil sie Holz verkaufen und Felder für die Landwirtschaft brauchen. Auch Straßen und Minen zerstören den Wald.}

\all{Wenn der Wald kleiner wird, leiden die Tiere und die Menschen. Viele Arten verlieren ihr Zuhause. Die Bauern bemerken, dass der Regen unregelmäßiger wird. Als die letzte Dürre kam, vertrockneten viele Felder, so dass die Familien hungern mussten.}

\all{Experten sagen: Wenn alle Länder den Wald schützen würden, hätte die Erde eine bessere Zukunft. Junge Leute in Kamerun pflanzen schon neue Bäume, denn sie wollen ihre Umwelt retten.}
\end{textebox}

\textbf{Fragen zum Text :}
\begin{enumerate}
\item \all{Warum ist der Regenwald in Zentralafrika wichtig für das Klima?}
\item \all{Warum fällen die Menschen die Bäume? Nennen Sie zwei Gründe.}
\item \all{Was passiert mit den Tieren, wenn der Wald kleiner wird?}
\item \all{Was bemerken die Bauern?}
\item \all{Was machen die jungen Leute in Kamerun für die Umwelt?}
\end{enumerate}

\textbf{Questions de langue :}
\begin{enumerate}
\item Relève dans le texte une phrase au subjonctif II et explique ce qu'elle exprime.
\item Relève une subordonnée de conséquence et souligne la conjonction.
\end{enumerate}
\end{exercice}

\begin{exercice}[Thème et version]
\textbf{A. Version} — Traduis en français :
\begin{enumerate}
\item \all{„Wenn alle Menschen die Umwelt schützen würden, hätten wir eine bessere Zukunft.“}
\item \all{„Es war so heiß, dass die Flüsse austrockneten.“}
\end{enumerate}

\textbf{B. Thème} — Traduis en allemand :
\begin{enumerate}
\item « Si nous recyclions plus, la ville serait plus propre. »
\item « Quand la sécheresse est arrivée, les paysans ont perdu leurs récoltes. »
\end{enumerate}
\end{exercice}

\begin{exercice}[Expression écrite type Bac]
\textbf{Sujet :} \all{Was können Jugendliche gegen den Klimawandel tun?}

Rédige un texte de \textbf{10 à 12 lignes} en allemand. Tu dois :
\begin{itemize}
\item présenter le problème du changement climatique en une ou deux phrases ;
\item proposer au moins \textbf{trois actions concrètes} pour les jeunes (à l'école, à la maison, en ville) ;
\item employer \textbf{au moins une phrase au subjonctif II} (par exemple : \all{Wenn wir \ldots{} würden, \ldots{}}) ;
\item employer \textbf{au moins une subordonnée de conséquence} avec \all{so dass} ;
\item terminer par une conclusion personnelle.
\end{itemize}

\textbf{Mots utiles :} \all{der Müll trennen, Wasser sparen, Bäume pflanzen, das Fahrrad fahren, Plastik vermeiden, die Umwelt schützen}.
\end{exercice}

\newpage

\coursec{Corrigés des exercices}

\begin{corrige}[Exercice 1 — QCM : Umwelt und Klima]
\begin{enumerate}
\item \textbf{b} — \all{die Erderwärmung} signifie \emph{le réchauffement climatique} (littéralement « le réchauffement de la Terre » : \all{die Erde} + \all{die Erwärmung}).
\item \textbf{b} — Le pluriel de \all{der Wald} est \all{die Wälder} (pluriel en \all{-er} avec Umlaut, typique des noms masculins monosyllabiques).
\item \textbf{a} — \all{Die Gletscher schmelzen} veut dire « Les glaciers fondent. » (\all{schmelzen} = fondre, verbe fort : \all{schmilzt, schmolz, ist geschmolzen}).
\item \textbf{b} — \all{so dass} introduit une subordonnée de conséquence (« si bien que »). \all{als} et \all{bevor} sont des conjonctions temporelles.
\end{enumerate}
\end{corrige}

\begin{corrige}[Exercice 2 — Vrai ou faux ?]
\begin{enumerate}
\item \all{richtig} — \all{CO$_2$ ist ein Treibhausgas, das die Wärme in der Atmosphäre hält und das Klima erwärmt.} (Le CO$_2$ est un gaz à effet de serre qui retient la chaleur.)
\item \all{falsch} — \all{Eine Dürre bedeutet, dass es lange nicht regnet und der Boden trocken wird.} (Une sécheresse, c'est l'absence prolongée de pluie, pas un excès de pluie.)
\item \all{falsch} — \all{Der Konjunktiv II drückt das Irreale aus, zum Beispiel Wünsche und Hypothesen.} (Le subjonctif II exprime l'irréel, non la réalité.)
\item \all{richtig} — \all{Nach „so dass“ steht das konjugierte Verb am Ende des Nebensatzes.} (Après \all{so dass}, le verbe conjugué se place en fin de subordonnée.)
\end{enumerate}
\end{corrige}

\begin{corrige}[Exercice 3 — Vocabulaire : articles et pluriels]
\begin{center}
\begin{tabular}{|l|l|l|}
\hline
\rowcolor{popblueL} Nom & Article & Pluriel \\
\hline
Umweltverschmutzung & \all{die} & \all{die Umweltverschmutzungen} \\
\hline
Klima & \all{das} & \all{die Klimate} (ou \all{Klimata}, rare) \\
\hline
Dürre & \all{die} & \all{die Dürren} \\
\hline
Gletscher & \all{der} & \all{die Gletscher} \\
\hline
Müll & \all{der} & \all{—} (nom sans pluriel courant) \\
\hline
Ernte & \all{die} & \all{die Ernten} \\
\hline
\end{tabular}
\end{center}
\textbf{Points de vigilance :} les noms en \all{-ung} sont toujours féminins (\all{die}) et prennent \all{-en} au pluriel ; \all{der Gletscher} a un pluriel identique au singulier ; \all{der Müll} est un nom de matière, généralement sans pluriel.
\end{corrige}

\begin{corrige}[Exercice 4 — Conjugaison : subjonctif II]
\begin{enumerate}
\item \all{Wenn wir weniger Plastik \textbf{kaufen würden}, \textbf{hätten} wir weniger Müll.} (« Si nous achetions moins de plastique, nous aurions moins de déchets. »)
\item \all{Ich \textbf{wäre} froh, wenn es \textbf{regnen würde}.} (« Je serais content(e) s'il pleuvait. »)
\item \all{Wenn die Menschen mehr Bäume \textbf{pflanzen würden}, \textbf{würde} die Luft sauberer.} (« Si les gens plantaient plus d'arbres, l'air deviendrait plus propre. »)
\item \all{Ohne Autos \textbf{könnten} wir besser atmen.} (« Sans voitures, nous pourrions mieux respirer. »)
\end{enumerate}
\textbf{Rappel :} on forme le subjonctif II avec \all{würde} + infinitif pour la plupart des verbes ; mais \all{sein} $\rightarrow$ \all{wäre}, \all{haben} $\rightarrow$ \all{hätte}, \all{können} $\rightarrow$ \all{könnte} ont des formes propres qu'il faut employer.
\end{corrige}

\begin{corrige}[Exercice 5 — Transformation : du réel à l'irréel]
\begin{enumerate}
\item \all{Wenn die Temperatur \textbf{steigen würde}, \textbf{würden} die Gletscher \textbf{schmelzen}.} (« Si la température montait, les glaciers fondraient. »)
\item \all{Wenn wir den Müll \textbf{trennen würden}, \textbf{würden} wir die Umwelt \textbf{schützen}.} (« Si nous trions les déchets, nous protégerions l'environnement. »)
\item \all{Wenn es nicht \textbf{regnen würde}, \textbf{würden} die Felder \textbf{vertrocknen}.} (« S'il ne pleuvait pas, les champs se dessécheraient. »)
\item \all{Wenn ich Zeit \textbf{hätte}, \textbf{würde} ich einen Baum \textbf{pflanzen}.} (« Si j'avais le temps, je plantera un arbre. »)
\end{enumerate}
\textbf{Méthode :} à l'indicatif, la condition est réelle ou probable ; au subjonctif II, elle devient imaginaire. Attention à la place du verbe : dans la subordonnée avec \all{wenn}, le verbe conjugué va à la fin ; dans la principale qui suit, il occupe la première place (\all{würden die Gletscher schmelzen}).
\end{corrige}

\begin{corrige}[Exercice 6 — Relie les phrases]
\begin{enumerate}
\item \all{Die Temperatur steigt, so dass die Gletscher schmelzen.} (« La température monte, si bien que les glaciers fondent. »)
\item \all{Als der Regen ausblieb, vertrockneten die Felder.} (« Quand la pluie fit défaut, les champs se desséchèrent. » — \all{als} pour un événement unique au passé.)
\item \all{Die Bauern ernten, bevor der Regen kommt.} (« Les paysans récoltent avant que la pluie n'arrive. »)
\item \all{Es war sehr heiß, so dass die Flüsse austrockneten.} (« Il faisait si chaud que les rivières se sont asséchées. »)
\end{enumerate}
\textbf{Point de vigilance :} dans toutes ces subordonnées, le verbe conjugué se place \textbf{en fin de phrase} ; quand la subordonnée est en tête (phrase 2), la principale commence par son verbe conjugué (\all{vertrockneten die Felder}).
\end{corrige}

\begin{corrige}[Exercice 7 — Als, wenn, nachdem ou bevor ?]
\begin{enumerate}
\item \all{\textbf{Als} ich ein Kind war, gab es mehr Wald.} (\all{als} : événement unique, période révolue, au passé.)
\item \all{\textbf{Nachdem} der Regen ausgeblieben war, begann die Dürre.} (\all{nachdem} + plus-que-parfait : antériorité — la pluie a fait défaut \emph{avant} le début de la sécheresse.)
\item \all{\textbf{Wenn} es im Frühling regnet, freuen sich die Bauern.} (\all{wenn} : action répétée, chaque printemps.)
\item \all{Wir müssen handeln, \textbf{bevor} es zu spät ist.} (\all{bevor} : « avant qu'il ne soit trop tard ».)
\end{enumerate}
\textbf{Rappel :} \all{als} = passé unique ; \all{wenn} = présent, futur ou répétition au passé ; \all{nachdem} = antériorité ; \all{bevor} = postériorité.
\end{corrige}

\begin{corrige}[Exercice 8 — Texte à trous : Le climat du Cameroun]
\all{In Kamerun verändert sich \textbf{das Klima} immer schneller. Wegen \textbf{der Erderwärmung} wird es jedes Jahr wärmer. Im Norden des Landes gibt es oft \textbf{Dürren} : es will monatelang nicht \textbf{regnen}. Die Bauern verlieren ihre \textbf{Ernten} und die Tiere finden kein Wasser. Auch \textbf{der Wald} wird kleiner, weil die Menschen viele Bäume fällen. Wenn wir nichts tun würden, hätten unsere Kinder eine schwere Zukunft.}

\textbf{Traduction :} « Au Cameroun, le climat change de plus en plus vite. À cause du réchauffement climatique, il fait chaque année plus chaud. Dans le nord du pays, il y a souvent des sécheresses : il ne veut pas pleuvoir pendant des mois. Les paysans perdent leurs récoltes et les animaux ne trouvent pas d'eau. La forêt aussi diminue, parce que les gens abattent beaucoup d'arbres. Si nous ne faisions rien, nos enfants auraient un avenir difficile. »

\textbf{Points de vigilance :} \all{wegen} + génitif/datif (\all{wegen der Erderwärmung}) ; \all{es gibt oft Dürren} exige le pluriel ; la dernière phrase est un bel exemple de subjonctif II (\all{würden \ldots{} hätten}).
\end{corrige}

\begin{corrige}[Exercice 9 — Texte de compréhension : la déforestation]
\textbf{Fragen zum Text — réponses modèles :}
\begin{enumerate}
\item \all{Der Regenwald ist wichtig für das Klima, weil die Bäume CO$_2$ aus der Luft nehmen.}
\item \all{Die Menschen fällen die Bäume, weil sie Holz verkaufen wollen und weil sie Felder für die Landwirtschaft brauchen.}
\item \all{Wenn der Wald kleiner wird, verlieren viele Arten ihr Zuhause.}
\item \all{Die Bauern bemerken, dass der Regen unregelmäßiger wird.}
\item \all{Die jungen Leute in Kamerun pflanzen neue Bäume, weil sie ihre Umwelt retten wollen.}
\end{enumerate}

\textbf{Questions de langue :}
\begin{enumerate}
\item Phrase au subjonctif II : \all{„Wenn alle Länder den Wald schützen würden, hätte die Erde eine bessere Zukunft.“} Elle exprime une \textbf{hypothèse irréelle} : les pays ne protègent pas (encore) tous la forêt, mais on imagine la conséquence si c'était le cas.
\item Subordonnée de conséquence : \all{„\ldots{} vertrockneten viele Felder, \textbf{so dass} die Familien hungern mussten.“} La conjonction \all{so dass} introduit la conséquence (les familles ont dû souffrir de la faim) ; le verbe conjugué \all{mussten} est en fin de phrase.
\end{enumerate}
\textbf{Barème indicatif (Bac) :} compréhension 10 points (2 par question : 1 pour le contenu exact, 1 pour la correction de la langue) ; questions de langue 4 points (2 par question : relevé exact + explication).
\end{corrige}

\begin{corrige}[Exercice 10 — Thème et version]
\textbf{A. Version}
\begin{enumerate}
\item « Si tout le monde protégeait l'environnement, nous aurions un avenir meilleur. » (Subjonctif II \all{würden schützen / hätten} $\rightarrow$ conditionnel français.)
\item « Il faisait si chaud que les rivières se sont asséchées. » (\all{so \ldots{} dass} : conséquence ; \all{austrockneten} : prétérit de \all{austrocknen}.)
\end{enumerate}

\textbf{B. Thème}
\begin{enumerate}
\item \all{Wenn wir mehr recyceln würden, wäre die Stadt sauberer.} (Irréel : \all{würden + Infinitiv} ; \all{sein} au subjonctif II : \all{wäre}.)
\item \all{Als die Dürre kam, verloren die Bauern ihre Ernten.} (Événement unique au passé : \all{als} ; prétérits : \all{kam}, \all{verloren} — verbe fort \all{verlieren, verlor, verloren}.)
\end{enumerate}
\textbf{Points de vigilance :} ne jamais traduire « si » par \all{ob} ; « quand » au passé unique = \all{als}, pas \all{wenn} ; le verbe conjugué de la principale suit immédiatement la subordonnée placée en tête.
\end{corrige}

\begin{corrige}[Exercice 11 — Expression écrite type Bac]
\textbf{Proposition de rédaction modèle :}

\all{Der Klimawandel ist ein großes Problem für unsere Welt. Die Erde wird wärmer, die Gletscher schmelzen und viele Regionen leiden unter Dürren. Aber auch Jugendliche können etwas tun.}

\all{Zuerst können wir zu Hause Wasser und Strom sparen, so dass wir weniger Energie verbrauchen. Zweitens sollten wir den Müll trennen und Plastik vermeiden, zum Beispiel mit einer Stofftasche beim Einkaufen. Drittens können wir in der Stadt das Fahrrad fahren oder zu Fuß gehen statt das Auto zu nehmen. An unserer Schule in Yaoundé pflanzen wir außerdem Bäume im Schulhof.}

\all{Wenn alle Jugendlichen so handeln würden, hätte unsere Erde eine bessere Zukunft. Meiner Meinung nach ist jeder kleine Schritt wichtig, denn die Umwelt gehört uns allen.}

\textbf{Traduction :} « Le changement climatique est un grand problème pour notre monde. La Terre se réchauffe, les glaciers fondent et beaucoup de régions souffrent de sécheresses. Mais les jeunes aussi peuvent agir. D'abord, nous pouvons économiser l'eau et l'électricité à la maison, si bien que nous consommons moins d'énergie. Ensuite, nous devrions trier les déchets et éviter le plastique, par exemple avec un sac en tissu pour les courses. Enfin, nous pouvons faire du vélo ou marcher en ville au lieu de prendre la voiture. Dans notre école à Yaoundé, nous plantons en outre des arbres dans la cour. Si tous les jeunes agissaient ainsi, notre Terre aurait un avenir meilleur. À mon avis, chaque petit pas est important, car l'environnement nous appartient à tous. »

\textbf{Barème indicatif (sur 20) :}
\begin{itemize}
\item Respect du sujet et de la longueur (10--12 lignes) : 4 points.
\item Trois actions concrètes présentées : 3 points.
\item Une phrase au subjonctif II correcte : 3 points.
\item Une subordonnée avec \all{so dass} correcte (verbe en fin) : 3 points.
\item Conclusion personnelle : 2 points.
\item Correction grammaticale et richesse du vocabulaire : 5 points.
\end{itemize}

\textbf{Points de vigilance :} verbe conjugué en deuxième position dans la principale et en fin de subordonnée ; majuscule à tous les noms (\all{der Klimawandel, die Erde, das Fahrrad}) ; éviter le faux ami \all{die Jugend} (la jeunesse) pour « les jeunes » $\rightarrow$ \all{die Jugendlichen} ; ne pas oublier l'Umlaut : \all{würden, hätte}.
\end{corrige}

\newpage

\coursec{Fiche bilan}

\begin{popfigure}
\begin{tikzpicture}[mindmap, grow cyclic, every node/.style={concept}, text width=2.6cm, align=center,
  root concept/.append style={concept color=popblue, font=\bfseries, text width=3cm},
  level 1 concept/.append style={level distance=3.6cm, sibling angle=60, font=\small},
  level 2 concept/.append style={level distance=2.4cm, font=\footnotesize, text width=2.2cm}]
\node[root concept, align=center]{AL13\\Pollution et réchauffement climatique}
  child[concept color=popgreen]{ node[align=center]{Texte support\\„Die Erde wird wärmer“}
    child[concept color=popgreen!60]{ node {Compréhension globale puis détaillée} }
  }
  child[concept color=poporange]{ node[align=center]{Lexique\\Umwelt und Klima}
    child[concept color=poporange!60]{ node[align=center]{die Umwelt-\\verschmutzung, die Dürre, das CO$_2$} }
  }
  child[concept color=poppink]{ node[align=center]{Grammaire 1\\Konjunktiv II}
    child[concept color=poppink!60]{ node {würde + Infinitiv ; hätte, wäre, könnte} }
  }
  child[concept color=popblue!70]{ node[align=center]{Grammaire 2\\Conséquence et temps}
    child[concept color=popblue!50]{ node {sodass ; als, wenn, nachdem} }
  }
  child[concept color=popgold]{ node[align=center]{Méthode\\Commentaire de texte}
    child[concept color=popgold!60]{ node {Lire, repérer, résumer, produire} }
  }
  child[concept color=poppurple]{ node[align=center]{Expression\\écrite et orale}
    child[concept color=poppurple!60]{ node {Réponses, résumé, court paragraphe} }
  };
\end{tikzpicture}
\legende{Carte mentale de la leçon AL13 : la pollution et le réchauffement climatique}
\end{popfigure}

\begin{aretenir}[L'essentiel de la leçon AL13]
\textbf{1. Lexique clé à connaître par cœur}
\begin{center}
\begin{tabularx}{\linewidth}{|X|X|}\hline
\rowcolor{popblueL} \textbf{Deutsch} & \textbf{Français} \\ \hline
die Umweltverschmutzung (-en) & la pollution de l'environnement \\ \hline
die Erderwärmung (-en) / der Klimawandel (-) & le réchauffement climatique / le changement climatique \\ \hline
das Klima (-s / Klimata) & le climat \\ \hline
das Kohlendioxid (CO$_2$) & le dioxyde de carbone \\ \hline
der Wald („-er) & la forêt \\ \hline
die Dürre (-n) & la sécheresse \\ \hline
der Treibhauseffekt (-e) & l'effet de serre \\ \hline
die Überschwemmung (-en) & l'inondation \\ \hline
der Abfall („-e) / der Müll (-s) & le déchet / les ordures \\ \hline
die erneuerbaren Energien (Pl.) & les énergies renouvelables \\ \hline
die Umwelt schützen / verschmutzen / recyceln & protéger / polluer / recycler l'environnement \\ \hline
steigen (stieg, ist gestiegen) & augmenter, monter \\ \hline
die Temperatur (-en) steigt / sinkt & la température augmente / baisse \\ \hline
der Regenwald („-er) & la forêt tropicale \\ \hline
die Naturkatastrophe (-n) & la catastrophe naturelle \\ \hline
\end{tabularx}
\end{center}

\textbf{2. Grammaire à connaître par cœur}
\begin{center}
\begin{tabularx}{\linewidth}{|X|X|}\hline
\rowcolor{popgreenL} \textbf{Règle} & \textbf{Exemple} \\ \hline
Konjunktiv II présent : \cle{würde} + infinitif (infinitif rejeté en fin de phrase) & \all{Wenn ich reich wäre, würde ich Bäume pflanzen.} « Si j'étais riche, je plantera des arbres. » \\ \hline
Formes irrégulières à connaître : \all{hätte, wäre, gäbe, könnte, müsste, sollte, wüsste} & \all{Ohne Wälder gäbe es mehr CO$_2$.} « Sans forêts, il y aurait plus de CO$_2$. » \\ \hline
Konjunktiv II passé : \all{hätte / wäre} + participe II (auxiliaire comme au Perfekt) & \all{Wir hätten früher handeln müssen.} « Nous aurions dû agir plus tôt. » \\ \hline
Conséquence : \cle{sodass} (ou \all{so dass}) + verbe en fin ; \all{so ... dass} & \all{Es regnete so stark, dass die Straßen überflutet wurden.} « Il a plu si fort que les rues ont été inondées. » \\ \hline
Temps : \cle{als} (passé, action unique), \cle{wenn} (présent, futur, répétition), \cle{nachdem} + Plusquamperfekt & \all{Als ich Kind war, gab es mehr Regen.} « Quand j'étais enfant, il y avait plus de pluie. » \\ \hline
Dans toute subordonnée : le verbe conjugué va \textbf{en dernière position} & \all{..., weil sich das Klima verändert.} « ..., parce que le climat change. » \\ \hline
\end{tabularx}
\end{center}

\textbf{3. Conjugaison du Konjunktiv II présent (formes à maîtriser)}
\begin{center}
\begin{tabular}{|l|l|l|l|}\hline
\rowcolor{poporangeL} \textbf{Personne} & \textbf{haben} & \textbf{sein} & \textbf{würde + Infinitiv} \\ \hline
ich & hätte & wäre & würde schützen \\ \hline
du & hättest & wärst / wärest & würdest schützen \\ \hline
er/sie/es & hätte & wäre & würde schützen \\ \hline
wir & hätten & wären & würden schützen \\ \hline
ihr & hättet & wärt / wäret & würdet schützen \\ \hline
sie/Sie & hätten & wären & würden schützen \\ \hline
\end{tabular}
\end{center}

\textbf{4. Méthodes express}
\begin{itemize}
\item \textbf{Comprendre le texte} : 1) lire le titre et anticiper ; 2) première lecture globale (qui ? quoi ? où ?) ; 3) deuxième lecture avec surlignage des mots-clés ; 4) répondre en phrases complètes en allemand.
\item \textbf{Résumé} : reprendre uniquement les idées principales, au présent, avec ses propres mots ; pas de citation longue, pas d'opinion personnelle.
\item \textbf{Production guidée} : introduction (thème) + 2 ou 3 arguments avec \all{denn, weil, deshalb} + conclusion avec une proposition au Konjunktiv II : \all{Man sollte ... / Wir müssten ...}
\item \textbf{Thème} : vérifier la place du verbe (2\textsuperscript{e} position en principale, finale en subordonnée), le genre des noms et la majuscule à tous les noms.
\end{itemize}
\end{aretenir}

\begin{attention}[Pièges fréquents des francophones]
\begin{itemize}
\item Ne traduis pas le conditionnel français par un Präteritum : « je plantera\emph{is} » = \all{ich würde pflanzen}, jamais \all{ich pflanzte}.
\item \all{würde} se conjugue et se place en 2\textsuperscript{e} position ; l'infinitif va \textbf{à la fin} : \all{Ich würde mehr recyceln.} (et non \all{Ich würde recyceln mehr}).
\item \all{als} et \all{wenn} ne sont pas interchangeables : au passé, pour une action unique, seul \all{als} convient : \all{Als die Dürre kam, verließen viele Bauern das Land.}
\item Après \all{nachdem}, le passé exige le Plusquamperfekt : \all{Nachdem der Regen gekommen war, kehrten die Bauern zurück.}
\item Attention au genre : \all{das Klima}, \all{das CO$_2$}, mais \all{die Dürre}, \all{die Erderwärmung}, \all{der Wald}.
\item \all{die Umwelt} = l'environnement (naturel) ; ne pas confondre avec \all{die Umgebung} (les alentours, le voisinage).
\end{itemize}
\end{attention}

\begin{aretenir}[Checklist avant le Bac]
\begin{itemize}
\item[$\square$] Je connais l'article et le pluriel des mots clés : \all{die Umweltverschmutzung, die Erderwärmung, das Klima, der Wald („-er), die Dürre (-n)}.
\item[$\square$] Je sais former le Konjunktiv II présent avec \all{würde} + infinitif en fin de phrase.
\item[$\square$] Je connais par cœur les formes irrégulières : \all{hätte, wäre, gäbe, könnte, müsste, sollte}.
\item[$\square$] Je sais exprimer un souhait ou un conseil : \all{Wir sollten die Umwelt schützen.}
\item[$\square$] Je sais former le Konjunktiv II passé : \all{hätte / wäre} + participe II.
\item[$\square$] Je distingue \all{als} (passé, action unique) et \all{wenn} (présent, futur, habitude).
\item[$\square$] J'utilise \all{sodass} et \all{so ... dass} pour exprimer la conséquence, avec le verbe en fin.
\item[$\square$] Je place toujours le verbe conjugué en dernière position dans une subordonnée.
\item[$\square$] Je mets une majuscule à tous les noms et j'utilise les guillemets allemands „ ... “.
\item[$\square$] Je sais résumer un texte en 4 ou 5 phrases au présent, avec mes propres mots.
\item[$\square$] Je sais écrire un court paragraphe d'opinion sur l'environnement avec \all{meiner Meinung nach, deshalb, außerdem}.
\item[$\square$] Je relis ma production : place du verbe, genre, cas, orthographe des Umlaute (ä, ö, ü, ß).
\end{itemize}
\end{aretenir}

\findecours

\end{document}
